\documentclass[11pt,a4paper]{article}
\usepackage[T1]{fontenc}
\usepackage[utf8]{inputenc}
\usepackage{lmodern}
\usepackage[margin=27mm]{geometry}
\usepackage{amsmath,amssymb,amsthm,mathtools,mathrsfs,aliascnt}
\usepackage{enumitem,booktabs,longtable,tabularx,xcolor}
\usepackage[final]{microtype}
\usepackage[unicode]{hyperref}
\usepackage[nameinlink,capitalize,noabbrev]{cleveref}
\usepackage{xurl}
\setlength{\emergencystretch}{2em}
\hypersetup{colorlinks=true,linkcolor=blue!45!black,citecolor=green!35!black,urlcolor=blue!50!black,pdftitle={Complete Active Open-Problem Catalogue: Non-Archimedean Pluripotential Theory},pdfauthor={}}
\setcounter{tocdepth}{2}
\setlist[enumerate,1]{label=\textup{(\arabic*)},ref=\textup{(\arabic*)}}
\setlist[description]{style=nextline,leftmargin=1.5em,labelindent=0pt,itemsep=.35\baselineskip}
\newtheorem{theorem}{Theorem}[section]
\newaliascnt{proposition}{theorem}
\newtheorem{proposition}[proposition]{Proposition}
\aliascntresetthe{proposition}
\newaliascnt{lemma}{theorem}
\newtheorem{lemma}[lemma]{Lemma}
\aliascntresetthe{lemma}
\newaliascnt{corollary}{theorem}
\newtheorem{corollary}[corollary]{Corollary}
\aliascntresetthe{corollary}
\theoremstyle{definition}
\newtheorem{problem}{Problem}[section]
\newaliascnt{definition}{theorem}
\newtheorem{definition}[definition]{Definition}
\aliascntresetthe{definition}
\newaliascnt{example}{theorem}
\newtheorem{example}[example]{Example}
\aliascntresetthe{example}
\theoremstyle{remark}
\newaliascnt{remark}{theorem}
\newtheorem{remark}[remark]{Remark}
\aliascntresetthe{remark}
\crefname{problem}{problem}{problems}
\Crefname{problem}{Problem}{Problems}
\renewcommand{\index}[1]{}
\newcommand{\etalchar}[1]{$^{#1}$}
% Global mathematical notation.  Part-specific symbols should be introduced
% locally unless they occur in at least three parts of the book.

\newcommand{\NA}{\mathrm{NA}}
\newcommand{\PSH}{\mathrm{PSH}}
\newcommand{\QPSH}{\mathrm{QPSH}}
\newcommand{\FS}{\mathrm{FS}}
\newcommand{\PL}{\mathrm{PL}}
\newcommand{\An}{\mathrm{an}}
\newcommand{\bethan}{\mathrm{\beth}}
\newcommand{\triv}{\mathrm{triv}}
\newcommand{\loc}{\mathrm{loc}}
\newcommand{\nef}{\mathrm{nef}}
\newcommand{\amp}{\mathrm{amp}}
\newcommand{\lin}{\mathrm{lin}}
\newcommand{\divis}{\mathrm{div}}
\newcommand{\model}{\mathrm{mod}}
% Upright mathematical constants and differential symbols.  Use \rd for
% ordinary differentials (including measure differentials), \dc and \ddc for
% the complex/model operators, and the primed variants in the CLD calculus.
\newcommand{\rd}{\mathrm d}
\newcommand{\rmi}{\mathrm i}
\newcommand{\rme}{\mathrm e}
\newcommand{\dc}{\rd^{\mathrm c}}
\newcommand{\ddc}{\rd\dc}
\newcommand{\dprime}{\rd'}
\newcommand{\ddprime}{\rd''}
\newcommand{\dprimeddprime}{\dprime\ddprime}
\newcommand{\Rbar}{\overline{\mathbf R}}
\newcommand{\one}{\mathbf 1}
\newcommand{\uscsup}[1][]{\operatorname*{sup}_{#1}\!{}^{*}}

\DeclareMathOperator{\Spec}{Spec}
\DeclareMathOperator{\Spf}{Spf}
\DeclareMathOperator{\Proj}{Proj}
\DeclareMathOperator{\Div}{Div}
\DeclareMathOperator{\ord}{ord}
\DeclareMathOperator{\val}{val}
\DeclareMathOperator{\Supp}{Supp}
\DeclareMathOperator{\Sing}{Sing}
\DeclareMathOperator{\Reg}{Reg}
\DeclareMathOperator{\Pic}{Pic}
\DeclareMathOperator{\NS}{NS}
\DeclareMathOperator{\Nef}{Nef}
\DeclareMathOperator{\Amp}{Amp}
\DeclareMathOperator{\Int}{Int}
\DeclareMathOperator{\RelInt}{RelInt}
\DeclareMathOperator{\Trop}{Trop}
\DeclareMathOperator{\trop}{trop}
\DeclareMathOperator{\Sk}{Sk}
\DeclareMathOperator{\MA}{MA}
\DeclareMathOperator{\Capa}{Cap}
\DeclareMathOperator{\vol}{vol}
\DeclareMathOperator{\Tr}{Tr}
\DeclareMathOperator{\DHm}{DH}
\DeclareMathOperator{\lct}{lct}
\DeclareMathOperator{\Aut}{Aut}
\DeclareMathOperator{\Ent}{Ent}
\DeclareMathOperator{\DF}{DF}

\newcommand{\cE}{\mathcal E}
\newcommand{\cH}{\mathcal H}
\newcommand{\cI}{\mathcal I}
\newcommand{\cL}{\mathcal L}
\newcommand{\cM}{\mathcal M}
\newcommand{\cO}{\mathcal O}
\newcommand{\cX}{\mathcal X}
\newcommand{\cY}{\mathcal Y}

\title{Complete Active Open-Problem Catalogue\\[.5em]\large Non-Archimedean Pluripotential Theory}
\author{}
\date{Repository snapshot: 2 October 2026\\Revised catalogue: 5 October 2026}
\begin{document}
\maketitle
\begin{abstract}
This standalone catalogue contains every active open-problem item in the
current book: 42 separately addressable targets in twelve mathematical
families. Each family includes its full problem statement, individual
atomic targets, completion criteria, and the detailed state of knowledge
recorded in the manuscript, with positive cases, negative obstructions,
remaining gaps and primary-source locators. The catalogue preserves
conditional and source-scoped results and does not assert that a recent
preprint claim has been independently proved in the book.
\end{abstract}
\tableofcontents
\clearpage
\section*{Scope, dates and how to read this catalogue}
\addcontentsline{toc}{section}{Scope, dates and how to read this catalogue}
``Complete'' refers to the active problem inventory of this book. The list
is selective as a survey of the whole subject. Stable IDs such as P02.1
are editorial addresses, not theorem numbers; gaps in the numbering are
retired or historically skipped targets and are not reassigned.

The general literature search in Chapter 72 has a cutoff of 24 August
2026, with targeted source updates through 19 September 2026 and dated
verification records and author-supplied integrations through 26 September
2026. The 2 October editorial reconciliation and the 3 October LFS descent
scope correction do not constitute a new exhaustive literature search.
The correction distinguishes rational finite-maximum closure from the
real-coefficient case and local envelopes from globally formed envelopes.
This revised edition also corrects the Chebyshev question against the
published Section~6.4 (p.~1325): it assumes positive ideal volume and
asks independence from an auxiliary K\"ahler-current potential.
The 5 October category correction uses the named BMPS refinement and
positive one-cycle inputs to exclude compact positive-dimensional
support in its finite globally embedded category. The resulting literal
regularity scope and the remaining real-PA targets are reconciled below.
This targeted source check is not an exhaustive literature update.
The 2 October source and PDF remain the historical edition. No active
atom is added or retired by this source correction.
``Open'' below means open in the exact
residual scope recorded by those sources. No atom is marked newly solved
by this export.

Each family begins with concise atomic descriptions and completion
criteria. Its full mathematical statement and detailed state-of-the-art
record follow. The latter is drawn from the active Chapter 72 text and its
included statement fragments. Proof bodies are omitted; mathematical
statements, conventions, source locators and domain restrictions are
retained. Dated verification paragraphs are historical evidence, not new
verdicts. Later supplied-source retirements supersede earlier checkpoint
residuals, as the surrounding text explicitly records.

References of the form ``Book Theorem 70.2.1'' point to the accompanying
monograph. References without the word ``Book'' point inside this catalogue.
The bibliography is embedded, so this file needs neither external TeX
inputs nor an external bibliography database. Local theorem numbering is
specific to this catalogue. The twelve problem families remain genuinely
open; a settled boundary case inside a family does not settle the family.

\section*{Additional foundational proof issue: homogeneous Siu division}
\addcontentsline{toc}{section}{Additional foundational proof issue: homogeneous Siu division}
The following source issue was highlighted by the author on 4 October
2026. It is separate from the unchanged inventory of 42 research atoms.
Let $X$ be a normal integral projective variety over an algebraically closed
trivially valued field, let $\theta\in N^1(X)_{\mathbf R}$, and let $E$ be
an effective $\mathbf Q$-Cartier divisor. For homogeneous
$\psi\in\PSH_X(\theta)$ with $\psi\leq\log|s_E|$, the division assertion is
\[
 \psi=\log|s_E|+u,
 \qquad u\in\PSH_{\mathrm{hom},X}(\theta-[E]).
\]
The identity is an additive identity on the entire analytification;
subtracting two infinite values at proper support is not its definition.
Uniqueness follows from divisorial determination. The issue is existence
at the critical normalized bound $\sup_{\{\log|s_E|=-1\}}\psi=-1$.

Boucksom--Jonsson's 2022 \emph{Global pluripotential theory over a trivially
valued field}, Theorem~6.21, pp.~736--737, claims this result in the stated
normal scope. Its proof incorrectly infers that
$\psi_i-t_i\log|s_E|$ decreases from the decrease of $\psi_i$, despite
varying $t_i$. The displayed approximation therefore does not justify
its appeal to decreasing psh closure. This is a failure of the proof;
no counterexample to the theorem is asserted.

The strict bound $\psi\leq t_0\log|s_E|$ with $t_0>1$ has a separate
fixed-coefficient proof. In smooth characteristic zero, Proposition~17
of Boucksom--Jonsson's 2024 \emph{Non-Archimedean Green's functions and
Zariski decompositions} uses the independently established envelope property
to prove division. The 2024 addendum's Theorem~A independently proves
envelopes for smooth varieties when the characteristic is zero or the
dimension is at most two; its Siu corollary itself cites the original
division theorem and is not an independent repair in general normal scope.

The remaining task for the general proof route is to prove critical
homogeneous division without using the envelope conclusion it is intended
to help prove, or to replace every critical use by independently justified
arguments. The book records this extra assertion as $\mathrm{(SD)}_k$.
Its general envelope assembly and downstream deductions through that
assembly are conditional on this input. The source checks above do not
constitute an exhaustive claim about the latest literature or a new
production-model verification.

\section*{Common notation and formulation boundaries}
\addcontentsline{toc}{section}{Common notation and formulation boundaries}
\begin{description}
\item[Fields and analytic spaces.]
\(K\) denotes a complete non-Archimedean field. Trivial or nontrivial valuation,
strictness, goodness, separatedness, pure dimension and relative boundary
are retained exactly where a statement requires them. \(X^{\mathrm{an}}\)
is Berkovich analytification. ``Ordinary open'' refers to its ordinary
topology, rather than an admissible G-cover. A statement on one fixed
polycylinder is weaker than an every-open statement obtained from an actual
basis.
\item[Semivaluations and normalized local boundaries.]
We use \(v=-\log|\cdot|\), allow \(+\infty\) on the support of a
semivaluation, and retain that support. For a local ring \(R\) and a primary
ideal \(\mathfrak a\), a normalized link is a section
\(\{v:v\text{ is centered},\ v(\mathfrak a)=1\}\).
Changing \(\mathfrak a\) is a radial change with a stated cocycle; it is not
an identification of all normalizers. The symbol \(X^\beth\) denotes the
valuative/trivially valued carrier used by the raw ideal transform.
\item[Positivity classes.]
\(\PSH\) means the specified plurisubharmonic class in each theory.
Classical psh is the upper-semicontinuous curve-test notion, with the
extended \(-\infty\) convention where stated. Real PL means the real linear
extension of the local integral piecewise-linear class. Semipositivity is
its residue/model nefness condition. An LFS function is locally a
nonnegative real linear combination of finite maxima of negative logarithms
of analytic units. The rational LFS subclass requires rational
coefficients. Finite maxima of arbitrary real-coefficient LFS functions
need not be real PL.
Ordinary CLD smoothness uses smooth tropicalizations; weak smoothness
permits harmonic tropicalizations. These classes are distinct.
CLD-v2 denotes the second generated CLD psh stage in the book.
\item[Currents and products.]
\(\dprimeddprime=\dprime\ddprime\) is the CLD Hessian. A continuous
function defines the order-zero current \([h]\). CLD means
Chambert--Loir--Ducros; GJR means Gubler--Jell--Rabinoff; BFJ means
Boucksom--Favre--Jonsson. Positivity of one Hessian current does not, by
itself, construct mixed wedge products or prove classical psh recognition.
\item[Global energy and mass.]
\(\mathcal E^1\) denotes the finite-energy class of the named theory.
Reference weights, raw intersection mass and probability normalization
must be fixed before comparing energies or measures. Quotienting by
\(\mathbf R\) means uniqueness modulo additive constants. Zero top mass
cannot be divided into a probability measure. The affine-cone probability
operator is an explicitly separate convention.
\item[Complex singularity data.]
\(\alpha\in H^{1,1}_{\mathrm{BC}}(X,\mathbf R)\) is a real Bott--Chern class.
MP means Mesquita--Piccione, DXZ means Darvas--Xia--Zhang, and BJ means
Boucksom--Jonsson. \(\mathcal I\)-equivalence records multiplier-ideal data;
\(P\)-equivalence records the finer envelope/singularity data. The latter
cannot generally be recovered from the former, even at positive mass.
\item[Polyhedral normalization.]
PA denotes piecewise-affine functions. BMPS means
Botero--Mazzon--Pille-Schneider. Literal integral slopes and the real scalar
extension are separate conventions. In dimension \(d\), the polyhedral
intersection convention differs from the real Alexandrov measure by the
factor \(d!\) and the face multiplicity. Saturated tangent lattices and
analytic multiplicities are part of an analytic comparison.
\end{description}

\section*{All 42 active items}
\addcontentsline{toc}{section}{All 42 active items}
All entries below have status \textbf{Open} in their residual scope.
The tier distribution is
\[ (T1,T2,T3,T4,T5)=(0,3,7,15,17). \]
Tiers are editorial difficulty classes: T2 denotes a focused missing
bridge, T3 a substantial theorem in a developed framework, T4 coupled
comparison inputs, and T5 a foundational or programme-level target.
They do not constitute mathematical conclusions.
\begin{longtable}{@{}p{.105\textwidth}p{.07\textwidth}p{.76\textwidth}@{}}
\toprule ID & Tier & Target (links to its full dossier)\\\midrule\endfirsthead
\toprule ID & Tier & Target\\\midrule\endhead
\bottomrule\endfoot
\hyperref[atom:P02.1]{P02.1} & T3 & Decreasing semipositive real-PL regularization\\
\hyperref[atom:P02.2]{P02.2} & T3 & Decreasing locally Fubini--Study regularization\\
\hyperref[atom:P02.3]{P02.3} & T3 & Decreasing weakly smooth psh regularization\\
\hyperref[atom:P02.4]{P02.4} & T4 & Decreasing ordinary CLD-smooth psh regularization\\
\hyperref[atom:P03.1]{P03.1} & T2 & Classical positivity of regularized local envelopes\\
\hyperref[atom:P03.2]{P03.2} & T3 & CLD-v2 positivity of regularized local envelopes\\
\hyperref[atom:P03.3]{P03.3} & T3 & Continuity of local obstacle envelopes\\
\hyperref[atom:P04.1]{P04.1} & T4 & Common local products beyond algebraizable modifications\\
\hyperref[atom:P04.2]{P04.2} & T4 & Bounded products for arbitrary classical psh functions\\
\hyperref[atom:P04.3]{P04.3} & T5 & Finite-energy and unbounded local products\\
\hyperref[atom:P04.4]{P04.4} & T5 & Common local products on singular and nonnormal germs\\
\hyperref[atom:P05.1]{P05.1} & T3 & Polyhedral regularity, with coefficient conventions fixed\\
\hyperref[atom:P05.2]{P05.2} & T3 & Polyhedral envelope orthogonality\\
\hyperref[atom:P05.5]{P05.5} & T4 & Global polarized realization of polyhedral measures\\
\hyperref[atom:P06.2]{P06.2} & T2 & Recognition of continuous zero-Hessian functions\\
\hyperref[atom:P12.2]{P12.2} & T4 & An independently geometric singular valuative psh cone\\
\hyperref[atom:P12.3]{P12.3} & T4 & An independently geometric singular envelope theory\\
\hyperref[atom:P12.4]{P12.4} & T5 & Model-independent bounded singular-local products\\
\hyperref[atom:P12.5]{P12.5} & T5 & Finite energy on singular local boundaries\\
\hyperref[atom:P12.6]{P12.6} & T4 & Geometric finite-map functoriality and conductor descent\\
\hyperref[atom:P13.1]{P13.1} & T4 & Comparison for nonproportional Reeb polarizations\\
\hyperref[atom:P13.2]{P13.2} & T4 & The image of the fixed-Reeb Monge--Ampère operator\\
\hyperref[atom:P14.1]{P14.1} & T4 & Intrinsic raw-transform image beyond the known regimes\\
\hyperref[atom:P14.4]{P14.4} & T4 & A P-sensitive big non-nef positivity theory\\
\hyperref[atom:P14.5]{P14.5} & T5 & Monge--Ampère measures in the big non-nef realization\\
\hyperref[atom:P14.6]{P14.6} & T5 & Energy in the big non-nef realization\\
\hyperref[atom:P14.8]{P14.8} & T5 & Zero-mass and formal-data reconstruction\\
\hyperref[atom:P15.1]{P15.1} & T5 & Path-independent unnormalized zero-volume energy\\
\hyperref[atom:P15.2]{P15.2} & T5 & Path-independent zero-volume Monge--Ampère theory\\
\hyperref[atom:P15.3]{P15.3} & T4 & P-sensitive lower-dimensional transcendental Okounkov bodies\\
\hyperref[atom:P16.1]{P16.1} & T5 & Finite-energy equations on singular or unibranch varieties\\
\hyperref[atom:P16.2]{P16.2} & T5 & Proper nonprojective algebraizable equation theories\\
\hyperref[atom:P16.3]{P16.3} & T5 & Finite-energy equations in positive residue characteristic\\
\hyperref[atom:P16.4]{P16.4} & T4 & Finite-energy equations over arbitrary complete valued fields\\
\hyperref[atom:P16.5]{P16.5} & T5 & Finite-energy equations for nonample positivity classes\\
\hyperref[atom:P16.6]{P16.6} & T4 & Arbitrary-support finite-energy measures\\
\hyperref[atom:P16.7]{P16.7} & T5 & Finite-energy equations on nonalgebraizable analytic spaces\\
\hyperref[atom:P19.1]{P19.1} & T5 & Entropy slopes of arbitrary finite-energy weak rays\\
\hyperref[atom:P19.2]{P19.2} & T4 & Approximation preserving energy, entropy and singularities\\
\hyperref[atom:P19.3]{P19.3} & T5 & Hybrid compactness and measures for degenerating families\\
\hyperref[atom:P19.4]{P19.4} & T5 & Equivariant, transcendental and singular hybrid extensions\\
\hyperref[atom:P21.3]{P21.3} & T2 & Bedford--Taylor independence for continuous equal-curvature pairs\\
\end{longtable}
\clearpage\section{Regularization of arbitrary classical psh functions}
\label{family:1}
\noindent\textbf{Active items:} P02.1, P02.2, P02.3, P02.4. \textbf{Status: Open in the residual scope.}

\subsection{Atomic targets}
\subsubsection*{P02.1 (T3): Decreasing semipositive real-PL regularization}
\label{atom:P02.1}
\addcontentsline{toc}{subsection}{P02.1: Decreasing semipositive real-PL regularization}
Characterize the higher-dimensional good strict analytic spaces on which every classically psh function is, near every point, the pointwise decreasing limit of semipositive real PL functions. The starting function is arbitrary in the classical curve-test class; invariant continuous toric functions form only a known special case.

\subsubsection*{P02.2 (T3): Decreasing locally Fubini--Study regularization}
\label{atom:P02.2}
\addcontentsline{toc}{subsection}{P02.2: Decreasing locally Fubini--Study regularization}
For the same arbitrary classical psh starting class, characterize the spaces admitting local decreasing approximations by locally Fubini--Study functions. This asks for the Fubini--Study conclusion itself, rather than merely real-PL or smooth approximation.

\subsubsection*{P02.3 (T3): Decreasing weakly smooth psh regularization}
\label{atom:P02.3}
\addcontentsline{toc}{subsection}{P02.3: Decreasing weakly smooth psh regularization}
Characterize the higher-dimensional good strict spaces on which every classically psh function has local decreasing weakly smooth psh approximants. Weak smoothness is defined through harmonic tropicalizations and must not be identified with ordinary CLD smoothness.

\subsubsection*{P02.4 (T4): Decreasing ordinary CLD-smooth psh regularization}
\label{atom:P02.4}
\addcontentsline{toc}{subsection}{P02.4: Decreasing ordinary CLD-smooth psh regularization}
Characterize the spaces on which every classical psh function has local decreasing approximations by ordinary CLD-smooth psh functions. The ordinary smooth requirement is part of the target; a weakly smooth sequence alone does not solve it.

\subsection{What constitutes a solution}
Give precise geometric hypotheses for the requested local approximation, a construction in the specified approximant class and a proof of decreasing convergence, or counterexamples that establish the sharp failure range. Treat scalar extension and overlaps explicitly; a net must not be silently replaced by a sequence.

\subsection{Full statement and detailed state of the art}
\label{sec:ch72-classical-regularization}

\begin{problem}\label{prob:frontier-classical-regularization}
On which higher-dimensional good strict analytic spaces is every classically
psh function locally a decreasing limit of semipositive real PL, locally
Fubini--Study, weakly smooth psh, or CLD-smooth psh functions?
\end{problem}

The four targets must be separated.  Real PL functions on the relative
interior satisfy the three-way comparison of
Book Theorem~69.8.1.  On smooth boundaryless curves,
\cite[Proposition~15.11]{GJR26} gives decreasing weakly smooth
regularization locally on each relatively compact open subset.  Wanner's
published CLD-smooth regularization had the additional residue-field or curve
hypotheses stated in \cite[Corollary~5.4]{Wanner19};
Book Theorem~19.3.4 removes them for
continuous current-psh functions on smooth proper algebraic curves.
Weakly smooth harmonic functions need not be CLD-smooth.

The attempted higher-dimensional proof combines local formal models,
regularized maxima, and harmonic tropicalizations.  Curvewise approximation
does not by itself glue: the approximants must be compatible on overlapping
tropical charts and after scalar extension.  Nor may a nonmetrizable net be
replaced by a sequence without a countable extraction argument.  The general
problem is \emph{open}; the curve and PL cases are verified boundary cases.

There is a strictly invariant higher-dimensional semipositive-real-PL
endpoint as well.  On an ordinary open convex chamber of a split torus, a
finite continuous maximal-compact-torus-invariant psh function descends to a
finite convex tropical potential.  On every displayed good strictly analytic
inverse image of a nonempty ordinary open set compactly contained in the
chamber, finite tangent-plane approximations of gently shifted
radial smoothings give one decreasing semipositive real-PL sequence.  The
explicit pullbacks are CLD-psh, hence classically psh, and the real-PL
comparison then gives semipositivity without separatedness, boundarylessness,
or pure-dimensionality.  This begins with the invariant current-psh class,
not every classically psh function.  Its real-PL conclusion by itself does
not identify the sequence as locally Fubini--Study or weakly smooth; the two
separate companion arguments below add those regularities.


\begin{corollary}[Invariant toric decreasing semipositive real-PL regularization]
\label{cor:p021-invariant-toric-semipositive-rpl-regularization}
Let \(K\) be a complete nontrivially valued non-Archimedean field, let
\(T=\mathbf G_{m,K}^n\) be a split torus with maximal compact torus \(S\),
and let \(O\subset N_{\mathbf R}\simeq\mathbf R^n\) be a nonempty ordinary
open convex set.  Put \(V=\operatorname{trop}^{-1}(O)\).  Let \(B\) be a
nonempty ordinary open set with \(\overline B\Subset O\), and assume that
\(W:=\operatorname{trop}^{-1}(B)\) is good and strictly \(K\)-analytic.
If \(u:V\to\mathbf R\) is a finite
continuous \(S\)-invariant psh function in the CLD-current sense, there is a
sequence of finite
\(S\)-invariant semipositive real-PL functions
\[
                    v_j:W\longrightarrow\mathbf R
\]
that decreases pointwise and locally uniformly to \(u|_W\).
\end{corollary}




There is also a strictly invariant higher-dimensional locally
Fubini--Study endpoint.  Rather than regularizing each real-PL approximant
on a potentially changing neighbourhood, finite supporting planes are
rounded directly to the character lattice on one fixed relatively compact
tropical open set.
Their maxima are locally Fubini--Study and retain a quantitative
interstage gap, so they form one decreasing sequence on that open set.



\begin{corollary}[Invariant toric decreasing locally Fubini--Study regularization]
\label{cor:p022-invariant-toric-lfs-regularization}
Let \(K\) be a complete nontrivially valued non-Archimedean field, let
\(T=\mathbf G_{m,K}^n\) be a split torus with maximal compact torus \(S\),
and let \(O\subset N_{\mathbf R}\simeq\mathbf R^n\) be a nonempty ordinary
open convex set.  Put \(V=\operatorname{trop}^{-1}(O)\).  If
\(u:V\to\mathbf R\) is a finite continuous \(S\)-invariant psh function,
then, for every nonempty ordinary open set \(B\) with \(\overline B\Subset O\),
there is a sequence of finite \(S\)-invariant locally Fubini--Study
functions
\[
                    v_j:\operatorname{trop}^{-1}(B)\longrightarrow\mathbf R
\]
that decreases pointwise and locally uniformly to
\(u|_{\operatorname{trop}^{-1}(B)}\).  Thus this class belongs locally to
the first decreasing locally Fubini--Study stage.
\end{corollary}




There is also a strictly invariant higher-dimensional CLD-smooth endpoint.
On one open convex chamber of a split torus, the finite continuous
maximal-compact-torus-invariant psh functions pull back from finite convex
tropical potentials.  Radial mollification on the tropical chamber therefore
gives a decreasing sequence of ordinary CLD-smooth psh functions after
pullback.  This is a current-psh invariant subfamily, not a result for all
classically psh functions.


\begin{corollary}[Invariant toric decreasing CLD-smooth regularization]
\label{cor:p024-invariant-toric-cld-smooth-regularization}
Let \(K\) be a complete nontrivially valued non-Archimedean field, let
\(T=\mathbf G_{m,K}^n\) be a split torus with maximal compact torus \(S\),
and let \(O\subset N_{\mathbf R}\simeq\mathbf R^n\) be a nonempty ordinary
open convex set.  Put \(V=\operatorname{trop}^{-1}(O)\).  If
\(u:V\to\mathbf R\) is a finite continuous \(S\)-invariant psh function,
then, for every nonempty ordinary open set \(B\) with \(\overline B\Subset O\),
there is a sequence of finite \(S\)-invariant CLD-smooth psh functions
\[
                    u_j:\operatorname{trop}^{-1}(B)\longrightarrow\mathbf R
\]
that decreases pointwise and locally uniformly to
\(u|_{\operatorname{trop}^{-1}(B)}\).  Thus \(u\) belongs locally to the
first decreasing CLD-smooth psh stage.
\end{corollary}




The same invariant-toric sequence also has the weakly smooth regularity
required by the third target.  This uses only the inclusion of ordinary
CLD-smooth forms into weakly smooth forms and retains the same
CLD-current positivity; it is not a comparison of two positivity notions.


\begin{corollary}[Invariant toric decreasing weakly smooth regularization]
\label{cor:p023-invariant-toric-weakly-smooth-regularization}
Under the hypotheses of
\Cref{cor:p024-invariant-toric-cld-smooth-regularization}, the sequence
\((u_j)_j\) there is a decreasing sequence of finite \(S\)-invariant weakly
smooth psh functions on \(\operatorname{trop}^{-1}(B)\), converging pointwise
and locally uniformly to \(u|_{\operatorname{trop}^{-1}(B)}\).
Consequently, this finite continuous invariant current-psh class belongs
locally to the first decreasing weakly smooth psh stage.
\end{corollary}




The convexity imposed above is only a device for making one common
regularization neighbourhood; it is not a germwise restriction on the
toric domain.  Indeed, the four conclusions localize simultaneously on an
arbitrary tropical open set.
\begin{corollary}[Invariant toric germwise regularization on arbitrary tropical domains]
\label{cor:p02-invariant-toric-germwise-regularization}
Let \(K\) be a complete nontrivially valued non-Archimedean field, let
\(T=\mathbf G_{m,K}^n\) be a split torus with maximal compact torus \(S\),
let \(\Omega\subset N_{\mathbf R}\simeq\mathbf R^n\) be an ordinary open
set, and put \(U=\operatorname{trop}^{-1}(\Omega)\).  If
\(u:U\to\mathbf R\) is finite, continuous, \(S\)-invariant, and psh in the
CLD-current sense, then every \(x\in\Omega\) has ordinary open convex
neighbourhoods
\[
                  x\in B\Subset O\Subset\Omega
\]
for which \(u|_{\operatorname{trop}^{-1}(B)}\) is a pointwise and locally
uniform decreasing limit, respectively, of finite \(S\)-invariant locally
Fubini--Study functions, weakly smooth psh functions, and CLD-smooth psh
functions.  If in addition \(\operatorname{trop}^{-1}(B)\) is good and
strictly \(K\)-analytic, it is also such a limit of finite \(S\)-invariant
semipositive real-PL functions.
Thus all four displayed invariant-toric endpoint conclusions hold germwise
on every tropical open domain.
\end{corollary}

\clearpage\section{Upper-semicontinuous local obstacle envelopes}
\label{family:2}
\noindent\textbf{Active items:} P03.1, P03.2, P03.3. \textbf{Status: Open in the residual scope.}

\subsection{Atomic targets}
\subsubsection*{P03.1 (T2): Classical positivity of regularized local envelopes}
\label{atom:P03.1}
\addcontentsline{toc}{subsection}{P03.1: Classical positivity of regularized local envelopes}
Find sharp normal or geometrically unibranch higher-dimensional hypotheses under which the upper-semicontinuous regularization of the local-Fubini--Study/semipositive-real-PL envelope is classically psh for every upper-semicontinuous obstacle. The perfected-unibranch reduced strict curve range is already proved. Necessity must account for the branched, nonsplit-node, and codimension-two pinching obstructions.

\subsubsection*{P03.2 (T3): CLD-v2 positivity of regularized local envelopes}
\label{atom:P03.2}
\addcontentsline{toc}{subsection}{P03.2: CLD-v2 positivity of regularized local envelopes}
For the same obstacle envelope, determine the additional hypotheses ensuring that its upper-semicontinuous regularization belongs to the second CLD psh stage for every upper-semicontinuous obstacle. Classical positivity alone is not the requested conclusion.

\subsubsection*{P03.3 (T3): Continuity of local obstacle envelopes}
\label{atom:P03.3}
\addcontentsline{toc}{subsection}{P03.3: Continuity of local obstacle envelopes}
For every finite continuous obstacle, determine sharp additional hypotheses ensuring continuity of the regularized local-Fubini--Study/semipositive-real-PL envelope. A modulus-controlled subfamily or an invariant toric envelope does not establish continuity of the full competitor family.

\subsection{What constitutes a solution}
Prove the asserted conclusion for the full competitor family and every obstacle in the stated class. Account for branches, curve pullback, scalar extension and relative boundary, and distinguish positivity from continuity. A sufficient property of one selected equicontinuous subfamily is not a classification of all envelopes.

\subsection{Full statement and detailed state of the art}
\label{sec:ch72-local-envelope}

The three unrestricted affirmative targets formerly listed here are false,
even for a continuous real-valued obstacle on a reduced strict affinoid curve.
Indeed, Book Proposition~15.15.1 constructs
\[
 U=\mathcal M\bigl(K\langle x,y\rangle/(xy)\bigr)
\]
and a continuous obstacle whose raw local-FS/semipositive-real-PL envelope
equals \(0\) on one punctured branch and a negative constant on the other
branch, including the node.  Upper-semicontinuous regularization changes only
the nodal value.  The result is discontinuous and violates the affinoid
Shilov-boundary maximum principle, so it is neither classically psh nor
CLD-v2 psh.

The same reducible obstruction persists in every higher dimension.  For
\(d\geq1\), Book Proposition~15.15.2
takes the product with the closed strict \(d\)-polydisc and pulls the
obstacle back from \(U\).  On the zero section, the raw envelope is exactly
the one-dimensional raw envelope.  After regularization its value at the
nodal section point is \(0\), while its value at the radius-\(r\) Gauss point
on the other branch is at most \(-a\).  Pulling a hypothetical classical psh
regularization back to that branch disc contradicts its singleton-Shilov
maximum principle.  The product is reduced, strict, and pure of dimension
\(d+1\).  Thus reducibility causes classical-psh and CLD-v2 failure in every
positive dimension; no continuity statement for the product envelope is
inferred.  The exact statement and proof have SHA-256 prefixes
\texttt{cd22f9b56b743506} and \texttt{93ea55ab2a14a152}.  An independent
exact-byte hostile audit and a dependency-closed production run using
\texttt{gpt-5.6-sol} at \texttt{xhigh} both passed; the production verdict
was \texttt{correct}, with zero critical errors and zero gaps.

Even pointwise-bijective normalization is not a product-stable substitute
for normality.  The nonsplit quadratic node over \(\mathbf Q_3\) in
Book Proposition~15.15.3 is irreducible
and has a singleton normalization fibre at every Berkovich point.  After
product with the strict open disc, a point with quadratic residue field
reveals two analytic branches.  A finite continuous obstacle then has a
regularized local-FS/semipositive-real-PL envelope which violates the
singleton-Shilov maximum principle on one of them.  Thus the failure occurs
already on an irreducible strict analytic surface; it is not a reducibility
artifact.

The supplied pinched-bidisc example
Book Theorem~15.18.1 is seminormal and normal
away from a codimension-two point, but the upper regularization of its
Fubini--Study envelope fails the curve test there. Thus codimension-one
normality or unibranchness is insufficient; the general local-envelope
classification below remains open.

\begin{problem}\label{prob:frontier-unibranch-local-envelope}
Determine sharp hypotheses on a good strict analytic domain \(U\), beyond
the perfected-unibranch curve range of
Book Theorem~15.15.1, under which the
upper-semicontinuous regularization of
\[
 \sup\{h:h\text{ is locally Fubini--Study or semipositive real PL on }U,
              \ h\leq f\}
\]
is classically psh for every upper-semicontinuous obstacle
\(f:U\to\mathbf R\cup\{-\infty\}\).  Determine separately the additional
hypotheses under which it is CLD-v2 psh for every such obstacle.  For
continuous real-valued obstacles, determine the hypotheses under which the
regularized envelope is continuous.  Any answer must account for
analytic-curve pullback, scalar extension, relative boundary, and branch
compatibility.
\end{problem}

The empty-family case is harmless in the extended convention used here.  If
\(f\equiv-\infty\), every displayed competitor is real-valued and the raw
supremum, as well as its upper-semicontinuous regularization, is identically
\(-\infty\).  This is subharmonic in the extended convention of
\cite[Definition~3.4 and \S3.4.1]{GR25Classical}, and hence classically psh
by its Definition~4.1.  This differs from the older finite-component
convention.  A nonempty-family assumption is retained only when a particular
family-level result needs it.

The already-usc positive result remains valid: if the raw supremum is upper
semicontinuous, direct the finite subsets of the competitor family by
inclusion, take their maxima, and apply
\cite[Theorem~6.5]{GR25Classical}.  The counterexample shows exactly why one
cannot insert upper-semicontinuous regularization into that argument on a
branched curve.  On the positive side,
Book Theorem~15.15.1 proves the classical-psh
conclusion for every reduced separated strict analytic curve whose
normalization becomes bijective after completed perfect closure and
reduction.  In particular, over a perfect field it covers normal curves and
curves with bijective normalization, with relative boundary allowed.  Its
Book Corollary~15.15.1 applies this directly
to the local-FS/semipositive-real-PL obstacle envelope.  The product
obstructions above remove both reducible spaces and the weaker
pointwise-bijective-normalization condition from any unconditional
higher-dimensional positive classification.  The residual classical-psh
problem is therefore sharp necessity in the normal or geometrically
unibranch higher-dimensional range, together with extensions beyond the
stated perfected-normalization hypothesis.  The CLD-v2 classification and
the continuity atom also remain open.  Consequently the inventory count is
unchanged.

There is a separate family-level sufficient condition in every dimension.
If the exact nonempty competitor family is finite-valued, has finite
pointwise supremum, and has a common local modulus, then
Book Proposition~15.15.4 makes the raw
supremum continuous and classically psh, hence equal to its usc
regularization.  The exact stored packet passed the local independent
verifier with final verdict \texttt{correct}, zero critical errors, and zero
gaps (item \texttt{pi\_2640874e3c1d992c29bd2646}).  For the local
Fubini--Study/semipositive-real-PL envelope this applies only when its
entire minorant family satisfies those explicit conditions.  It does not
make them automatic, does not assert CLD-v2 pshness, and does not change any
P03 count.

The graph-tube theorem
Book Theorem~15.15.2 gives a different unrestricted
positive sector on every ordinary open subset of analytic affine space.
For any nonempty locally upper-bounded family of classically psh functions,
the upper-semicontinuous regularization of its pointwise supremum passes
every smooth strict analytic-curve test defined over the ground field.  In
particular, Book Corollary~15.15.3 applies to the full
local-Fubini--Study/semipositive-real-PL minorant family of every
upper-semicontinuous obstacle on such a chart.  The proof uses a relative
polydisc tube around the graph of the curve, its finite relative-Gauss
boundary, the affinoid maximum principle, and local openness of the tube
projection.

This does not yet give the scalar-stable classical-psh conclusion required
by P03.1.  If \(s\) is the raw supremum, the theorem after an analytic
extension \(K'/K\) controls \((s\circ\pi_{K'/K})^*\), whereas the pullback of
the ground-field envelope is \(s^*\circ\pi_{K'/K}\).  The exact identity
\[
                    (s\circ\pi_{K'/K})^*
                         =s^*\circ\pi_{K'/K}
\]
for every analytic extension upgrades the result to classical psh by
Book Corollary~15.15.2.  It holds when the scalar
projection is open, but
Book Proposition~15.15.5 records that scalar projections
need not be open even on the analytic affine line.  The example is only a
failure of this topological shortcut, not a psh-envelope counterexample.
Thus the theorem isolates scalar-regularization compatibility as the exact
remaining obstruction on smooth affine charts and changes no P03 atom or
total count.

The compact-fiber formula
Book Corollary~15.15.4 sharpens that obstruction.  Put
\[
       v_{K'}:=(s\circ\pi_{K'/K})^*.
\]
Since scalar projection is closed and surjective with compact fibers,
Book Corollary~15.15.5 gives the exact identity
\[
 s^*\circ\pi_{K'/K}(y)
   =\max_{z\in\pi_{K'/K}^{-1}(\pi_{K'/K}(y))}v_{K'}(z).
\]
Consequently the missing scalar compatibility is equivalent to
fiberwise constancy of \(v_{K'}\).  The graph-tube theorem already makes
\(v_{K'}\) curvewise psh over \(K'\), so the residual affine problem is to
prove this constancy, or to prove directly that its fiberwise maximum is
classically psh.

This formulation is sharp at the topological level.
Book Proposition~15.15.6 shows that a continuous
surjection commutes with usc regularization for every bounded raw function
if and only if it is open, and
Book Proposition~15.15.7 realizes a strict defect on
a type--\(4\) scalar fiber.  The defective datum is only a bounded
\(\{0,1\}\)-valued function: it is not shown to be psh or to arise from a
psh obstacle family.  Thus no counterexample or retirement is obtained,
and all three P03 atoms remain active.

For the affine Perron family above, the completed algebraic closure itself
causes no such defect.
Book Theorem~15.15.3 proves that, with
\(L=\widehat{K^{\mathrm{alg}}}\), scalar
projection \(V_L\to V\) is open and
\[
       s^*\circ\pi_{L/K}=(s\circ\pi_{L/K})^*.
\]
Moreover, classical psh-ness of this envelope is equivalent before and
after passage to \(L\).  Hence the stronger arbitrary-family affine Perron
route can be pursued over a complete algebraically closed field; only
further analytic extensions remain.  This does not identify the full
local-Fubini--Study or semipositive-real-PL obstacle competitor families
before and after scalar extension, so it retires no P03 atom and leaves the
total count unchanged.

The remaining obstruction is confined to a canonical closed set.
For any nonempty locally upper-bounded classically psh family,
Book Theorem~15.15.4 makes the regularized supremum
classically psh on
\[
   V\setminus\overline{\{s^*>s\}}.
\]
Moreover,
Book Corollary~15.15.6 shows that both
the strict defect and its closure pull back exactly to
\(L=\widehat{K^{\mathrm{alg}}}\).  Thus a residual failure on a smooth affine
chart must meet this closed defect support.  Nothing here shows that the
support is small or removable, or controls it after a further analytic
extension; all three P03 atoms and the total count remain unchanged.

On an ordinary splitting neighborhood,
Book Theorem~15.15.5 identifies the LFS envelope formed
there with the envelope on its complete inverse image, including empty
families. The proof transports competitors along individual sheets. It
does not identify restrictions of globally formed envelopes. In fact,
Book Remark~15.15.1 recalls the published irrational-bend
example: a finite maximum of real-coefficient LFS functions need not even
be real PL. Thus a global full-real LFS descent theorem cannot be inferred
from finite-maximum closure. The global section-palette and rational
finite-flat results below have separate proofs. P03.1--P03.3 and the total
count are unchanged.

Finite symmetry gives a second exact partial descent result.  If a finite
flat surjection between good strict analytic spaces carries a finite group
action over the base and every set-theoretic fiber is one orbit, then
Book Theorem~15.15.6 descends the fiberwise maximum of every
semipositive rational PL function.  Consequently
Book Theorem~15.15.7 proves exact pullback compatibility
for the raw semipositive rational-PL envelope. For the full classes it gives
the inequalities from pullback and inclusion of competitor families. A
bound of the cover LFS envelope by the pulled-back base real-PL envelope
requires the additional condition that each cover LFS competitor have a
real-PL orbit maximum. This condition is not automatic:
Book Proposition~15.15.8 gives a two-sheet fold
presentation with an LFS competitor whose maximum has an irrational bend
and is not real PL. No
finite orbit presentation is produced for scalar extension, and no LFS
descent or envelope regularity is proved.  Thus all three P03 atoms and the
total count remain unchanged.

The full real-PL branch nevertheless has an exact invariant-subfamily
description.  Under the same orbit-presentation hypothesis,
Book Theorem~15.15.8 identifies base semipositive real-PL
minorants with invariant cover minorants; group averaging sends every cover
minorant into that invariant family.  If the family is nonempty,
Book Theorem~15.15.9 writes the full P03 pullback defect
as a finite nonnegative fiberwise-constant symmetry-breaking term, with an
exact centered formula and an upper bound by orbit oscillation.  This turns
finite-orbit compatibility into a concrete cofinality criterion, but proves
neither vanishing nor regularity of the defect and supplies no scalar orbit
presentation.  The P03 atoms and total count therefore remain unchanged.

The symmetry defect vanishes in a further geometric range.
Book Theorem~15.15.10 treats finite surjections
with finitely many global analytic sections covering every set-theoretic
fiber.  Restriction along a section through a selected cover point transfers
each competitor value to a global base competitor, so both the full
semipositive-real-PL envelope and the P03 union envelope commute exactly with
pullback.  In an orbit presentation, one global section suffices by taking
its group translates.  The fold example in
Book Corollary~15.15.8 has exact envelope
descent even though its fiberwise maximum is not real PL.  Thus maximum
closure is not necessary, but no section palette is constructed for scalar
extension or for a general finite flat cover.  The P03 atoms and total count
remain unchanged.

The corresponding local statement needs only one analytic section through
the selected cover point.  By
Book Theorem~15.15.11, the local raw semipositive
real-PL envelopes then agree exactly.  The global pullback gain is therefore,
by Book Theorem~15.15.12, the base PL
globalization defect minus the cover PL globalization defect.  In a finite
orbit presentation, group translates of one local section cover the entire
central fiber.  Consequently base local-to-global PL cofinality forces both
cover cofinality and exact P03 compatibility on that fiber.  This removes the
local obstruction but neither constructs local sections for the relevant
scalar maps nor proves base cofinality in the unresolved higher-dimensional
range.  The P03 atoms and total count remain unchanged.

At every interior point, the local PL and LFS branches are nevertheless
cofinal.  By Book Theorem~15.15.13, local uniform LFS
approximation can be shifted downward by its uniform error, preserving the
obstacle.  The defect identity in Book Theorem~15.15.14 then
writes the PL globalization defect as the LFS globalization defect minus the
global PL--LFS branch gap.  Thus vanishing of the LFS globalization defect
collapses all four local and global envelopes.  Combined with a local section
and the finite orbit action, this forces exact P03 compatibility on the whole
fiber; if local sections exist near every base point, base LFS
local-to-global cofinality closes the cover globally.  The argument proves
neither that such local sections exist for scalar projection nor that local
LFS minorants globalize in the unresolved range.  The P03 atoms and total
count remain unchanged.

The section gate can nevertheless be removed for every finite flat
surjection between separated good boundaryless strict analytic spaces.
The characteristic-polynomial formula in
Book Lemma~15.15.4 shows that fiberwise maximum
preserves rational LFS functions, and
Book Proposition~15.15.9 approximates every local LFS
function from below by that rational subclass.  Consequently
Book Theorem~15.15.16 proves exact local LFS
and PL pullback, provided an LFS baseline is available near the complete
fiber.  With a global baseline,
Book Theorem~15.15.17 makes the base PL
globalization defect the sum of the cover defect and the global pullback
gain.  In the finite-orbit range this yields the four-defect identity at
every cover point, without a local section.  The argument neither proves
base local-to-global cofinality nor constructs the finite-flat presentation
needed for arbitrary scalar extensions, so P03.1--P03.3 and the total count
remain unchanged.

There is an exact conservation law behind this finite-flat result.  If the
obstacle has a global rational LFS minorant, define its rational
globalization defect and the global gap from rational LFS to full real LFS.
Book Theorem~15.15.18 proves that the rational
globalization defect is invariant under finite flat pullback.  The LFS
pullback gain is exactly the increase of the rational-coefficient gap, while
the full P03 gain is exactly the increase of that gap plus the PL--LFS branch
gap.  Thus full compatibility is equivalent to preservation of this total
two-stage gap.  Vanishing of the base rational globalization defect forces
all rational, LFS, and PL defects and both branch gaps to vanish on the
complete fiber.  The conservation law does not force that base vanishing or
construct the scalar-extension presentation, so the P03 atoms remain open.

The same ledger is additive in every finite flat tower.  By
Book Theorem~15.15.19, the rational
globalization defect is constant along a compatible chain, while the LFS and
full P03 gains are sums of nonnegative one-step gains.  Consequently
Book Corollary~15.15.16 shows that exact
compatibility for the composite is equivalent to exact compatibility at
every stage: no later finite flat cover can cancel an intermediate defect.
This constrains any tower-based scalar-extension strategy but neither
constructs such a tower nor makes the base rational defect vanish, so the P03
atoms remain open.

The base rational defect now has an exact boundary-collar formula.
For a local rational LFS minorant \(v\) with positive pointwise excess over
the global rational envelope,
Book Theorem~15.15.20 identifies that excess with the
infimum, over all sufficiently small relatively compact domains and all
global rational minorants, of the maximum boundary difference \(v-b\).
Hence Book Corollary~15.15.18 characterizes zero rational
defect by arbitrarily small absorption of every local competitor on some
surrounding boundary.  The stronger obstacle-level condition in
Book Definition~15.15.12 asks a global rational minorant
to approximate the obstacle itself on such a boundary.  It forces all
rational, LFS, and PL defects and both global branch gaps to vanish at the
central point, and
Book Corollary~15.15.20 propagates this
collapse through every finite flat tower.  This is a concrete sufficient
criterion, not a theorem that arbitrary normal or geometrically unibranch
higher-dimensional obstacles satisfy it; the P03 atoms and total count are
therefore unchanged.

The collar obstruction now has an exact connected-carrier form.
For every local rational competitor with positive globalization gap,
Book Theorem~15.15.21 shows that each superlevel set below
that gap has a compact connected component from the central point to the
boundary of every smaller relatively compact domain.  Conversely,
Book Corollary~15.15.21 proves that one
competitor and one positive level whose distinguished component escapes
across every admissible smaller domain are equivalent to a positive rational
globalization defect.  Thus a prospective counterexample
cannot be supported at an isolated or locally trapped high-defect component;
it must propagate through every surrounding collar.  A theorem trapping
these components in the normal or geometrically unibranch
higher-dimensional range would annihilate the complete defect ledger, but no
such trapping theorem is proved here.  The P03 atoms and total count remain
unchanged.

The scale-dependent carriers can now be made coherent.
Book Theorem~15.15.22 proves that the positive
part of the difference between a local rational minorant and the global
rational envelope has the same maximum on a relatively compact domain as on
its boundary.  Hence boundary domination by the global envelope propagates
through the domain.  More sharply,
Book Corollary~15.15.22 places all truncated
components through the point inside one intrinsic connected superlevel
component.  That component crosses every smaller boundary, is not relatively
compact in its local domain, and accumulates on the boundary when that domain
is relatively compact.
Book Corollary~15.15.23 makes the existence of
such a carrier that is not relatively compact equivalent to positivity of the
complete rational globalization defect.  The unresolved step is therefore
to exclude these coherent end-reaching carriers, not merely isolated or
scale-dependent defects, in the normal or geometrically unibranch
higher-dimensional range.  No such exclusion is proved here, so the P03
atoms and total count remain unchanged.

The end obstruction now also has an exact numerical boundary trace.
For a rational competitor on a relatively compact local domain,
Book Theorem~15.15.23 identifies the supremum of
its positive excess over the global rational envelope with the maximum,
over the actual boundary, of its interior upper trace.  Hence asymptotic
boundary domination is removable.
Book Corollary~15.15.24 recovers the pointwise excess as
the supremal height of a connected superlevel carrier whose closure reaches
that boundary, and
Book Corollary~15.15.25 characterizes
vanishing of the rational defect on an open set by zero boundary trace for
every relatively compact local competitor.  Under finite flat pullback,
Book Corollary~15.15.26 preserves the
trace at every boundary lift and bounds the number of pullback components of
each compact truncated carrier by the cardinality of the central fiber.
Thus a finite flat cover cannot dilute the end height or replace one compact
carrier by infinitely many unanchored branches.  These facts still do not
force the boundary trace to vanish in the unresolved normal or geometrically
unibranch higher-dimensional range, so the P03 atoms and total count remain
unchanged.

The truncation levels no longer require different endpoints.
For a fixed root point,
Book Theorem~15.15.24 intersects the nested
compact boundary-contact sets of all strict superlevels and obtains a
nonempty terminal boundary core.  The associated rooted trace is upper
semicontinuous, is bounded by the full boundary trace, and has maximum
exactly the pointwise defect.  Thus every positive witnessing competitor has
one boundary point reached at every strict level below its defect height.
Moreover,
Book Corollary~15.15.28 shows that every lift
of such a terminal point has one point in the central fiber lying in the
same pulled-back compact component at all levels.  Hence one may choose
both the terminal endpoint downstairs and a finite-fiber anchor upstairs
independently of the level.  This persistent-end certificate is still
compatible with a positive defect and therefore does not prove vanishing in
the unresolved normal or geometrically unibranch higher-dimensional range.
The P03 atoms and total count remain unchanged.

The finiteness of the pullback fiber is not needed for this persistent
component anchor.  By
Book Theorem~15.15.25, every component of the
inverse image of a compact connected set under a proper open surjection of
locally compact Hausdorff spaces maps onto the whole set.  The nested form
Book Theorem~15.15.26 then intersects those
components with one fixed compact fiber, and
Book Corollary~15.15.29 applies this to all
terminal defect carriers at once.  Thus infinite or nondiscrete fibers are
not by themselves an obstruction to persistent anchoring.  For scalar
extension, openness of the projection and compatibility of the analytic
defect remain separate issues, so this topological strengthening does not
prove P03.1--P03.3 or change the active count.

There is a complementary scalar-specific anchor over an algebraically closed
ground field.  Poineau's universal-point theorem gives every scalar projection
a canonical continuous section.  By
Book Lemma~15.15.9, projection is open at
each point of the section image and usc regularization commutes with pullback
there.  More strongly,
Book Theorem~15.15.27 shows that the
component through a section-selected endpoint contains the section of the
whole connected base carrier and therefore has the explicit common anchor
\(\sigma_{K'/K}(x)\) for every decreasing carrier family.
Book Corollary~15.15.31 applies this to all
strict terminal defect levels without properness, global openness, compactness,
or a fiber-cardinality hypothesis.  It controls only the canonical endpoint
lift and supplies no scalar-fiber constancy or analytic defect compatibility;
P03.1--P03.3 and the active count remain unchanged.

The canonical section now also gives an exact normal form for the scalar
regularization defect.  With
\(r_{K'/K}:=\sigma_{K'/K}\circ\pi_{K'/K}\),
Book Theorem~15.15.28 proves
\[
 (a\circ\pi_{K'/K})^*(r_{K'/K}(y))
   =a^*(\pi_{K'/K}(y))
   =\max_{z\in\pi_{K'/K}^{-1}(\pi_{K'/K}(y))}
                    (a\circ\pi_{K'/K})^*(z).
\]
Thus the strict defect is exactly the failure of the upstairs regularization
to be invariant under this continuous idempotent retraction, and it is
disjoint from the canonical section image.  For the affine Perron family,
Book Corollary~15.15.32 gives
\(u\circ\pi_{K'/K}=v_{K'}\circ r_{K'/K}\) and reduces classical psh-ness
exactly to subharmonicity of these retracted pullbacks on every extended
smooth strict curve.  The retraction is not proved analytic, and composition
with it is not proved to preserve subharmonicity.  Hence this isolates the
remaining analytic obligation but does not retire P03.1--P03.3 or change the
active count.

The curvewise defect split
Book Theorem~15.15.29 narrows this obligation
further.  For an extended smooth strict curve, let \(p_g\) be the intrinsic
usc regularization of the pulled raw supremum, let \(h_g\) be the pullback of
the extended-field regularization, and let \(q_g\) be the pullback of the
ground-field envelope.  Then
\[
                         p_g\leq h_g\leq q_g,
\]
both \(p_g\) and \(h_g\) are subharmonic, and \(q_g\) is subharmonic off the
closure of \(\{h_g<q_g\}\).  Thus the separate gap \(\{p_g<h_g\}\) cannot by
itself support failure: every local failure point must be accumulated by
the genuine canonical-retraction defect.  Book Corollary~15.15.33 turns this into a necessary
curve witness for any affine P03.1 counterexample.  No converse,
removability theorem, or vanishing of that defect is proved, so the P03
atoms and active count remain unchanged.

The isolated part of this curvewise obstruction is nevertheless
nonremovable.  By
Book Theorem~15.15.30, every isolated point
of the strict canonical-retraction defect is a strict local maximum of the
retracted pullback and hence an actual curve-test failure point.  If that
defect locus is closed discrete or locally finite, it equals the failure
locus, and the pullback is subharmonic exactly when the defect is empty.
Consequently a classically psh envelope can have no isolated scalar-defect
point on any extended test curve.  No theorem makes the defect locally
finite or excludes a nonempty defect locus without isolated points, so
P03.1--P03.3 and the active count remain unchanged.

There is nevertheless an invariant toric slice for a nonempty
obstacle family.  If \(V=\operatorname{trop}^{-1}(O)\) is the inverse image
of an ordinary open convex set in the dense tropical vector space of a split
torus, \(f=\varphi\circ\operatorname{trop}\), and
\(\varphi:O\to\mathbf R\cup\{-\infty\}\) is upper semicontinuous with a
finite convex minorant, then
\Cref{prop:p033-invariant-toric-real-pl-envelope} treats the raw envelope
over \emph{all} finite continuous maximal-compact-torus-invariant psh
minorants.  The minorant forces \(\varphi\) to be finite; on every compact
ball it gives a lower bound while upper semicontinuity gives an upper bound.
Convex chord slopes then make the tropical supremum automatically continuous.
Under tropicalization this is exactly the envelope of finite convex minorants
\(a\leq\varphi\).  Every supporting affine function of that convex envelope
is itself a real-PL minorant, so the full invariant psh envelope agrees with
the smaller invariant real-PL envelope.  The tropical toric correspondence of
\cite[Theorem~5.3.11]{BGJKK25} makes the common envelope continuous psh, while
\Cref{cor:p024-invariant-toric-cld-smooth-regularization} places it locally
in the first decreasing CLD-smooth stage, hence makes it classically psh by
Book Theorem~15.14.1.  Together with the following strict-domain
refinement, this settles the restricted invariant-competitor
classical-psh and CLD-v2 questions for every such upper-semicontinuous
obstacle, and the continuity question when the obstacle is continuous.  It
does not control the full noninvariant local-Fubini--Study/semipositive-real-PL
envelope, so it changes no open-problem count.


\begin{proposition}[Invariant toric upper-semicontinuous obstacle envelope]
\label{prop:p033-invariant-toric-real-pl-envelope}
Let $K$ be a complete nontrivially valued non-Archimedean field, let
$T=\mathbf G_{m,K}^n$ be a split torus with tropicalization
\(\operatorname{trop}:T^{\mathrm{an}}\to N_{\mathbf R}\simeq\mathbf R^n\), and let
\(O\subset N_{\mathbf R}\) be a nonempty ordinary open convex set.  Put
\(V:=\operatorname{trop}^{-1}(O)\), and let $S$ be the maximal compact torus
of $T^{\mathrm{an}}$.  Let
\(\varphi:O\to\mathbf R\cup\{-\infty\}\) be upper semicontinuous, set
\(f:=\varphi\circ\operatorname{trop}\), and write
\[
 \mathcal C_\varphi
 :=\{a:O\to\mathbf R\mid a\text{ is finite and convex},\ a\leq\varphi\}.
\]
Set
\[
 \mathcal C_\varphi^{\mathrm{RPL}}
 :=\{a\in\mathcal C_\varphi\mid a\text{ is locally real piecewise affine}\}.
\]
Let \(\mathcal P^S_f(V)\) be the set of finite continuous \(S\)-invariant
psh functions \(u\) on \(V\) satisfying \(u\leq f\).  Assume that
\(\mathcal C_\varphi\) is nonempty.  Then the functions
\[
 p:=\sup_{a\in\mathcal C_\varphi}a,
 \qquad
 P^S_Vf:=\sup_{u\in\mathcal P^S_f(V)}u,
 \qquad
 P^{S,\mathrm{RPL}}_V f
 :=\sup_{a\in\mathcal C_\varphi^{\mathrm{RPL}}}
        (a\circ\operatorname{trop})
\]
satisfy
\[
 P^{S,\mathrm{RPL}}_V f=P^S_Vf=p\circ\operatorname{trop}.
\]
The function $p$ is finite, continuous, and convex on $O$.  Consequently
\(P^S_Vf\) is finite, continuous, \(S\)-invariant, and psh in the
CLD-current sense.  It belongs locally to the first decreasing CLD-smooth
psh stage, and hence is CLD-v2 psh and classically psh; moreover it is locally a
uniform limit of smooth psh functions.  In particular,
\[
 \bigl(P^S_Vf\bigr)^*=P^S_Vf.
\]
\end{proposition}








Convexity of the tropical domain is needed above to compare with one global
real-PL envelope.  It is not needed for the envelope over all continuous
invariant psh minorants: on an arbitrary tropical open set, local convexity
is the exact replacement.
\begin{corollary}[Invariant toric psh envelope on an arbitrary tropical open set]
\label{cor:p033-invariant-toric-locally-convex-envelope}
Let \(K\), \(T\), \(S\), and \(\operatorname{trop}\) be as in
\Cref{prop:p033-invariant-toric-real-pl-envelope}.  Let
\(\Omega\subset N_{\mathbf R}\) be a nonempty ordinary open set, put
\(U:=\operatorname{trop}^{-1}(\Omega)\), and let
\(f=\varphi\circ\operatorname{trop}\), where
\(\varphi:\Omega\to\mathbf R\cup\{-\infty\}\) is upper semicontinuous.
Write
\[
 \mathcal C^{\mathrm{loc}}_\varphi
 :=\left\{a:\Omega\to\mathbf R\,\middle|\,
     \substack{a\text{ is finite, continuous, and locally convex,}\\
               a\leq\varphi}\right\},
\]
and set
\[
                         p:=\sup_{a\in\mathcal C^{\mathrm{loc}}_\varphi}a,
\]
where locally convex means convex on every sufficiently small ordinary open
convex neighbourhood.  Let \(\mathcal P^S_f(U)\) be the family of finite
continuous \(S\)-invariant CLD-current-psh functions on \(U\) bounded above
by \(f\), and set \(P^S_Uf:=\sup_{u\in\mathcal P^S_f(U)}u\).  If
\(\mathcal C^{\mathrm{loc}}_\varphi\) is nonempty, then
\[
                       P^S_Uf=p\circ\operatorname{trop}.
\]
The function \(p\) is finite, continuous, and locally convex; in particular
\(P^S_Uf\) is finite, continuous, \(S\)-invariant, and CLD-current-psh.
For every \(x\in\Omega\), there are ordinary open convex neighbourhoods
\(x\in B\Subset O\Subset\Omega\) on whose inverse image \(P^S_Uf\) is a
pointwise and locally uniform decreasing limit, respectively, of finite
\(S\)-invariant locally Fubini--Study functions, weakly smooth psh functions,
and CLD-smooth psh functions.  If \(\operatorname{trop}^{-1}(B)\) is good
and strictly \(K\)-analytic, it is also such a limit of finite
\(S\)-invariant semipositive real-PL functions.  In particular it belongs
locally to the first decreasing CLD-smooth psh stage, and hence is CLD-v2 psh
and classically psh.
\[
                         \bigl(P^S_Uf\bigr)^*=P^S_Uf.
\]
This is an assertion only for the displayed continuous invariant psh
competitor family.  It neither identifies an invariant real-PL or
locally-Fubini--Study envelope on a nonconvex domain nor admits any
noninvariant competitor.  In particular, the displayed decreasing
locally-Fubini--Study or real-PL approximants are not asserted to remain
below \(f\).
\end{corollary}


The following strict-domain refinement identifies the raw envelope over the
entire \(S\)-invariant part of the exact local-Fubini--Study/semipositive-real-PL
competitor union.  It does not admit any noninvariant competitor.

\begin{corollary}[Strict invariant semipositive-real-PL envelope]
\label{cor:p033-invariant-toric-semipositive-real-pl-envelope}
In the setting of
\Cref{prop:p033-invariant-toric-real-pl-envelope}, assume in addition that
$V$ is good, separated, strictly \(K\)-analytic, and boundaryless.  Let
\(\mathcal E^S_f(V)\) be the set of \(S\)-invariant functions
\(u:V\to\mathbf R\) with \(u\leq f\) which are locally Fubini--Study or
semipositive real PL.  Let \(\mathcal S^S_f(V)\) consist of the finite
continuous \(S\)-invariant real PL functions \(u:V\to\mathbf R\) that are
semipositive and satisfy \(u\leq f\).  Then
\(\mathcal E^S_f(V)=\mathcal S^S_f(V)\).
\[
 \sup_{u\in\mathcal E^S_f(V)}u
 =\sup_{u\in\mathcal S^S_f(V)}u
 =P^{S,\mathrm{RPL}}_Vf=P^S_Vf=p\circ\operatorname{trop}.
\]
Thus the invariant continuous-psh envelope is also the raw envelope over
the full \(S\)-invariant part of the local-Fubini--Study/semipositive-real-PL
competitor class in \Cref{prob:frontier-unibranch-local-envelope}.  More
precisely, it is finite and continuous and belongs locally to the first
decreasing CLD-smooth psh stage, hence is CLD-v2 psh and classically psh.
\end{corollary}




There is a further LFS-only refinement on every ordinary open convex set
compactly contained in the tropical chamber whose inverse image is good, separated,
strictly analytic, and boundaryless.  The local convex minorant envelope can
be recovered by a pointwise increasing sequence of finite maxima of
character-lattice affine minorants.  A common-denominator maximum keeps each
term locally Fubini--Study and below the obstacle, so the raw invariant
envelope over the LFS class alone agrees with the local convex envelope.
This still excludes every noninvariant
competitor, but it has the same upper-semicontinuous scope with a finite
convex minorant as the preceding invariant result.

\begin{corollary}[Invariant toric locally Fubini--Study minorant envelope]
\label{cor:p033-invariant-toric-lfs-minorant-envelope}
In the setting of \Cref{prop:p033-invariant-toric-real-pl-envelope}, fix an
nonempty ordinary open convex set \(B\) with \(\overline B\Subset O\), and assume that
\(V_B:=\operatorname{trop}^{-1}(B)\) is good, separated, strictly analytic,
and boundaryless.  Put
\[
 \mathcal C_{\varphi,B}
 :=\{a:B\longrightarrow\mathbf R\mid
       a\text{ is finite and convex},\ a\leq\varphi|_B\},
 \qquad
 p_B:=\sup_{a\in\mathcal C_{\varphi,B}}a.
\]
Let \(\mathcal L^S_{f,B}\) be the family of finite \(S\)-invariant locally
Fubini--Study functions \(h:V_B\to\mathbf R\) with \(h\leq f|_{V_B}\).
Then
\[
              \sup_{h\in\mathcal L^S_{f,B}}h
              =p_B\circ\operatorname{trop}.
\]
The function \(p_B\) is finite, continuous, and convex.  Moreover, there is
a pointwise nondecreasing sequence in \(\mathcal L^S_{f,B}\) that converges
locally uniformly to \(p_B\circ\operatorname{trop}\).  Thus the invariant
raw envelope over locally Fubini--Study minorants alone is already continuous on this
strict boundaryless relatively compact convex toric domain.  This applies to every
upper-semicontinuous obstacle in the preceding proposition's nonempty-minorant
scope.  It does not assert anything about noninvariant minorants.
\end{corollary}




The LFS-minorant equality is also genuinely local in the tropical domain.
The following form does not compare a global LFS envelope on a nonconvex
domain; it records the exact raw envelope after passage to one strict
boundaryless convex ball.
\begin{corollary}[Germwise invariant toric LFS obstacle envelope]
\label{cor:p033-invariant-toric-germwise-lfs-envelope}
Let \(K\), \(T\), \(S\), and \(\operatorname{trop}\) be as in
\Cref{prop:p033-invariant-toric-real-pl-envelope}.  Let
\(\Omega\subset N_{\mathbf R}\) be a nonempty ordinary open set, put
\(U:=\operatorname{trop}^{-1}(\Omega)\), and let
\(f=\varphi\circ\operatorname{trop}\), where
\(\varphi:\Omega\to\mathbf R\cup\{-\infty\}\) is upper semicontinuous.
Assume that \(\mathcal C^{\mathrm{loc}}_\varphi\) from
\Cref{cor:p033-invariant-toric-locally-convex-envelope} is nonempty.  Fix
\(x\in\Omega\), and let \(B\) be an ordinary open convex ball with
\(x\in B\), \(\overline B\Subset\Omega\), and
\(V_B:=\operatorname{trop}^{-1}(B)\) good, separated, strictly
\(K\)-analytic, and boundaryless.  Set
\[
 \mathcal C_{\varphi,B}
 :=\{a:B\longrightarrow\mathbf R\mid
       a\text{ is finite and convex},\ a\leq\varphi|_B\},
 \qquad
 p_B:=\sup_{a\in\mathcal C_{\varphi,B}}a,
\]
and let \(\mathcal L^S_{f,B}\) be the finite \(S\)-invariant locally
Fubini--Study functions on \(V_B\) bounded above by \(f|_{V_B}\).  Then
\[
                  \sup_{h\in\mathcal L^S_{f,B}}h
                  =p_B\circ\operatorname{trop}.
\]
The function \(p_B\) is finite, continuous, and convex, and the displayed
envelope is the locally uniform limit of a pointwise nondecreasing sequence
in \(\mathcal L^S_{f,B}\).  This identifies only the raw envelope on
\(V_B\); it neither identifies the corresponding raw envelope on \(U\) nor
admits a noninvariant competitor.
\end{corollary}


The same germwise ball also recovers the full invariant part of the original
local-Fubini--Study/semipositive-real-PL competitor union, not merely its LFS
subfamily.
\begin{corollary}[Germwise full invariant toric competitor envelope]
\label{cor:p033-invariant-toric-germwise-full-envelope}
In the setting of
\Cref{cor:p033-invariant-toric-germwise-lfs-envelope}, let
\(\mathcal E^S_{f,B}\) be the family of \(S\)-invariant functions
\(u:V_B\to\mathbf R\), bounded above by \(f|_{V_B}\), which are locally
Fubini--Study or semipositive real PL.  Let \(\mathcal S^S_{f,B}\) consist
of the finite continuous \(S\)-invariant semipositive real-PL functions on
\(V_B\) bounded above by \(f|_{V_B}\).  Then
\[
 \sup_{u\in\mathcal E^S_{f,B}}u
 =\sup_{u\in\mathcal S^S_{f,B}}u
 =\sup_{h\in\mathcal L^S_{f,B}}h
 =p_B\circ\operatorname{trop}.
\]
Thus this raw invariant local-Fubini--Study/semipositive-real-PL envelope is
finite, continuous, CLD-v2 psh, and classically psh on \(V_B\).  It is the
locally uniform limit from below of a pointwise nondecreasing sequence in
\(\mathcal L^S_{f,B}\).  No assertion is made for the corresponding global
envelope on \(U\) or for a noninvariant competitor.
\end{corollary}

\clearpage\section{A common local Monge--Ampère product}
\label{family:3}
\noindent\textbf{Active items:} P04.1, P04.2, P04.3, P04.4. \textbf{Status: Open in the residual scope.}

\subsection{Atomic targets}
\subsubsection*{P04.1 (T4): Common local products beyond algebraizable modifications}
\label{atom:P04.1}
\addcontentsline{toc}{subsection}{P04.1: Common local products beyond algebraizable modifications}
Extend the existing CLD/formal-BFJ common-model comparison to arbitrary nonalgebraizable analytic modifications. Specify the carrier, boundary datum, raw mass, and the precise class of compatible potentials.

\subsubsection*{P04.2 (T4): Bounded products for arbitrary classical psh functions}
\label{atom:P04.2}
\addcontentsline{toc}{subsection}{P04.2: Bounded products for arbitrary classical psh functions}
Construct a common local Monge--Ampère product for arbitrary bounded classically psh functions outside the jointly realizable model closure. Prove compatibility with each existing endpoint rather than declaring a product by classical positivity alone.

\subsubsection*{P04.3 (T5): Finite-energy and unbounded local products}
\label{atom:P04.3}
\addcontentsline{toc}{subsection}{P04.3: Finite-energy and unbounded local products}
Construct a finite-energy or unbounded extension of the common local product, with an explicit polar-mass rule. Distinguish a nonpluripolar extension from an extension retaining polar mass and state the hypotheses under which uniqueness is intended.

\subsubsection*{P04.4 (T5): Common local products on singular and nonnormal germs}
\label{atom:P04.4}
\addcontentsline{toc}{subsection}{P04.4: Common local products on singular and nonnormal germs}
Extend the local construction to singular and nonnormal germs, compatibly with normalization and conductor descent. The analytic local product and the normalized-link product have different dimensions and must be compared with their exact boundary normalizers.

\subsection{What constitutes a solution}
Every extension must specify its boundary datum and prove model independence, monotone continuity, locality and functoriality in its exact domain. Match the raw intersection normalization on the known common-model locus. State the polar rule before claiming an unbounded extension; literal restriction cannot preserve an unchanged probability normalization.

\subsection{Full statement and detailed state of the art}
\label{sec:ch72-local-ma}

\begin{problem}\label{prob:frontier-local-ma}
Construct a local product that simultaneously extends the continuous CLD
Bedford--Taylor product, the bounded formal BFJ measure on a normalized
valuation link, and the largest justified class of classically psh functions.
It should have a stated boundary datum, model independence, monotone
continuity, locality, and functoriality.
\end{problem}

The two existing endpoints have different signatures.  CLD works on analytic
opens with tropical charts and compactly supported test forms; the local BFJ
operator of Book Theorem~70.6.1 takes \(n-1\) potentials on the link
\(\nu(\mathfrak m)=1\) and uses exceptional intersection theory.  Classical
psh positivity, by itself, constructs no wedge product.  The following
result identifies the two products on one explicit model locus and then on a
jointly realizable bounded monotone closure.  It does not identify either
product for an arbitrary curve-test psh function.


Let \((X,0)\) be a smooth complex analytic germ of dimension \(n\geq2\), put
\(R=\mathcal O_{X,0}\), let \(\mathfrak m\) be its maximal ideal, and set
\(d=n-1\).  Its normalized valuation link is
\begin{equation}
 V_{\mathfrak m}
 :=\{v:R\longrightarrow[0,+\infty]\mid
       v\text{ is centered at }0,\ v(\mathfrak m)=1\}.
 \label{eq:frontier-local-ma-link}
\end{equation}
For a prime divisor \(E\) above \(0\), write
\begin{equation}
 b_E:=\ord_E(\mathfrak m),
 \qquad v_E:=b_E^{-1}\ord_E.
 \label{eq:frontier-local-ma-divisorial-normalization}
\end{equation}
Fix \(K=\mathbf C((t))\), normalized by \(|t|=\rme^{-1}\), and regular
parameters \(z_1,\ldots,z_n\) of \(R\).  We choose a closed good, possibly
non-strict, \(K\)-analytic polydisc \(B\) of polyradius \(\rho\), where
\(|t|<\rho<1\).  Taylor series of the chosen complex analytic germs converge
on \(B\), and the shell below has a collar in \(B\).  For an
\(\mathfrak m\)-primary ideal
\(\mathfrak a=(f_1,\ldots,f_N)\), put
\begin{equation}
 u_{\mathfrak a}:=\log\max_j|f_j|,
 \qquad
 u_{\mathfrak m}:=\log\max_i|z_i|,
 \qquad
 \chi_{\mathfrak m}:=\max\{u_{\mathfrak m},-1\}.
 \label{eq:frontier-local-ma-model-functions}
\end{equation}
Changing generators leaves the associated model metric unchanged after a
common admissible refinement.  Set
\begin{equation}
S_{\mathfrak m}:=\{u_{\mathfrak m}=-1\}.
\label{eq:frontier-local-ma-shell}
\end{equation}
In this subsection an \emph{algebraizable log model} means the
analytification, after shrinking the germ, of a projective modification of a
Zariski neighborhood \(U\subset\mathbf A^n_{\mathbf C}\), with the relevant
divisors and line bundles descending to that algebraic model.

For model psh potentials \(u_1,\ldots,u_d\), let
\(T=\dprimeddprime u_1\wedge\cdots\wedge\dprimeddprime u_d\).  We define
their swept boundary product by
\begin{equation}
 \begin{split}
 \MA^{\partial}_{\mathrm{CLD},\mathfrak m}
       (u_1,\ldots,u_d)
 :={}&T\wedge\dprimeddprime\chi_{\mathfrak m}\\
 &-\one_{\{u_{\mathfrak m}>-1\}}
       \bigl(T\wedge\dprimeddprime u_{\mathfrak m}\bigr),
 \end{split}
 \label{eq:frontier-local-ma-swept-product}
\end{equation}
where the restricted measure in the second line is extended by zero.  The
two products are the continuous Bedford--Taylor products in the revised CLD
psh sheaf.  The subtraction and the level \(-1\) are part of the boundary
datum; in particular, \eqref{eq:frontier-local-ma-swept-product} is neither a
probability normalization nor the raw \(d\)-fold product of restrictions to
the shell.

The definition is meaningful for the model potentials above.  Indeed, if
\(\mathfrak a\) is \(\mathfrak m\)-primary, then
\(\mathfrak m^p\subset\mathfrak a\subset\mathfrak m^q\) for suitable
positive integers \(p,q\), so \(u_{\mathfrak a}\) is bounded on a collar of
\(S_{\mathfrak m}\).  On \(\{u_{\mathfrak m}>-1\}\), locality gives
\(\dprimeddprime\chi_{\mathfrak m}
=\dprimeddprime u_{\mathfrak m}\), and the two terms cancel.  On
\(\{u_{\mathfrak m}<-1\}\), the cutoff is constant and both terms vanish.
Thus \eqref{eq:frontier-local-ma-swept-product} is supported on the shell.
This argument does not assert that
\(\dprimeddprime\chi_{\mathfrak m}\) itself is shell-supported; that
assertion is false in dimension greater than one.  The reference factor is
essential: \(u_{\mathfrak m}|_{V_{\mathfrak m}}=-1\), whose intrinsic
\(\dprimeddprime\) is zero, whereas the local BFJ measure of this datum has
mass one.

\begin{lemma}\label{lem:frontier-local-ma-gauss}
There are maps
\[
 \iota:V_{\mathfrak m}\hookrightarrow S_{\mathfrak m},
 \qquad
 p:S_{\mathfrak m}\longrightarrow V_{\mathfrak m},
 \qquad p\circ\iota=\mathrm{id},
\]
such that, for every \(\mathfrak m\)-primary ideal \(\mathfrak a\),
\begin{equation}
             u_{\mathfrak a}(\iota(v))=-v(\mathfrak a).
 \label{eq:frontier-local-ma-gauss-ideal}
\end{equation}
If \(\pi:Y\to X\) is an algebraizable log resolution of finitely many
\(\mathfrak m\)-primary ideals and
\(\mathfrak m\mathcal O_Y=\mathcal O_Y(-D)\), then every face
\begin{equation}
 \Delta(Y,D)=
 \left\{\sum_{j=1}^r w_jE_j:\ w_j\geq0,
                       \ \sum_{j=1}^r b_jw_j=1\right\}
 \label{eq:frontier-local-ma-dual-face}
\end{equation}
of its normalized dual complex is carried by \(\iota\) onto a polyhedral
skeleton in a formal model of the shell.
\end{lemma}



\begin{lemma}\label{lem:frontier-local-ma-common-model-core}
Let \(\mathfrak a\) be \(\mathfrak m\)-primary and suppose that
\(\mathfrak a\mathcal O_Y=\mathcal O_Y(-A)\) on a principalization.  Then
\(u_{\mathfrak a}\) is a continuous semipositive model function, hence a
member of the revised CLD psh sheaf on a collar of the shell; its pullback to
every analytic curve is subharmonic; and its restriction to the Gauss
section is the BFJ model formal psh function
\(g_{\mathfrak a}(v)=-v(\mathfrak a)\).  The same conclusion holds for a
rational semipositive model function associated with a relatively nef
exceptional \(\mathbf Q\)-Cartier b-divisor carried by an algebraizable
common log model.
\end{lemma}



\begin{theorem}\label{thm:frontier-local-ma-model-comparison}
Let \(\mathfrak a_1,\ldots,\mathfrak a_d\) be
\(\mathfrak m\)-primary ideals.  On one log resolution write
\[
 \mathfrak m\mathcal O_Y=\mathcal O_Y\!\left(-\sum_Eb_EE\right),
 \qquad
 \mathfrak a_i\mathcal O_Y=\mathcal O_Y(-A_i).
\]
Then
\begin{equation}
 \MA^{\partial}_{\mathrm{CLD},\mathfrak m}
       (u_{\mathfrak a_1},\ldots,u_{\mathfrak a_d})
 =\sum_E b_E\bigl((-A_1)\cdots(-A_d)\cdot E\bigr)
       \delta_{\iota(v_E)}.
 \label{eq:frontier-local-ma-cld-atomic}
\end{equation}
After applying \(p\), this is exactly
\begin{equation}
 \MA_{\mathrm{BFJ},\mathfrak m}
       (g_{\mathfrak a_1},\ldots,g_{\mathfrak a_d})
 =\sum_E b_E\bigl((-A_1)\cdots(-A_d)\cdot E\bigr)
       \delta_{v_E}.
 \label{eq:frontier-local-ma-bfj-atomic}
\end{equation}
The same comparison holds for rational semipositive model functions on an
algebraizable common log model.
\end{theorem}



\begin{proposition}\label{prop:frontier-local-ma-independence}
The common model measure of
\Cref{thm:frontier-local-ma-model-comparison} is independent of the log
resolution, formal model, and tropical charts.  With the boundary datum
\((\mathfrak m,-1)\) fixed, it is local on \(V_{\mathfrak m}\).  Literal
restriction to a relatively open subset is restriction of the measure, with
no renormalization of its mass.
\end{proposition}



Define the joint bounded closure \(\mathcal C_b^{\mathrm{com}}\) to consist
of bounded continuous BFJ formal psh functions \(g\) for which there are
decreasing semipositive model functions \(g^{(k)}\downarrow g\), each in
the preceding ideal/algebraizable model locus, whose homogeneous Gauss
extensions \(u^{(k)}\), on one fixed collar of the shell, satisfy all three
conditions:
\begin{enumerate}
\item they belong to the revised CLD psh sheaf;
\item they decrease locally boundedly to a continuous revised-CLD psh
      function \(u\);
\item their restrictions to all analytic curves are subharmonic.
\end{enumerate}
For a finite collection we require a common collar and synchronized common
algebraizable model refinements.  For a further decreasing sequence,
``joint'' means that one decreasing two-index array works simultaneously for
the outer and model indices.

\begin{proposition}\label{prop:frontier-local-ma-bounded-closure}
For \(g_1,\ldots,g_d\in\mathcal C_b^{\mathrm{com}}\), the limit
\begin{equation}
 \MA^{\mathrm{com}}_{\mathfrak m}(g_1,\ldots,g_d)
 :=\lim_{k\to\infty}
 p_*\MA^{\partial}_{\mathrm{CLD},\mathfrak m}
       (u_1^{(k)},\ldots,u_d^{(k)})
 \label{eq:frontier-local-ma-bounded-limit}
\end{equation}
exists, is independent of all approximants, and equals the BFJ mixed local
intersection measure.  It is continuous along decreasing sequences admitting
one joint bounded model approximation of the preceding kind.
\end{proposition}



\begin{proposition}\label{prop:frontier-local-ma-reference-ideal}
Let \(I\) be an \(\mathfrak m\)-primary ideal.  For ideal functions, and for
rational semipositive functions carried by an algebraizable common log
model, the preceding construction and comparison remain valid with the link
\(V_I=\{v:v(I)=1\}\), the shell \(u_I=-1\), and the cutoff
\(\max\{u_I,-1\}\).  The joint bounded extension remains valid with the
same approximation hypothesis.  We denote the resulting measure by
\(\MA_I^{\mathrm{com}}\).
\end{proposition}



\begin{proposition}\label{prop:frontier-local-ma-finite-map}
Let \(F:(X',0')\to(X,0)\) be a finite map of smooth
\(n\)-dimensional germs of local degree \(q\), and put
\(I=\mathfrak m\mathcal O_{X',0'}\).  Normalize the source link first by
\(I\), and write
\(F_\#:V_I\to V_{\mathfrak m}\), \(w\mapsto w\circ F^\#\).  For ideal
model functions, for rational functions whose target models and pulled-back
source models both lie in the algebraizable model locus, and for members of the joint bounded
closure whose pullbacks remain in that closure,
\begin{equation}
 (F_\#)_*\MA^{\mathrm{com}}_I(F^*g_1,\ldots,F^*g_d)
 =q\,\MA^{\mathrm{com}}_{\mathfrak m}(g_1,\ldots,g_d).
 \label{eq:frontier-local-ma-finite-projection}
\end{equation}
On the standard source link \(V_{\mathfrak m'}\), set
\[
 c(v):=v(I),\qquad R_I(v):=v/c(v),\qquad
 \widehat F(v):=(v\circ F^\#)/v(I)=F_\#(R_I(v)).
\]
For ideal model functions, the exact standard-link formula is
\begin{equation}
 \widehat F_*\!\left(
 c\,\MA^{\mathrm{com}}_{\mathfrak m'}
   (g_{F^*\mathfrak a_1},\ldots,g_{F^*\mathfrak a_d})\right)
 =q\,\MA^{\mathrm{com}}_{\mathfrak m}
   (g_{\mathfrak a_1},\ldots,g_{\mathfrak a_d}).
 \label{eq:frontier-local-ma-standard-link}
\end{equation}
For an \'etale map of pointed complex germs, \(q=1\),
\(I=\mathfrak m'\), and \(c\equiv1\).
\end{proposition}



\begin{proposition}\label{prop:frontier-local-ma-no-probability}
No rule can both replace every nonzero literal local restriction of the
common measure by a probability measure and form a sheaf under literal
restriction.
\end{proposition}



The next example shows why the bounded hypothesis cannot simply be deleted.

\begin{example}\label{ex:frontier-local-ma-polar-endpoint}
On \((\mathbf C^2,0)\), let
\(v_s(x)=1\), \(v_s(y)=s\) for \(1\leq s<+\infty\), and let
\(v_\infty\) be the endpoint with \(v_\infty(y)=+\infty\).  For
\(\mathfrak a_k=(y,x^k)\),
\begin{equation}
 g_k(v_s)=-v_s(\mathfrak a_k)=-\min\{s,k\}\downarrow-s=:g(v_s),
 \label{eq:frontier-local-ma-unbounded-ray}
\end{equation}
and \(g(v_\infty)=-\infty\).  Resolving \(\mathfrak a_k\) gives a chain with
\(v_j(x)=1\), \(v_j(y)=j\), exceptional intersections
\(E_j^2=-2\) for \(j<k\), \(E_k^2=-1\), and adjacent intersection one.
The ideal divisor has coefficient \(j\) along \(E_j\).  Hence
\((-A)\cdot E_j=0\) for \(j<k\) and
\((-A)\cdot E_k=1\), so
\begin{equation}
                  \MA_{\mathrm{BFJ}}(g_k)=\delta_{v_k}.
 \label{eq:frontier-local-ma-unbounded-atoms}
\end{equation}
The full decreasing limit is \(\delta_{v_\infty}\).  On the other hand,
\(g_k=\max\{g,-k\}\), and the nonpluripolar truncation
\[
 \one_{\{g>-k\}}\MA_{\mathrm{BFJ}}(\max\{g,-k\})
\]
is zero: the only atom in
\eqref{eq:frontier-local-ma-unbounded-atoms} lies on \(\{g=-k\}\).  Thus
full monotone continuation and nonpluripolar continuation disagree by a
polar endpoint mass.
\end{example}

\begin{theorem}\label{thm:frontier-common-local-ma}
For a smooth complex germ and the fixed datum
\((\mathfrak m,|t|=\rme^{-1},u_{\mathfrak m}=-1)\), the Gauss realization
of \Cref{lem:frontier-local-ma-gauss} identifies the swept CLD boundary
measure with the BFJ local measure on ideal functions and on rational
semipositive functions carried by an algebraizable common log model.  This
identification is model- and chart-independent within that locus, local
without probability renormalization, and functorial under finite maps in the
input domains specified in \Cref{prop:frontier-local-ma-finite-map}, with the
reference-change factor in \textup{(Book 72.25)}.
It extends uniquely to
\(\mathcal C_b^{\mathrm{com}}\) with joint decreasing monotone continuity.
Without an additional convention governing polar mass, the model comparison
does not determine an extension to arbitrary unbounded functions.
\end{theorem}



Consequently the model-locus comparison and the explicitly defined joint
bounded closure are \emph{verified}.  The following parts of the problem
remain \emph{open}: a product for an arbitrary
bounded classically psh function not known to admit the joint realization; a
capacity or finite-energy criterion selecting the treatment of polar mass;
comparison for a rational model carried only by a nonalgebraizable analytic
modification; and extensions to singular or nonnormal germs.  No primary
source located by the cutoff states the comparison theorem above; its proof
is a new Rethlas-verified manuscript argument, not a claim about priority in
the literature.  The Bingener comparison recorded in
Book Theorem~59.2.1 and based on
\cite[Aussage~2.8]{Bingener76Stein} algebraizes suitable projective germs over
a local or Stein algebra; it does not by itself produce a finite-type
\(\mathbf C\)-model and an admissible CLD formal model for an arbitrary
analytic modification, so it does not close the stated algebraizability gap.

\clearpage\section{Balanced polyhedral regularity and realization}
\label{family:4}
\noindent\textbf{Active items:} P05.1, P05.2, P05.5. \textbf{Status: Open in the residual scope.}

\subsection{Atomic targets}
\subsubsection*{P05.1 (T3): Polyhedral regularity, with coefficient conventions fixed}
\label{atom:P05.1}
\addcontentsline{toc}{subsection}{P05.1: Polyhedral regularity, with coefficient conventions fixed}
Characterize balanced rational polyhedral pairs \((\mathbf X,\gamma)\) for which the BMPS envelope regularity property holds. Treat the literal integral-slope/rational-constant class and its explicit real scalar extension separately. In the finite globally embedded BMPS category, every nonempty positive-dimensional balanced space is noncompact, so the existing noncompact obstruction excludes literal regularity on every such pair. The higher-dimensional real-PA classification, source coefficient-convention repair and arbitrary nonseparable obstacles remain open.

\subsubsection*{P05.2 (T3): Polyhedral envelope orthogonality}
\label{atom:P05.2}
\addcontentsline{toc}{subsection}{P05.2: Polyhedral envelope orthogonality}
Characterize the pairs for which the regularized polyhedral envelope satisfies the exact contact identity \(\int(P_\gamma(u)-(\gamma+u))\,\MA_{\mathrm{poly}}(P_\gamma(u))=0\) for all admissible obstacles. Positive growth at infinity, balancing, or correct total mass alone does not establish orthogonality.

\subsubsection*{P05.5 (T4): Global polarized realization of polyhedral measures}
\label{atom:P05.5}
\addcontentsline{toc}{subsection}{P05.5: Global polarized realization of polyhedral measures}
Characterize polarized Monge--Ampère realizations of \((\mathbf X,\gamma)\). Construct global semipositive lifts of the relative potentials \(a_i-\gamma\) in a named BFJ domain, prove intersection-mass and component-calibration identities, and compare global non-Archimedean and polyhedral measures, including mixed products and boundary terms. Local recursive CLD/delta comparison is an established subcase, not the full target.

\subsection{What constitutes a solution}
Fix the PA coefficient convention before stating regularity. Orthogonality requires the contact integral, not just total mass. A global BFJ comparison must name the polarization, operator domain, lifts, saturated lattice indices, intersection normalization and boundary terms; the already resolved local recursive comparisons are not new atoms.

\subsection{Full statement and detailed state of the art}
\label{sec:ch72-polyhedral}

\begin{problem}\label{prob:frontier-polyhedral-regularity}
Let \(\mathbf X\) be a pure-dimensional balanced rational polyhedral space
and let \(\gamma\) be a piecewise-affine polyhedrally convex reference
function positive at infinity.  Characterize the pairs
\((\mathbf X,\gamma)\) for which the
regularity and orthogonality properties of Botero--Mazzon--Pille-Schneider
hold.  Under realizability hypotheses, compare their polyhedral
Monge--Amp\`ere operator with CLD, delta-form, and non-Archimedean
Monge--Amp\`ere measures.
\end{problem}

The bibliographic dates require a correction.  The first arXiv version was
submitted on 9 March 2026 and the date printed inside the manuscript is
10 March 2026, not 11 August 2026 \cite{BMPS26}.  We use that version
throughout.

Here are the conventions needed to compare it with
Book Appendix~G.  Normal vectors point from a codimension-one
face into an adjacent top-dimensional face, as in
\textup{(Book G.20)}, and a convex bend has positive divisor.
Botero--Mazzon--Pille-Schneider allow positive real balancing weights; a
tropical cycle in the narrower Mikhalkin--Rau convention has integral
weights.  Analytic comparison therefore uses the latter only after the
saturated tangent lattices and analytic tropical multiplicities have been
installed.  Finally, their real Monge--Amp\`ere measure is the Alexandrov
measure of Book Theorem~D.5.1, without a factorial.  Their
polyhedral intersection convention consequently gives
\begin{equation}
 \one_{\RelInt(\sigma)}\MA_{\mathrm{poly}}(\varphi)
 =d![\mathbf X](\sigma)\,
   \MA_{\mathbf R}(\varphi|_{\RelInt(\sigma)})
 \label{eq:frontier-polyhedral-real-local}
\end{equation}
on the relative interior of every top-dimensional face
\cite[Propositions~4.12 and~4.20]{BMPS26}.  This is exactly the factorial
firewall in Book Appendix~J.

We first separate the source results from the new deduction below.  For a
piecewise-affine polyhedrally convex reference \(\gamma\) positive at
infinity, \cite[Theorem~4.13 and Proposition~4.16]{BMPS26} constructs the
positive mixed operator
\[
 \MA_{\mathrm{poly}}:
 \mathrm{PC}_{\mathrm{reg}}(\mathbf X,\gamma)^d
 \longrightarrow\mathcal M^+(\mathbf X),
 \qquad
 \int_{\mathbf X}\MA_{\mathrm{poly}}=\deg(\gamma),
\]
and proves its continuity under uniform, in particular decreasing,
regularizable approximation.  Its Theorem~6.4 says that regularity is
equivalent to each of
\[
 \begin{split}
 &P_\gamma(u)\in\mathrm{PC}_{\mathrm{reg}}(\mathbf X,\gamma)
       \quad\text{for every bounded piecewise-affine }u,\\
 &\mathrm{PAPC}(\mathbf X,\gamma)
       \text{ is uniformly dense in }\mathrm{PC}(\mathbf X,\gamma),\\
 &\mathrm{PC}(\mathbf X,\gamma)
       =\mathrm{PC}_{\mathrm{reg}}(\mathbf X,\gamma).
 \end{split}
\]
The source states one-dimensional regularity in Proposition~6.13 and
complete-affine and conical regularization in Remark~6.6.  These statements
must be separated from its literal coefficient convention: Definition~2.28
requires integral linear parts and rational constants, while the proofs of
the cited regularization statements use arbitrary real affine functions.
The theorem and counterexample below give the precise one-dimensional
repair.  Orthogonality is the further
identity
\begin{equation}
 \int_{\mathbf X}
 \bigl(P_\gamma(u)-(\gamma+u)\bigr)
 \MA_{\mathrm{poly}}(P_\gamma(u))=0.
 \label{eq:frontier-polyhedral-orthogonality}
\end{equation}
The source proves \eqref{eq:frontier-polyhedral-orthogonality} on
one-dimensional polyhedrally smooth spaces and obtains existence and
uniqueness, modulo constants, for compactly supported positive measures of
mass \(\deg(\gamma)\)
\cite[Proposition~6.21 and Theorem~6.22]{BMPS26}.



Write $\mathrm{PA}_{\mathbf Z}(\mathbf X)$ for the literal class in
\cite[Definition~2.28]{BMPS26}: the linear part on every face is integral and
the constant term is rational.  Write
$\mathrm{PA}_{\mathbf R}(\mathbf X)$ for its real scalar extension, consisting
of continuous functions which are affine with arbitrary real coefficients on
the faces of a finite rational polyhedral structure.  Write
$\mathrm{PAPC}_{\mathbf Z}$ and
$\mathrm{PC}_{\mathrm{reg},\mathbf Z}$ for the source's literal classes, and
define
$\mathrm{PAPC}_{\mathbf R}(\mathbf X,\gamma)$ and
$\mathrm{PC}_{\mathrm{reg},\mathbf R}(\mathbf X,\gamma)$ from
$\mathrm{PA}_{\mathbf R}$ exactly as the source defines $\mathrm{PAPC}$ and
$\mathrm{PC}_{\mathrm{reg}}$ from $\mathrm{PA}_{\mathbf Z}$.  Thus the latter
space consists of uniform limits of decreasing sequences in
$\mathrm{PAPC}_{\mathbf R}(\mathbf X,\gamma)$.  We say that
$(\mathbf X,\gamma)$ has real-polyhedral regularity if
$P_\gamma(u)\in\mathrm{PC}_{\mathrm{reg},\mathbf R}(\mathbf X,\gamma)$ for
every bounded $u\in\mathrm{PA}_{\mathbf R}(\mathbf X)$.

\begin{theorem}
\label{thm:frontier-polyhedral-one-dimensional-real-regularity}
Let $\mathbf X$ be a one-dimensional balanced rational polyhedral space
admitting a finite structure, and let
$\gamma\in\mathrm{PA}_{\mathbf R}(\mathbf X)\cap
\mathrm{PC}(\mathbf X)$ be positive at infinity.  If
$u\in\mathrm{PA}_{\mathbf R,b}(\mathbf X)$ and $\Pi$ is a common finite
rational polyhedral structure on which $\gamma$ and $u$ are affine, then
\[
 P_\gamma(u)\in
 \mathrm{PAPC}_{\mathbf R}(\mathbf X,\gamma)
\]
and $P_\gamma(u)$ is affine on every edge of $\Pi$.  In particular,
$(\mathbf X,\gamma)$ has real-polyhedral regularity.
\end{theorem}

\begin{proposition}
\label{prop:frontier-polyhedral-integral-regularity-counterexample}
For the literal integral-slope class, regularity can fail even on the complete
affine line.  More precisely, let $\mathbf X=\mathbf R$,
$\gamma(x)=|x|$, and
\[
 u(x)=
 \begin{cases}
  0,&x\leq1,\\
  1-x,&1\leq x\leq2,\\
  -1,&x\geq2.
 \end{cases}
\]
Then $\gamma\in\mathrm{PAPC}_{\mathbf Z}(\mathbf R)$ is positive at infinity,
$u\in\mathrm{PA}_{\mathbf Z,b}(\mathbf R)$, and
\[
 P_\gamma(u)(x)=
 \begin{cases}
  -x,&x\leq0,\\
  x/2,&0\leq x\leq2,\\
  x-1,&x\geq2.
 \end{cases}
\]
Moreover every
$\psi\in\mathrm{PAPC}_{\mathbf Z}(\mathbf R,\gamma)$ satisfies
\[
 \lVert\psi-P_\gamma(u)\rVert_{L^\infty(\mathbf R)}\geq\frac14.
\]
Consequently
$P_\gamma(u)\notin\mathrm{PC}_{\mathrm{reg},\mathbf Z}
(\mathbf R,\gamma)$, and the literal integral-slope pair
$(\mathbf R,|\cdot|)$ does not have the regularity property of
\cite[Theorem~6.4 and Definition~6.5]{BMPS26}.
\end{proposition}

Thus the one-dimensional assertion in
\cite[Proposition~6.13]{BMPS26} is valid after real scalar extension but is
incompatible, as written, with the integral-slope convention in the source's
Definition~2.28.  The theorem and proposition make no higher-dimensional
regularity claim.








\begin{proposition}\label{prop:frontier-polyhedral-bounded-reference-shift}
Let \(\gamma\) and \(\widetilde\gamma\) be piecewise-affine
polyhedrally convex references on the same pure-dimensional balanced
rational polyhedral space \(\mathbf X\), both positive at infinity, and put
\[
             \delta:=\widetilde\gamma-\gamma\in
             \mathrm{PA}_{b}(\mathbf X).
\]
Then, as classes of functions on \(\mathbf X\),
\[
\begin{gathered}
 \mathrm{PC}(\mathbf X,\widetilde\gamma)
 =\mathrm{PC}(\mathbf X,\gamma),\qquad
 \mathrm{PAPC}(\mathbf X,\widetilde\gamma)
 =\mathrm{PAPC}(\mathbf X,\gamma),\\
 \mathrm{PC}_{\mathrm{reg}}(\mathbf X,\widetilde\gamma)
 =\mathrm{PC}_{\mathrm{reg}}(\mathbf X,\gamma).
\end{gathered}
\]
For every \(u\in\mathrm{PA}_{b}(\mathbf X)\), their envelopes satisfy
\[
 P_{\widetilde\gamma}(u)=P_\gamma(u+\delta).
 \tag{\(\star\)}\label{eq:frontier-polyhedral-bounded-shift-envelope}
\]
Consequently regularity holds for \((\mathbf X,\widetilde\gamma)\) if and
only if it holds for \((\mathbf X,\gamma)\); when these equivalent
regularity conditions hold, orthogonality holds for one reference if and
only if it holds for the other.

The same assertions hold in the real scalar extension after replacing
\(\mathrm{PA}_{b}\), \(\mathrm{PAPC}\), and
\(\mathrm{PC}_{\mathrm{reg}}\) throughout by their
\(\mathbf R\)-subscripted versions and assuming
\(\delta\in\mathrm{PA}_{\mathbf R,b}(\mathbf X)\).
\end{proposition}





\begin{theorem}
\label{thm:frontier-polyhedral-one-vertex-integral-nonregularity}
Let
\[
                  \mathbf X=\bigcup_{i=1}^r\mathbf R_{\geq0}u_i
\]
be a one-dimensional balanced rational fan with one vertex and primitive ray
directions $u_i$.  Let
$\gamma\in\mathrm{PAPC}_{\mathbf Z}(\mathbf X)$ be affine on every ray and
positive at infinity, and write
\[
                   \gamma(tu_i)=\gamma(0)+g_i t,
                   \qquad g_i\in\mathbf Z_{>0}.
\]
Fix a ray $j$ and an integer $N\geq2$.  Define the bounded literal
piecewise-affine function $u_{j,N}$ to vanish off that ray and, on the ray,
by
\[
 u_{j,N}(tu_j)=
 \begin{cases}
  0,&0\leq t\leq N-1,\\
  N-1-t,&N-1\leq t\leq N,\\
  -1,&t\geq N.
 \end{cases}
\]
Then
\begin{equation}
\label{eq:frontier-polyhedral-one-vertex-fractional-envelope}
 P_\gamma(u_{j,N})(tu_i)=
 \begin{cases}
  \gamma(tu_i),&i\neq j,\\
  \gamma(tu_j)-t/N,&i=j,\ 0\leq t\leq N,\\
  \gamma(tu_j)-1,&i=j,\ t\geq N.
 \end{cases}
\end{equation}
For every
$\psi\in\mathrm{PAPC}_{\mathbf Z}(\mathbf X,\gamma)$ one has
\[
 \left\lVert\psi-P_\gamma(u_{j,N})\right\rVert_{L^\infty(\mathbf X)}
 \geq\frac{N-1}{2N}.
\]
In particular, $(\mathbf X,\gamma)$ does not have the literal
integral-slope regularity property.
\end{theorem}

\begin{corollary}
\label{cor:frontier-polyhedral-one-vertex-integral-nonregularity}
Let $\mathbf X$ be as in
\Cref{thm:frontier-polyhedral-one-vertex-integral-nonregularity}, and let
$\gamma\in\mathrm{PAPC}_{\mathbf Z}(\mathbf X)$ be any positive-at-infinity
reference defined on a finite rational subdivision.  Then
$(\mathbf X,\gamma)$ does not have the literal integral-slope regularity
property.  Consequently no such pair lies in the domain of the literal
orthogonality property of \cite[Definition~6.9]{BMPS26}, whose definition
assumes regularity.
\end{corollary}

The corollary concerns the source's literal coefficient convention.  It does
not contradict real-polyhedral regularity or the real-polyhedral
orthogonality classification proved above, and it does not assign a
polyhedral Monge--Amp\`ere measure to the nonregularizable envelope in
\eqref{eq:frontier-polyhedral-one-vertex-fractional-envelope}.








Thus the P05.1 classification factors through bounded piecewise-affine
reference shifts within either fixed coefficient category.  On every
one-vertex balanced rational fan, the real scalar category has regularity,
whereas the literal integral-slope category never does.  In particular, the
source's conical regularization statement cannot be promoted to regularity
for arbitrary bounded integral-slope obstacles: that statement concerns
conical functions, while
\Cref{thm:frontier-polyhedral-one-vertex-integral-nonregularity} produces a
nonregularizable envelope for every positive-at-infinity literal reference.

The author-supplied one-vertex note is now explicitly withdrawn
\cite{SuppliedP05220260925}.  Its withdrawal records the same coefficient
conflict: the stated positive theorem used the literal integral-slope class
but invoked a real-affine envelope.  The criterion below is therefore not
cited from that note.  It is a book-local theorem in the explicit real
scalar extension, proved by the displayed contact and obstruction arguments
and using the local convexity test from
\cite[Definitions~3.2--3.3]{BMPS26} and the real-linear one-dimensional
Laplacian formula from
\cite[Definition~6.15 and Proposition~6.17]{BMPS26}.



Use the real scalar extension
$\mathrm{PA}_{\mathbf R}(\mathbf X)$ fixed in
\Cref{thm:frontier-polyhedral-one-dimensional-real-regularity}.  A
polyhedral convex combination is a finite barycentric combination
$x=\sum_j a_jx_j$ for which, on some common refinement, there are faces
$\tau$ and $\sigma_j$ with $x\in\tau$, $x_j\in\sigma_j$, and
$\tau\preceq\sigma_j$ for every $j$.  The envelope is
\[
 P_\gamma(u):=\sup\{\psi:\ \psi\leq\gamma+u,\quad
   \psi\in\mathrm{PC}(\mathbf X),\quad
   \psi\leq\gamma+O(1)\}.
\]
The preceding theorem proves, independently of the withdrawn supplied note, that
$P_\gamma(u)\in\mathrm{PA}_{\mathbf R}(\mathbf X)\cap
\mathrm{PC}(\mathbf X)$ on the same finite structure whenever
$u\in\mathrm{PA}_{\mathbf R,b}(\mathbf X)$.  For such a function, write
$\MA_{\mathrm{poly},\mathbf R}(f):=\Delta f$, where the
one-dimensional Laplacian uses the primitive integral-affine coordinate
and the balancing weight on each edge.  This is the real-linear extension
of the source's one-dimensional intersection formula
\cite[Definition~6.15 and Proposition~6.17]{BMPS26}.
We say that $(\mathbf X,\gamma)$ has real-polyhedral orthogonality if, for
every $u\in\mathrm{PA}_{\mathbf R,b}(\mathbf X)$,
\[
 \int_{\mathbf X}
 \bigl(P_\gamma(u)-(\gamma+u)\bigr)
 \MA_{\mathrm{poly},\mathbf R}(P_\gamma(u))=0.
\]

\begin{lemma}
\label{lem:frontier-polyhedral-one-vertex-convexity}
Let
\[
 \mathbf X=\bigcup_{i=1}^{r}\mathbf R_{\geq0}u_i
 \subset N_{\mathbf R}
\]
be the union of distinct rational rays with primitive generators $u_i$, put
\[
 K:=\left\{\lambda\in\mathbf R_{\geq0}^{r}:\
       \sum_{i=1}^{r}\lambda_i u_i=0\right\},
\]
and let $f$ be continuous and real-piecewise-affine on every ray.  If $s_i$
is the outgoing slope of $f$ at the origin along $u_i$, then $f$ is
polyhedrally convex if and only if every ray restriction is convex and
\[
                    \sum_i\lambda_i s_i\geq0
                    \quad\text{for every }\lambda\in K.
\]
If the ray $\mathbf R_{\geq0}u_i$ has balancing weight $w_i$, then
\[
             \MA_{\mathrm{poly},\mathbf R}(f)\{0\}=\sum_iw_i s_i.
\]
\end{lemma}

\begin{theorem}
\label{thm:frontier-polyhedral-one-vertex-orthogonality}
Let
$
 \mathbf X=\bigcup_{i=1}^{r}\mathbf R_{\geq0}u_i
$
be a one-dimensional balanced rational fan with distinct primitive ray
directions $u_i$ and positive rational weights $w_i$ such that
$
 \sum_iw_i u_i=0
$.
Let $\gamma\in\mathrm{PA}_{\mathbf R}(\mathbf X)$ be polyhedrally convex,
affine on every ray, with outgoing slopes $g_i>0$.  Then
$(\mathbf X,\gamma)$ has real-polyhedral orthogonality if and only if
\[
 \dim K=1,
 \qquad
 K=\left\{\lambda\in\mathbf R_{\geq0}^{r}:\
       \sum_i\lambda_i u_i=0\right\}.
\]
Equivalently,
$
 \dim_{\mathbf R}\operatorname{span}\{u_1,\ldots,u_r\}=r-1
$.
\end{theorem}

\begin{corollary}
\label{cor:frontier-polyhedral-one-vertex-bounded-shift}
In the setting of
\Cref{thm:frontier-polyhedral-one-vertex-orthogonality}, let $\gamma$ be
an arbitrary real-piecewise-affine polyhedrally convex reference positive at
infinity, without assuming that it is affine on each whole ray.  Thus
$\gamma$ is defined on a finite rational subdivision of the fan.  Then
$(\mathbf X,\gamma)$ has real-polyhedral orthogonality if and only if
$\dim K=1$.
\end{corollary}

\begin{example}
\label{ex:frontier-polyhedral-nonsaturated-tripod}
In $N=\mathbf Z^2$, take
\[
 u_1=(1,0),\qquad u_2=(1,2),\qquad u_3=(-3,-2),
 \qquad (w_1,w_2,w_3)=(2,1,1).
\]
The ray directions are primitive, their weighted sum is zero, and their
real span has rank two.  The lattice they generate has index two in $N$,
since the absolute values of the three pairwise determinants are $2,2,4$.
Thus the vertex is not polyhedrally smooth in the lattice sense, whereas
\Cref{cor:frontier-polyhedral-one-vertex-bounded-shift} gives
real-polyhedral orthogonality for every real-piecewise-affine
polyhedrally convex reference positive at infinity on this fan.
\end{example}










The same contact perturbation is local and therefore survives bounded
edges.  Applying it at every vertex whose cone of nonnegative balancing
relations is one-dimensional gives the following graph-level extension.



Let $\mathbf X$ be a one-dimensional balanced rational polyhedral space,
and fix a finite rational polyhedral structure $\Pi$ carrying its balancing
cycle.  For a vertex $v\in\Pi^{(0)}$, write $E(v)$ for the incident edges,
$u_{e/v}$ for the primitive outgoing integral direction of $e$ at $v$, and
$w_e>0$ for the balancing weight of $e$.  Put
\[
 K_v:=\left\{\lambda\in\mathbf R_{\geq0}^{E(v)}:\
             \sum_{e\in E(v)}\lambda_eu_{e/v}=0\right\}.
\]
The local weight vector $w_v=(w_e)_{e\in E(v)}$ belongs to the relative
interior of $K_v$.

\begin{theorem}
\label{thm:frontier-polyhedral-local-rank-one-graph}
Assume that $\dim K_v=1$ for every vertex $v\in\Pi^{(0)}$.  Let
$\gamma\in\mathrm{PA}_{\mathbf R}(\mathbf X)\cap
\mathrm{PC}(\mathbf X)$ be positive at infinity.  Then
$(\mathbf X,\gamma)$ has real-polyhedral orthogonality: for every
$u\in\mathrm{PA}_{\mathbf R,b}(\mathbf X)$,
\[
 \int_{\mathbf X}
 \bigl(P_\gamma(u)-(\gamma+u)\bigr)
 \MA_{\mathrm{poly},\mathbf R}(P_\gamma(u))=0.
\]
The conclusion allows bounded edges and cycles in the underlying graph.
It depends only on the local cones $K_v$ and is unchanged by subdividing
$\Pi$.
\end{theorem}

\begin{corollary}
\label{cor:frontier-polyhedral-trivalent-graph}
Every one-dimensional balanced rational polyhedral space admitting a finite
structure whose vertices have valence at most three satisfies the conclusion
of \Cref{thm:frontier-polyhedral-local-rank-one-graph} for every
real-piecewise-affine polyhedrally convex reference positive at infinity.
In particular, lattice saturation and the absence of bounded cycles are not
required.
\end{corollary}

The theorem is a sufficient criterion, not a converse.  At a vertex with
$\dim K_v>1$, the one-vertex construction gives a local noncontact mass
mechanism, but extending that germ through all bounded edges while retaining
global polyhedral convexity and the prescribed growth is a separate
realizability problem.








The remaining converse is a global realization question: a local
higher-rank obstruction must be extended through the graph without changing
the prescribed growth.  Positivity at infinity supplies exactly that
extension.  A sufficiently slow ambient rational-polyhedral norm can be
raised above the reference near any chosen vertex and then cut off by taking
the maximum with the reference.  Rational chordal interpolation retains its
vertex slopes and restores the required rational subdivision.



Continue with the finite one-dimensional balanced rational polyhedral
structure $\Pi$ and the local relation cones $K_v$ above.  We may insert one
bivalent vertex in every component without a vertex.  This does not change
any original relation cone or the orthogonality problem.

\begin{definition}
\label{def:frontier-polyhedral-mass-accessible}
Let $\gamma\in\mathrm{PA}_{\mathbf R}(\mathbf X)\cap
\mathrm{PC}(\mathbf X)$.  A vertex $v\in\Pi^{(0)}$ is
$\gamma$-mass-accessible if there is a function
\[
 \varphi\in\mathrm{PA}_{\mathbf R}(\mathbf X)\cap
 \mathrm{PC}(\mathbf X),
 \qquad
 \varphi-\gamma\in\mathrm{PA}_{\mathbf R,b}(\mathbf X),
\]
such that
\[
 \MA_{\mathrm{poly},\mathbf R}(\varphi)\{v\}
 =\sum_{e\in E(v)}w_es_{v,e}(\varphi)>0.
\]
\end{definition}

\begin{lemma}
\label{lem:frontier-polyhedral-ambient-vertex-accessibility}
Let $\gamma$ be positive at infinity.  Then every vertex of $\Pi$ is
$\gamma$-mass-accessible.
\end{lemma}

\begin{lemma}
\label{lem:frontier-polyhedral-accessible-mass-obstruction}
Suppose that $\dim K_v>1$ and that $v$ is
$\gamma$-mass-accessible.  Then $(\mathbf X,\gamma)$ does not have
real-polyhedral orthogonality.  More precisely, there is a bounded
real-piecewise-affine obstacle $u$ such that
\[
 \int_{\mathbf X}
 \bigl(P_\gamma(u)-(\gamma+u)\bigr)
 \MA_{\mathrm{poly},\mathbf R}(P_\gamma(u))<0.
\]
\end{lemma}

\begin{theorem}
\label{thm:frontier-polyhedral-graph-orthogonality-classification}
Let $\mathbf X$ be a one-dimensional balanced rational polyhedral space
admitting a finite structure $\Pi$, and let
$\gamma\in\mathrm{PA}_{\mathbf R}(\mathbf X)\cap
\mathrm{PC}(\mathbf X)$ be positive at infinity.  Then
$(\mathbf X,\gamma)$ has real-polyhedral orthogonality if and only if
\[
                         \dim K_v=1
                         \qquad\text{for every }v\in\Pi^{(0)}.
\]
Thus the repaired one-dimensional classification is local at the vertices:
bounded edges, graph cycles, and lattice saturation introduce no additional
condition.
\end{theorem}

Positivity at infinity enters the converse only through
\Cref{lem:frontier-polyhedral-ambient-vertex-accessibility}: it permits an
ambient polyhedral norm with strictly smaller recession slope on every
unbounded edge.  The statement neither repairs the literal integral-slope
class nor addresses higher-dimensional balanced spaces.










The same local relation cones also sharpen the negative result in the
literal coefficient category.  The one-vertex construction used positivity
of each ray slope.  In a graph, the exact replacement is strict positivity
of the whole slope vector on every nonzero local balancing relation.



Let $\mathbf X$ be a one-dimensional balanced rational polyhedral space
with a finite rational polyhedral structure $\Pi$.  For a vertex
$v\in\Pi^{(0)}$, retain the notation
\[
 K_v=\left\{\lambda\in\mathbf R_{\geq0}^{E(v)}:\
             \sum_{e\in E(v)}\lambda_eu_{e/v}=0\right\}.
\]

\begin{theorem}
\label{thm:frontier-polyhedral-strict-vertex-integral-nonregularity}
Let $\gamma\in\mathrm{PAPC}_{\mathbf Z}(\mathbf X)$ be positive at
infinity, and refine $\Pi$ so that $\gamma$ is affine on every edge.  Fix a
vertex $v$, and let $s_e\in\mathbf Z$ be the outgoing slope of $\gamma$ on
the edge germ $e\in E(v)$.  Suppose that
\begin{equation}
             \sum_{e\in E(v)}\lambda_es_e>0
             \qquad\text{for every }\lambda\in K_v\setminus\{0\}.
 \label{eq:frontier-polyhedral-strict-vertex}
\end{equation}
Then $(\mathbf X,\gamma)$ does not have the literal integral-slope
regularity property.

More precisely, choose a positive rational number $R$ such that the closed
edge-germ segments
\[
       x_e([0,R]),\qquad x_e(t):=v+tu_{e/v},
\]
have pairwise disjoint interiors, contain no other vertex, and lie in the
affine pieces of $\gamma$.  There are arbitrarily large integers $N\geq2$
such that
\begin{equation}
 \sum_e\lambda_e\left(s_e-\frac1N\right)\geq0
                 \qquad\text{for every }\lambda\in K_v.
 \label{eq:frontier-polyhedral-strict-vertex-choice}
\end{equation}
For any such $N$, put $c:=R/N$ and define the bounded literal
piecewise-affine obstacle $u_{v,R,N}$ by
\[
 u_{v,R,N}(x_e(t))=
 \begin{cases}
  c,&0\leq t\leq R-c,\\
  R-t,&R-c\leq t\leq R,
 \end{cases}
 \qquad
 u_{v,R,N}=0\quad\text{off the closed star}.
\]
Its envelope is
\begin{equation}
 \begin{aligned}
 P_\gamma(u_{v,R,N})(x_e(t))
   &=\gamma(x_e(t))+\frac{R-t}{N}
       &&(0\leq t\leq R),\\
 P_\gamma(u_{v,R,N})&=\gamma
       &&\text{off the closed star}.
 \end{aligned}
 \label{eq:frontier-polyhedral-strict-vertex-envelope}
\end{equation}
Every literal integral-slope approximant satisfies
\begin{equation}
 \left\lVert\psi-P_\gamma(u_{v,R,N})\right\rVert_{L^\infty(\mathbf X)}
 \geq \frac{(N-1)R}{2N^2},
 \qquad
 \psi\in\mathrm{PAPC}_{\mathbf Z}(\mathbf X,\gamma).
 \label{eq:frontier-polyhedral-strict-vertex-separation}
\end{equation}
\end{theorem}

\begin{corollary}
\label{cor:frontier-polyhedral-rank-one-graph-integral-nonregularity}
Assume that $\dim K_v=1$ at every vertex of a finite rational polyhedral
structure on $\mathbf X$.  Let
$\gamma\in\mathrm{PAPC}_{\mathbf Z}(\mathbf X)$ be positive at infinity
and have positive degree.  Then $(\mathbf X,\gamma)$ does not have the
literal integral-slope regularity property.
\end{corollary}

Thus, on the graph class where real-polyhedral orthogonality holds by
\Cref{thm:frontier-polyhedral-graph-orthogonality-classification}, every
positive-degree literal reference instead fails the regularity hypothesis
required by the source's literal orthogonality definition.  The two
statements concern different coefficient categories.








The degree-zero case has the opposite behaviour.  Positivity at infinity
makes every unbounded end contribute strictly positive degree.  When the
degree vanishes, the graph is compact, and balancing at a maximum vertex
forces every polyhedrally convex function to be componentwise constant.
In the globally embedded source category, the compact alternative is
vacuous in positive dimension; the source-input reduction below explains
this category boundary.



\begin{theorem}
\label{thm:frontier-polyhedral-zero-degree-integral-regularity}
Let $\mathbf X$ be a one-dimensional balanced rational polyhedral space
admitting a finite structure, and let
$\gamma\in\mathrm{PAPC}_{\mathbf Z}(\mathbf X)$ be positive at infinity.
Then the following conditions are equivalent:
\begin{enumerate}
\item $\deg(\gamma)=0$;
\item $\mathbf X$ is compact.
\end{enumerate}
When these conditions hold, every function in
$\mathrm{PC}(\mathbf X,\gamma)$ is constant on each connected component.
For every bounded literal piecewise-affine function
$u\in\mathrm{PA}_{\mathbf Z,b}(\mathbf X)$ and every connected component
$C$ of $\mathbf X$, one has
\begin{equation}
 P_\gamma(u)|_C
   =\min_{x\in C}\bigl(\gamma(x)+u(x)\bigr).
 \label{eq:frontier-polyhedral-zero-degree-envelope}
\end{equation}
The right-hand side is a rational constant.  In particular,
$P_\gamma(u)\in\mathrm{PAPC}_{\mathbf Z}(\mathbf X,\gamma)$ and
$(\mathbf X,\gamma)$ has the literal integral-slope regularity property.
\end{theorem}

Thus the zero-degree branch of the one-dimensional literal problem is
regular.  Together with
\Cref{thm:frontier-polyhedral-strict-vertex-integral-nonregularity}, this
leaves only positive-degree graphs for which every positive-mass vertex has
a higher-dimensional local relation cone and the reference slope vector lies
on a proper face of its dual cone.






The remaining positive-degree boundary admits a global obstruction which
does not depend on a strict reference vertex.  The point is to place a
rational convex combination strictly between the reference and a second
literal convex function.  Edgewise integral-slope concave roofs then give a
literal obstacle whose convex envelope is that combination.



\begin{theorem}
\label{thm:frontier-polyhedral-one-dimensional-literal-classification}
Let $\mathbf X$ be a one-dimensional balanced rational polyhedral space
admitting a finite structure, and let
$\gamma\in\mathrm{PAPC}_{\mathbf Z}(\mathbf X)$ be positive at infinity.
The following are equivalent:
\begin{enumerate}
\item $(\mathbf X,\gamma)$ has the literal integral-slope regularity
property;
\item $\deg(\gamma)=0$;
\item $\mathbf X$ is compact.
\end{enumerate}

More precisely, suppose that $\deg(\gamma)>0$.  Choose a rational point $v$
in a noncompact connected component, a rational number $C>\gamma(v)$, and
put
\[
                 \varphi:=\max\{\gamma,C\}.
\]
There is a finite rational subdivision on which $\gamma$ and $\varphi$ are
affine and an oriented bounded edge cell $e$ on which
$\varphi-\gamma$ has a nonzero integral slope $k_e$.  For every integer
$N>|k_e|$, there is a bounded literal piecewise-affine obstacle $u_N$ such
that
\begin{equation}
 P_\gamma(u_N)
 =\left(1-\frac1N\right)\gamma+\frac1N\varphi.
 \label{eq:frontier-polyhedral-literal-classification-envelope}
\end{equation}
If $L_e$ is the integral-affine length of $e$ and $q_e$ is the slope of the
right-hand side of
\eqref{eq:frontier-polyhedral-literal-classification-envelope} on $e$, then
every literal integral-slope polyhedrally convex function $\psi$ satisfies
\begin{equation}
 \left\lVert\psi-P_\gamma(u_N)\right\rVert_{L^\infty(\mathbf X)}
 \geq \frac{L_e}{4}\operatorname{dist}(q_e,\mathbf Z)>0.
 \label{eq:frontier-polyhedral-literal-classification-separation}
\end{equation}
Thus the failure in positive degree is independent of the dimensions and
faces of the local relation cones.
\end{theorem}






In higher dimension the edgewise concave roof has no canonical analogue, but
the vertex-contact mechanism does.  For each vertex, an integral
piecewise-affine lattice norm gives a cone majorizing the rational gap.
Taking the minimum of these cones and the original integral gap preserves
the prescribed vertex values while forcing the obstacle to vanish on every
unbounded cell.  This removes dimension, balancing, and local-relation
hypotheses from the noncompact literal obstruction.



\begin{theorem}
\label{thm:frontier-polyhedral-noncompact-literal-nonregularity}
Let $\mathbf X\subset N_{\mathbf R}$ be a noncompact rational polyhedral
space admitting a finite structure, and let
$\gamma\in\mathrm{PAPC}_{\mathbf Z}(\mathbf X)$ be positive at infinity.
Then $(\mathbf X,\gamma)$ does not have the literal integral-slope
regularity property.

More precisely, choose a rational point $x$ in a noncompact connected
component, a rational number $C>\gamma(x)$, and put
\[
                 \varphi:=\max\{\gamma,C\},
                 \qquad D:=\varphi-\gamma.
\]
There are a finite rational subdivision $\Pi$, a bounded positive-dimensional
cell $\sigma\in\Pi$, a primitive vector $w\in N_\sigma$, and a rational
segment
\[
             I=x_0+[0,L]w\subset\operatorname{relint}(\sigma)
\]
such that $D$ is affine on $\sigma$ and its linear part $k\in M_\sigma$
satisfies $k(w)\neq0$.  For every integer $N>|k(w)|$, there is a bounded
literal piecewise-affine function $u_N$ for which
\begin{equation}
 P_\gamma(u_N)
 =\gamma+\frac{D}{N}
 =\left(1-\frac1N\right)\gamma+\frac1N\varphi.
 \label{eq:frontier-polyhedral-noncompact-literal-envelope}
\end{equation}
If $q$ is the slope of this envelope along the primitive direction $w$ on
$I$, then every literal integral-slope polyhedrally convex function $\psi$
satisfies
\begin{equation}
 \left\lVert\psi-P_\gamma(u_N)\right\rVert_{L^\infty(\mathbf X)}
 \geq \frac{L}{4}\operatorname{dist}(q,\mathbf Z)>0.
 \label{eq:frontier-polyhedral-noncompact-literal-separation}
\end{equation}
In particular, the literal complete-affine regularity assertion fails on
$N_{\mathbf R}$ in every positive dimension.
\end{theorem}






There is an elementary source consequence that fixes the compact branch.
The space in \cite[Definitions~2.2 and~2.4]{BMPS26} is the support of a
finite complex embedded in one vector space, with the strictly positive
fundamental cycle of \cite[Definition~2.27]{BMPS26}.
\cite[Proposition~3.11 and Lemma~3.25]{BMPS26} provide a finite
simplicial refinement carrying a strictly positive one-cycle whenever
the dimension is positive. If the support were compact, all its edges
would be bounded. Fix a Euclidean inner product on the ambient vector
space. Write an edge as $q-p=\lambda_e u_e$, where
$\lambda_e>0$ and $u_e$ is its primitive direction, and let $c_e>0$
be its cycle weight. Balancing at the vertices gives
\[
 0=\sum_v\left\langle v,\sum_{e\ni v}c_eu_{e,v}\right\rangle
   =-\sum_{e=[p,q]}c_e\lambda_e\lVert u_e\rVert^2<0,
\]
a contradiction. Thus every nonempty positive-dimensional balanced
space in this source category is noncompact. This is a direct formal
consequence of the named source inputs.
The compactification of \cite[Definition~2.13]{BMPS26} and compact
tropical spaces defined only by local affine charts have different
hypotheses; compactness of either does not supply a compact support here.

This completely classifies the repaired real-piecewise-affine version of
P05.2 on every one-dimensional balanced rational polyhedral space admitting
a finite structure.  Orthogonality holds exactly when every local cone of
nonnegative balancing relations is one-dimensional.  The theorem allows
bounded edges and cycles and shows that lattice saturation is irrelevant.
The conditional degree-zero graph statements retain their mathematical
form, but their compact branch is empty in the embedded source category.
Every positive-degree graph reference fails literal regularity,
independently of the local relation cones.  More generally, literal
regularity fails on every noncompact finite rational polyhedral space in
every dimension, whether or not it is balanced; in particular, the
source's literal complete-affine assertion fails in every positive
dimension. Combined with the preceding source consequence, this excludes
literal regularity for every nonempty positive-dimensional finite balanced
pair in that category. The earlier strict-vertex theorem remains a local quantitative
form of this obstruction.  Thus every graph on which real-polyhedral
orthogonality holds is negative for literal regularity when the literal
degree is positive, but the literal failure is not confined to that graph
class.  The
one-vertex orthogonality theorem and its bounded-reference-shift corollary
are book-local real-scalar results, not results cited from the withdrawn
note.  The classification also explains
\Cref{ex:frontier-polyhedral-cross-orthogonality}: the four-ray cross has a
two-dimensional cone of nonnegative balancing relations.  Conversely,
\Cref{ex:frontier-polyhedral-nonsaturated-tripod} shows that lattice
smoothness is sufficient but not necessary in this class.  The
higher-dimensional real-PA classification, arbitrary nonseparable product
obstacles, and the source
coefficient-convention repair remain open, so P05.1 and P05.2 remain active.

The following elementary convex-envelope fact is the missing contact
mechanism on a complete affine space.


\begin{lemma}\label{lem:frontier-polyhedral-affine-contact}
Let \(f:\mathbf R^d\to\mathbf R\) be continuous and coercive, and let
\(P\) be its greatest lower-semicontinuous convex minorant.  Suppose that
\(P\) is polyhedral.  If \(\partial P(x)\) has positive
\(d\)-dimensional volume, then \(P(x)=f(x)\).
\end{lemma}



\begin{theorem}\label{thm:frontier-polyhedral-affine-orthogonality}
Let \(N\) have rank \(d\), let
\(\mathbf X=N_{\mathbf R}\) carry the constant balancing weight
\(c>0\), and let
\(\gamma\in\mathrm{PA}_{\mathbf R}(N_{\mathbf R})\) be polyhedrally
convex and positive at infinity.  Then
\((\mathbf X,\gamma)\) has real-polyhedral regularity in every dimension.
Moreover, for every \(u\in\mathrm{PA}_{\mathbf R,b}(N_{\mathbf R})\),
writing \(P=P_\gamma(u)\) and
\[
             \MA_{\mathrm{poly},\mathbf R}(P)
             :=d!c\,\MA_{\mathbf R}(P),
\]
one has
\[
 \int_{N_{\mathbf R}}
 \bigl(P-(\gamma+u)\bigr)\,
 \MA_{\mathrm{poly},\mathbf R}(P)=0.
\]
No strict-convexity hypothesis beyond positivity at infinity is needed.
\end{theorem}



The epigraph proof in \Cref{thm:frontier-polyhedral-affine-orthogonality}
is the real-affine argument underlying the complete-space discussion in
\cite[Remark~6.6]{BMPS26}; the contact deduction is book-local.  It cannot
be read as a theorem in the literal integral-slope category:
\Cref{prop:frontier-polyhedral-integral-regularity-counterexample} disproves
that interpretation already for \(N_{\mathbf R}=\mathbf R\).



Let $r\geq1$.  For $1\leq i\leq r$, let $\mathbf X_i$ be a rational
polyhedral space admitting a finite structure, and write
$p_i:\mathbf X:=\prod_i\mathbf X_i\to\mathbf X_i$ for the projection.  Give
$\mathbf X$ the product polyhedral structure.  Let
$\gamma_i\in\mathrm{PA}_{\mathbf R}(\mathbf X_i)\cap
\mathrm{PC}(\mathbf X_i)$ and
$u_i\in\mathrm{PA}_{\mathbf R,b}(\mathbf X_i)$, and put
\[
 a_i:=\gamma_i+u_i,
 \qquad
 P_i:=P_{\gamma_i}(u_i),
 \qquad
 \gamma:=\sum_{i=1}^r p_i^*\gamma_i,
 \qquad
 u:=\sum_{i=1}^r p_i^*u_i.
\]

\begin{theorem}
\label{thm:frontier-polyhedral-separable-product-envelope}
With the notation above, assume that
\[
 P_i\in
 \mathrm{PAPC}_{\mathbf R}(\mathbf X_i,\gamma_i)
 \qquad(1\leq i\leq r).
\]
Then
\[
 P_\gamma(u)=\sum_{i=1}^r p_i^*P_i.
 \tag{*}\label{eq:frontier-polyhedral-separable-product-envelope}
\]
In particular, the envelope of the separable obstacle $u$ belongs to
$\mathrm{PAPC}_{\mathbf R}(\mathbf X,\gamma)$.
\end{theorem}

\begin{corollary}
\label{cor:frontier-polyhedral-graph-affine-separable-regularity}
Assume in addition that every $\mathbf X_i$ is either a complete affine
space $N_{i,\mathbf R}$ or a one-dimensional balanced rational polyhedral
space admitting a finite structure, and that every $\gamma_i$ is positive
at infinity.  The product carries its product balancing cycle, the sum
reference $\gamma$ is positive at infinity, and
\eqref{eq:frontier-polyhedral-separable-product-envelope} holds for every
tuple of bounded real-piecewise-affine obstacles $u_i$.
Thus products of finite balanced graphs and complete affine spaces have
real-piecewise-affine envelopes for all separable bounded obstacles.  This
includes higher-dimensional products with branched singular factors.
\end{corollary}

The corollary concerns the separable obstacle class
$u=\sum_i p_i^*u_i$.  It does not assert the regularity property for an
arbitrary bounded real-piecewise-affine function on the product.








The next example shows why completeness of the ambient convex geometry
cannot be replaced by strict positivity of the reference alone.

\begin{example}\label{ex:frontier-polyhedral-cross-orthogonality}
Let
\[
 \mathbf X=(\mathbf R\times\{0\})\cup(\{0\}\times\mathbf R)
\]
with all four ray weights equal to one, and put
\(\gamma(x,y)=|x|+|y|\) on \(\mathbf X\).  This reference is positive at
infinity and strictly polyhedrally convex on the four-ray complex.  For
\(\varepsilon>0\), define the obstacle
\[
 f(x,0)=|x|+\varepsilon,
 \qquad
 f(0,y)=\bigl||y|-\varepsilon\bigr|,
 \qquad u=f-\gamma.
\]
Then \(u\) is bounded and piecewise affine, and
\[
 P_\gamma(u)(x,0)=|x|,
 \qquad
 P_\gamma(u)(0,y)=\max\{0,|y|-\varepsilon\}.
\]
\end{example}

Indeed, the displayed candidate is the restriction of the ambient convex
function
\(
 \max\{|x|,|y|-\varepsilon,0\}
\), hence is polyhedrally convex and lies below \(f\).  Conversely, if
\(q\leq f\) is an admissible polyhedrally convex function, convexity between
\((0,\varepsilon)\) and \((0,-\varepsilon)\) gives \(q(0,0)\leq0\).
On either horizontal ray, convexity and the asymptotic slope bound one give
\(q(x,0)\leq|x|\).  On the positive vertical ray, convexity first gives
\(q(0,y)\leq0\) for \(0\leq y\leq\varepsilon\), and then the same slope
bound gives \(q(0,y)\leq y-\varepsilon\) for
\(y\geq\varepsilon\); the negative ray is identical.  This proves the
envelope formula.  Its outgoing slopes at the origin have sum two, so
\[
 \MA_{\mathrm{poly}}(P_\gamma(u))\{(0,0)\}=2,
 \qquad P_\gamma(u)(0,0)-f(0,0)=-\varepsilon.
\]
All other atoms lie on the contact set, and the orthogonality integral is
therefore \(-2\varepsilon\).  Finally, a positive local one-cycle at the
origin has weights \(a,a,b,b>0\) on the two opposing horizontal and two
opposing vertical rays; intersection with \(\gamma\) has coefficient
\(2a+2b>0\).  The strict-convexity criterion of
\cite[Remark~3.9]{BMPS26} applies.  This is the source's Example~6.23 with
its sign and contact calculation made explicit.

There is a complementary smooth-fan warning in
\cite[Example~6.24, Claim~6.25, and Remark~6.26]{BMPS26}.  On the Bergman
fan of the uniform matroid \(U_{3,4}\), assign value zero to the opposite
rays generated by \(e_{12}\) and \(e_{34}=-e_{12}\), and value one to the
other primitive ray generators.  The resulting polyhedrally convex
reference \(\gamma\) admits a bounded piecewise-affine obstacle \(g\geq
\gamma\) for which
\[
 P_\gamma(g-\gamma)=\gamma,
 \qquad
 \MA_{\mathrm{poly}}(\gamma)\{0\}=4,
 \qquad
 \gamma(0)<g(0).
\]
The source proves the envelope identity cone by cone: equality first holds
on the line \(\mathbf R e_{12}\), then on the original rays by one-variable
convexity, and finally on each two-dimensional cone by joining its boundary
to the diagonal ray.  The mass is the direct iterated-divisor computation
on that fan.  This example is tropically smooth but \(\gamma\) is identically
zero on \(\mathbf R e_{12}\).  Hence it is neither strictly convex nor
positive at infinity.  It does not disprove the conjunction of a suitable
higher-dimensional smoothness condition with the strict growth hypotheses
in \Cref{prob:frontier-polyhedral-regularity}.

The same Bergman-fan mechanism nevertheless survives the growth condition
itself.  The point is to retain a finite flat segment of the exceptional
line, which forces contact at the origin, while making the reference grow
outside that segment.

\begin{proposition}\label{prop:frontier-polyhedral-bergman-positive-infinity-obstruction}
In the notation of \cite[Example~6.24 and Claim~6.25]{BMPS26}, let
\(g\geq\gamma\) be its obstacle with
\(\varepsilon\in\mathbf Q_{>0}\), and put \(u=g-\gamma\).
Use the real scalar category \(\mathrm{PA}_{\mathbf R}\) throughout.
Write
\(\ell(x_1,x_2,x_3)=x_1+x_2\), so that
\[
 \ell(e_{12})=2,
 \qquad
 \ell(e_{34})=-2,
 \qquad
 \bigl(\ell(e_1),\ell(e_2),\ell(e_3),\ell(e_4)\bigr)
   =(1,1,0,-2).
\]
For \(\lambda\in\mathbf Q\cap(0,\tfrac12)\) and
\(R\in\mathbf Q\) with \(R>2\lambda\varepsilon\), set
\[
 \gamma_{\lambda,R}
   :=\max\{\gamma,\lambda\ell-R,-\lambda\ell-R\}.
\]
Then \(\gamma_{\lambda,R}\) is piecewise-affine polyhedrally convex and
positive at infinity.  If
\[
 P:=P_{\gamma_{\lambda,R}}(u),
\]
then \(P=\gamma\) on a neighbourhood of the origin.  Moreover, either
\[
 P\notin
 \mathrm{PC}_{\mathrm{reg},\mathbf R}
     (\mathbf X,\gamma_{\lambda,R}),
\]
or
\[
 \int_{\mathbf X}
 \bigl(P-(\gamma_{\lambda,R}+u)\bigr)\,
 \MA_{\mathrm{poly},\mathbf R}(P)
 \leq -4\varepsilon<0.
\]
Consequently \((\mathbf X,\gamma_{\lambda,R})\) cannot have both
real-polyhedral regularity and real-polyhedral orthogonality.
\end{proposition}



On the primitive exceptional ray \(e_{12}\), the eventual slope of
\(\gamma_{\lambda,R}\) is \(2\lambda\in(0,1)\).  This explains the
real scalar convention in the proposition: the growth correction has a
nonintegral slope.  The literal noncompact obstruction is instead supplied
by \Cref{thm:frontier-polyhedral-noncompact-literal-nonregularity}.

Thus a higher-dimensional smoothness-and-orthogonality theorem requires
more than a polyhedrally convex reference positive at infinity: even on
this Bergman fan, that growth condition is compatible with a forced
noncontact Monge--Amp\`ere atom.  The proposition leaves open whether an
additional combined intrinsic-smoothness and strict-convexity package restores
orthogonality.

Balancing also does not, by itself, produce an analytic space.  Once a
proper calibrated tropical morphism has been supplied, however, the
recursive comparison follows from the calibration and the stated
generated-psh pullback condition.  Genuine algebraic tropicalizations give
the principal local analytic case under the value-group hypothesis in the
following corollary.







\providecommand{\an}{\mathrm{an}}

\begin{proposition}\label{prop:p053-projection-recursion-repaired-candidate}
Let \(d\geq1\), and let \((\mathbf X,[\mathbf X])\) be a pure
\(d\)-dimensional balanced rational polyhedral space, endowed with its
boundaryless tropical-space structure and its integration current
\(\delta_{[\mathbf X]}\).  Let \(Y\) be a pure \(d\)-dimensional
boundaryless CLD tropical space, let \([Y]\) be its positive closed current
of integration, and let

\[
 h:Y\longrightarrow\mathbf X
\]

be a topologically proper morphism of tropical spaces.  Assume that

\begin{equation}
 h_*[Y]=\delta_{[\mathbf X]}
 \label{eq:p053-repaired-top-calibration}
\end{equation}

and that

\begin{equation}
 h^*\operatorname{PAPC}(\mathbf X)
 \subseteq \operatorname{CPsh}(Y).
 \label{eq:p053-repaired-psh-pullback}
\end{equation}

For \(a_1,\ldots,a_d\in\operatorname{PAPC}(\mathbf X)\), set

\begin{equation}
 S_0=[Y],\qquad
 S_j=\dprimeddprime\bigl((a_j\circ h)S_{j-1}\bigr)
 \quad (1\leq j\leq d),
 \label{eq:p053-repaired-source-recursion}
\end{equation}

and

\begin{equation}
 c_0=[\mathbf X],\qquad c_j=a_j\cdot c_{j-1}.
 \label{eq:p053-repaired-target-recursion}
\end{equation}

Then every \(S_j\) is positive and closed, every \(c_j\) is an effective
balanced tropical cycle, and

\begin{equation}
 h_*S_j=\delta_{c_j}
 \quad (0\leq j\leq d).
 \label{eq:p053-repaired-intermediate-pushforward}
\end{equation}

In particular,

\begin{equation}
 h_*\!\left(
   \dprimeddprime[a_1\circ h]\wedge\cdots\wedge
   \dprimeddprime[a_d\circ h]\wedge[Y]
 \right)
 =\delta_{a_1\cdots a_d\cdot[\mathbf X]}
 =\MA_{\mathrm{poly}}(a_1,\ldots,a_d).
 \label{eq:p053-repaired-pa-comparison}
\end{equation}

Both iterated products in \eqref{eq:p053-repaired-pa-comparison} are
independent of the displayed order.

Let now \(\gamma\) be a piecewise-affine polyhedrally convex reference
function that is positive at infinity.  If

\[
 a_1,\ldots,a_d\in\mathrm{PC}_{\mathrm{reg}}(\mathbf X,\gamma),
\]

then every \(a_i\circ h\) is continuous psh and

\begin{equation}
 h_*\!\left(
   \dprimeddprime[a_1\circ h]\wedge\cdots\wedge
   \dprimeddprime[a_d\circ h]\wedge[Y]
 \right)
 =\MA_{\mathrm{poly}}(a_1,\ldots,a_d)
 \label{eq:p053-repaired-regularizable-comparison}
\end{equation}

as Radon measures on \(\mathbf X\).  For every simplicial polyhedral
structure \(\Pi\) carrying \(\gamma\), the measure on the right, viewed on
the associated BMPS compactification \(\mathbf X_\Pi\), has zero mass on
\(\mathbf X_\Pi\setminus\mathbf X\).
\end{proposition}

\begin{corollary}\label{cor:p053-algebraic-chart-repaired-candidate}
Let \(K\) be an algebraically closed complete field with a nontrivial
non-Archimedean absolute value, with
\(\mathbf Q\subseteq-\log|K^\times|\), and let \(U\) be an irreducible very affine
\(K\)-variety of dimension \(d\).  Use the convention

\[
 h=\operatorname{trop}_U=(-\log|f_1|,\ldots,-\log|f_n|):
 U^{\an}\longrightarrow\mathbf X:=\operatorname{Trop}(U)
\]

for a closed embedding of \(U\) in a split torus.  On the CLD side this is
the sign-reversed moment chart.  Equip \(\mathbf X\) with its saturated
lattices and analytic tropical multiplicities, and assume that it is a
rational balanced polyhedral space in the sense of
Botero--Mazzon--Pille-Schneider.  Then
\eqref{eq:p053-repaired-top-calibration} and
\eqref{eq:p053-repaired-psh-pullback} hold for \(Y=U^{\an}\).  Consequently,
all conclusions of
\Cref{prop:p053-projection-recursion-repaired-candidate} apply.

For \(a_i\in\operatorname{PAPC}(\mathbf X)\), put \(u_i=a_i\circ h\),
endow the trivial line bundle with the piecewise-smooth delta-metric

\[
 \lVert1\rVert_i=\mathrm e^{-u_i},
\]

and write \(\Theta_i=c_1(\mathcal O_U,\lVert\cdot\rVert_i)\) for its first
Chern delta-form.  Each \(\Theta_i\) is a positive delta-form.  More
precisely, for every ordering and every \(1\leq j\leq d\), the product of
the first \(j\) factors is a positive generalized delta-form; this assertion
is made for these particular Chern forms and not for arbitrary products of
positive generalized delta-forms.  The top-degree delta-form
\(\Theta_1\wedge\cdots\wedge\Theta_d\) therefore induces a positive,
locally finite Radon measure \(\mu_\delta\) on \(U^{\an}\), and

\begin{equation}
 [\mu_\delta]
 =\dprimeddprime[u_1]\wedge\cdots\wedge
   \dprimeddprime[u_d]\wedge[U^{\an}],
 \qquad
 h_*\mu_\delta=\MA_{\mathrm{poly}}(a_1,\ldots,a_d).
 \label{eq:p054-repaired-delta-comparison}
\end{equation}

No finiteness of the total mass of \(\mu_\delta\) is asserted without an
additional properness or support hypothesis.

More generally, take
\(a_i\in\mathrm{PC}_{\mathrm{reg}}(\mathbf X,\gamma)\) and arbitrary
decreasing approximations

\[
 a_{i,k}\in\operatorname{PAPC}(\mathbf X,\gamma),
 \qquad a_{i,k}\searrow a_i.
\]

Let \(\mu_{\delta,k}\) be the locally finite delta-form measure associated
with \(a_{1,k},\ldots,a_{d,k}\).  Then \(\mu_{\delta,k}\) converges vaguely
on \(U^{\an}\) to

\begin{equation}
 \mu_{\mathrm{CLD}}
 :=\dprimeddprime[u_1]\wedge\cdots\wedge
   \dprimeddprime[u_d]\wedge[U^{\an}].
 \label{eq:p054-repaired-source-limit}
\end{equation}

The limit is independent of all approximating sequences and satisfies

\begin{equation}
 h_*\mu_{\mathrm{CLD}}
 =\MA_{\mathrm{poly}}(a_1,\ldots,a_d).
 \label{eq:p054-repaired-limit-pushforward}
\end{equation}

No assertion is made that \(\mu_{\mathrm{CLD}}\) is represented by a finite
delta-form or has finite total mass without an additional properness or
support hypothesis.
\end{corollary}






The proposition and its algebraic-chart corollary remove the recursive
projection and local delta-form comparison gaps formerly isolated as
P05.3 and P05.4 under their explicit operator-domain hypotheses.  In the
algebraic corollary, \(\mathbf Q\subseteq-\log|K^\times|\) ensures that
the rational polyhedral weights pull back to real PL functions, so the
book's PL comparison gives membership in the generated psh sheaf.  For
a general valued field this membership remains an additional input,
even when the ordinary current is positive.  The propositions do not
supply that general-field operator-domain bridge, a proper polarized
realization, a global semipositive lift, or the mass calibration required
by P05.5.

The final boundary step in a polarized realization is nevertheless formal
once the local comparison and the exact total masses are known:




\begin{proposition}
\label{prop:p055-boundary-mass-completion}
Let \(j:X\hookrightarrow\overline X\) be an open embedding of locally compact
Hausdorff spaces.  Let \(\overline Y\) be a locally compact Hausdorff space,
let \(\overline h:\overline Y\to\overline X\) be a proper continuous map, and
put
\[
 U:=\overline h^{-1}(X),\qquad h:=\overline h|_U:U\longrightarrow X.
\]
Let \(\mu_{\mathrm{NA}}\) be a finite positive Radon measure on
\(\overline Y\), and let \(\mu_{\mathrm{poly}}\) be a finite positive Radon
measure on \(X\).  Suppose that
\begin{equation}
 h_*\bigl(\mu_{\mathrm{NA}}|_U\bigr)=\mu_{\mathrm{poly}}
 \label{eq:p055-local-pushforward}
\end{equation}
and
\begin{equation}
 \mu_{\mathrm{NA}}(\overline Y)
 =\mu_{\mathrm{poly}}(X)<\infty.
 \label{eq:p055-equal-mass}
\end{equation}
Then
\begin{equation}
 \mu_{\mathrm{NA}}(\overline Y\setminus U)=0
 \label{eq:p055-zero-source-boundary}
\end{equation}
and
\begin{equation}
 \overline h_*\mu_{\mathrm{NA}}=j_*\mu_{\mathrm{poly}}
 \label{eq:p055-global-pushforward}
\end{equation}
as finite positive Radon measures on \(\overline X\), where
\(j_*\mu_{\mathrm{poly}}\) denotes extension by zero across
\(\overline X\setminus X\).

Consequently, in a polarized realization for P05.5, assume that
\[
 \mu_{\mathrm{NA}}=\MA_\theta(u),\qquad
 \mu_{\mathrm{poly}}=\MA_{\mathrm{poly}}(a),
\]
that \eqref{eq:p055-local-pushforward} holds, and that established mass
formulae give
\[
 \mu_{\mathrm{NA}}(\overline Y)=(L^d)
 =\deg(\gamma)=\mu_{\mathrm{poly}}(X).
\]
Then the global comparison is \eqref{eq:p055-global-pushforward}.  This
consequence asserts none of the required realization, polarization,
plurisubharmonicity, local-comparison, or mass-formula hypotheses.
\end{proposition}




At the model level, the local comparison, mass identity, and calibration
inputs can be packaged into an explicit common-measure statement.  The
following proposition also records exactly where the GK comparison is not
being asserted under the standing discretely valued field hypotheses.




\begin{proposition}
\label{prop:p055-model-calibration-transfer}
Assume the BFJ field and variety hypotheses of
\textup{(Book 24.1)}, with the variety denoted by $Y$ and
$\dim Y=d$.  Let $\mathbf X$ be a pure $d$-dimensional balanced rational
polyhedral space, let $\gamma$ be a piecewise-affine polyhedrally convex
reference function positive at infinity, and let
$a_1,\ldots,a_d\in\mathrm{PAPC}(\mathbf X,\gamma)$.  Write
\begin{equation}
 \mu_{\mathrm{poly}}
 :=\MA_{\mathrm{poly}}(a_1,\ldots,a_d)
 =\sum_{q\in Q}m_q\delta_q,
 \qquad m_q>0,
 \qquad \sum_{q\in Q}m_q=\deg(\gamma),
 \label{eq:p055-model-polyhedral-atoms}
\end{equation}
where $Q\subseteq\mathbf X$ is finite.

Let $j:\mathbf X\hookrightarrow\overline{\mathbf X}$ be an open embedding
into a locally compact Hausdorff space.  Let $L$ be an ample line bundle on
$Y$, let $\|\cdot\|_0$ be a semipositive model metric on $L$, and put
$\theta=c_1(L,\|\cdot\|_0)$.  For $1\leq i\leq d$, let
\begin{equation}
 \|\cdot\|_i=\mathrm e^{-u_i}\|\cdot\|_0
 \label{eq:p055-model-relative-metrics}
\end{equation}
be a semipositive model metric on $L$, and set
$\theta_i=\theta+\ddc u_i$.  Suppose that all the $\theta_i$ are determined
on one normal flat projective model $\mathcal Y$, write
\[
 \mathcal Y_s=\sum_E b_EE,
 \qquad x_E=b_E^{-1}\ord_E,
\]
and define the nonnegative intersection weights
\begin{equation}
 w_E:=b_E
 \bigl(\theta_{1,\mathcal Y}|_E\cdots
       \theta_{d,\mathcal Y}|_E\bigr)\geq0.
 \label{eq:p055-model-intersection-weights}
\end{equation}

Let $\overline h:Y^{\mathrm{an}}\to\overline{\mathbf X}$ be a proper
continuous map, put
\[
 U:=\overline h^{-1}(\mathbf X),
 \qquad h:=\overline h|_U,
\]
and assume that
\begin{equation}
 u_i|_U=(a_i-\gamma)\circ h
 \qquad(1\leq i\leq d).
 \label{eq:p055-model-lift}
\end{equation}
Assume moreover the exact mass identity
\begin{equation}
 (L^d)=\deg(\gamma)
 \label{eq:p055-model-mass-match}
\end{equation}
and the componentwise calibration
\begin{equation}
 \sum_{\substack{E:\ x_E\in U\\ h(x_E)=q}}w_E=m_q
 \qquad(q\in Q).
 \label{eq:p055-model-component-calibration}
\end{equation}

Then the BFJ mixed model measure, the Chambert--Loir mixed measure, and the
CLD product of the formal first Chern currents all coincide with
\begin{equation}
 \mu_{\mathrm{NA}}
 :=\MA_\theta(u_1,\ldots,u_d)
 =\sum_Ew_E\delta_{x_E}.
 \label{eq:p055-model-common-measures}
\end{equation}
They satisfy
\begin{equation}
 \mu_{\mathrm{NA}}(Y^{\mathrm{an}}\setminus U)=0
 \label{eq:p055-model-boundary-zero}
\end{equation}
and
\begin{equation}
 \overline h_*\mu_{\mathrm{NA}}
 =j_*\mu_{\mathrm{poly}}
 \label{eq:p055-model-global-comparison}
\end{equation}
as finite positive Radon measures on $\overline{\mathbf X}$.

The conclusion is conditional on the proper polarized realization, the
global model lifts \eqref{eq:p055-model-lift}, semipositivity, the mass
identity \eqref{eq:p055-model-mass-match}, and the intersection-theoretic
calibration \eqref{eq:p055-model-component-calibration}.  It asserts none
of those inputs for an arbitrary balanced polyhedral space.  It also makes
no assertion about a GK delta-form product on the original analytification
under these discretely valued field hypotheses.
\end{proposition}




Thus no separate compactification argument is needed after those two hard
inputs have been established.  The two propositions do not construct a
proper polarized realization or the componentwise calibration and therefore
narrow, but do not retire, P05.5.

The unconditional analytic result in the source is narrower.  For the
toric degenerations of Calabi--Yau complete intersections in its
\S7, let \(\mathbf X\) be the tropicalization, let
\(\gamma\) be the polarization weight, and let \(\mu_0\) be Lebesgue
measure on the essential skeleton.  Under the three assumptions
\[
 (L^d)=\deg(\gamma),\qquad
 \varphi\in\mathrm{PC}_{\mathrm{reg}}(\mathbf X,\gamma),
 \quad\MA_{\mathrm{poly}}(\varphi)=\mu_0,
\]
and semipositivity of
\(
 \Psi=\varphi_0+(\varphi-\gamma)\circ\operatorname{trop}
\), Proposition~7.2 of \cite{BMPS26} proves
\(
 \MA_{\NA}(\Psi)=\mu_0
\)
and, over \(\mathbf C((t))\), the comparison property.  Its Lemma~7.3
computes every maximal-cell weight as one from analytic degree divided by
lattice index.  The local equality is then
\eqref{eq:frontier-polyhedral-real-local} together with Vilsmeier's
real/non-Archimedean comparison \cite[Theorem~1.1]{Vil20}; equality away
from lower-dimensional skeleton faces and equality of total masses force
the global identity.  None of this is a realization theorem for an abstract
balanced space.

The source itself formulates three sharper tasks that organize this broad
problem: enlarge the notion of polyhedral smoothness so that higher-dimensional
regularity and orthogonality hold; characterize polarized Calabi--Yau
families whose tropicalizations satisfy those properties; and determine
which matroids have Bergman fans satisfying orthogonality
\cite[pp.~3--5, \S1.5, and Remark~6.27]{BMPS26}.  By contrast, the analytic
realization of the cross-fan was a diagnostic problem formulated in this
book, not a question attributed to Botero--Mazzon--Pille--Schneider.  Its
literal compact-source and weighted-current parts are now settled by
Book Proposition~69.2.3, Book Proposition~69.2.4.
The source's only stronger strict-positivity prospect is expressly assigned
to its upcoming sequel: its bibliography lists that sequel as ``in
preparation.''  It is therefore not used here as an available theorem on
regularity or orthogonality \cite[pp.~5 and 60]{BMPS26}.
The stronger upstairs-equality problem in
\Cref{prob:frontier-cross-fan-upstairs} is not merely book-local: after
removing an ill-defined delta-current alternative, it is precisely the
analytic specialization of the Bedford--Taylor question posed in
\cite[\S8.6.1 and Definition~8.6.2]{CLD25}.

Thus the audit gives six verdicts.  The real-scalar complete affine case in
\Cref{thm:frontier-polyhedral-affine-orthogonality} is \emph{verified},
whereas literal integral-slope regularity already fails on the affine line
by \Cref{prop:frontier-polyhedral-integral-regularity-counterexample};
strict positivity alone is \emph{disproved} by
\Cref{ex:frontier-polyhedral-cross-orthogonality}; positive growth alone,
even on the displayed Bergman fan, cannot force the conjunction of
regularity and orthogonality by
\Cref{prop:frontier-polyhedral-bergman-positive-infinity-obstruction}; the
recursive CLD and
delta-form comparison is \emph{verified} under the intrinsic proper
calibration hypotheses of
\Cref{prop:p053-projection-recursion-repaired-candidate}, with the genuine
algebraic tropicalization supplied, under its stated value-group
hypothesis, by \Cref{cor:p053-algebraic-chart-repaired-candidate}; the global polarized
non-Archimedean Monge--Amp\`ere comparison remains \emph{open} beyond the
realized source cases because the proper polarization, lift, mass, and
component calibration are additional inputs; and the main
higher-dimensional local contact problem remains \emph{open}.  The missing local
assertion is a contact theorem on singular fans: under an intrinsic
smoothness and strict-positivity package, positive top intersection mass of
an envelope should force contact with the obstacle.  Neither a proof nor a
counterexample satisfying both parts of that package is presently known.

\clearpage\section{Recognizing continuous zero-Hessian functions}
\label{family:5}
\noindent\textbf{Active items:} P06.2. \textbf{Status: Open in the residual scope.}

\subsection{Atomic targets}
\subsubsection*{P06.2 (T2): Recognition of continuous zero-Hessian functions}
\label{atom:P06.2}
\addcontentsline{toc}{subsection}{P06.2: Recognition of continuous zero-Hessian functions}
Let \(K\) be complete and nontrivially valued and \(X\) good, separated, strict and boundaryless. Outside the verified scopes, decide whether every continuous \(h\) with \(\dprimeddprime[h]=0\) is classically pluriharmonic (or harmonic where that notion is available), or give a counterexample in this scope. Scalar ascent of the current equation and regularization of both signs are separate missing gates. The unrestricted arbitrary-radius product extension is false.

\subsection{What constitutes a solution}
Prove recognition for general continuous zero-Hessian functions in a new stated scope, or construct a counterexample satisfying the original strict ambient hypotheses. Make scalar-current ascent and two-sign regularization separate obligations. Ordinary smooth, weakly smooth and G-local real-PL conclusions have distinct consequences.

\subsection{Full statement and detailed state of the art}
\label{sec:ch72-pluriharmonic}

The supplied singular descent theorem
Book Theorem~18.12.1 settles P06.1: classical
pluriharmonicity implies local real PL regularity throughout the original
scope, including arbitrary singular points. P06.2 retains the reverse
bridge from a current equation to the curve-test notion.

\begin{problem}\label{prob:frontier-pluriharmonic}
Let $K$ be complete and nontrivially valued and let $X$ be a good,
separated, strict, boundaryless $K$-analytic space. If a continuous $h$
satisfies $\dprimeddprime[h]=0$, must $h$ be classically pluriharmonic,
or harmonic under the hypotheses where a classical harmonic notion is
available, outside the positive cases recorded below?
\end{problem}

The residual problem has two logically separate gates.  One must first
control the zero-current equation after scalar extension, and one must then
regularize continuous positive-current functions into the CLD-psh
hierarchy.  The next proposition records that these two gates suffice and
that every counterexample must obstruct at least one of them.

\medskip


\begin{proposition}
\label{prop:p062-scalar-regularization-two-gate}
Let $K$ be complete and nontrivially valued, let $X$ be a good,
separated, strict, boundaryless $K$-analytic space, and let
$h\colon X\to\mathbf R$ be continuous with
$\dprimeddprime[h]=0$.  Suppose that there is an analytic extension
$L/K$ such that, for the scalar-extension map
$\pi\colon X_L\to X$ and $h_L:=h\circ\pi$, the following two conditions
hold:
\begin{enumerate}
\item the current equation ascends, so that
$\dprimeddprime[h_L]=0$;
\item current positivity regularizes for the two signs of $h_L$: if
$f\in\{h_L,-h_L\}$ and $\dprimeddprime[f]\geq0$, then $f$ is CLD-psh.
\end{enumerate}
Then $h$ is classically pluriharmonic and locally $\mathbf R$-PL.  If $X$
is pure-dimensional, then $h$ is also GJR harmonic.
\end{proposition}

\begin{corollary}
\label{cor:p062-counterexample-two-gate-obstruction}
Let $K$ be complete and nontrivially valued, let $X$ be a good, separated,
strict, boundaryless $K$-analytic space, and let
$h\colon X\to\mathbf R$ be continuous with $\dprimeddprime[h]=0$.
Suppose that $h$ is not classically pluriharmonic.  For every analytic
extension $L/K$, at least one of the following failures occurs:
\begin{enumerate}
\item $\dprimeddprime[h_L]\neq0$;
\item one of $h_L$ and $-h_L$ has nonnegative Hessian current but does
not belong to the CLD-psh hierarchy.
\end{enumerate}
For $L=K$, the first alternative is impossible.  Hence every negative
answer to P06.2 furnishes, in one of the two signs of $h$, a continuous
zero-current function outside the CLD-psh hierarchy.  If $X$ is
pure-dimensional, this is a counterexample to the converse of
Book Proposition~15.9.8 in its stated scope.
\end{corollary}








This reduction neither pulls a continuous current through an arbitrary
scalar extension nor proves positive-current regularization.  It identifies
those operations as distinct proof obligations and shows that a negative
answer to P06.2 produces a zero-current function outside the CLD-psh
hierarchy at the ground field.  In the pure-dimensional scope this would
disprove the missing nonsmooth converse in
Book Proposition~15.9.8.

The following dated verification history explains those positive cases.
Its former P06.1 residuals are superseded by the singular descent theorem.


The real-coefficient convention in the resolved P06.1 theorem is sharp.  On the strict
semistable annulus, \(\sqrt2\log|T|\) is classically pluriharmonic and
\(\mathbf R\)-PL but is neither \(\mathbf Q\)-PL nor integral PL; see
Book Proposition~M.1.1.

The quasi-smooth scope is now settled.  By
Book Theorem~M.1.1, every classically pluriharmonic function
on a good strict quasi-smooth analytic space is \(\mathbf R\)-PL and GJR
harmonic.  The proof first passes, after a finite ground-field extension, to
Temkin's rigid \(\acute{e}\)tale cover by an affinoid with a strictly
semistable formal model.  It then descends \(\mathbf R\)-PL regularity through
the full rigid \(\acute{e}\)tale cover by the de Jong--van der Put sheaf
criterion; finite-flat descent is used only for the finite
\(\acute{e}\)tale affinoid equalizers where its hypotheses hold.  The theorem
requires no global quasicompactness, separatedness, or boundarylessness.

There is also a direct finite-flat descent closure.  If a good strict
neighbourhood admits a finite, flat, surjective analytic cover by a
neighbourhood in any verified P06.1 regularity range, then
Book Corollary~M.1.1 descends \(\mathbf R\)-PL
regularity and GJR harmonicity for every classically pluriharmonic function.
It is an application of the finite-flat \(\mathbf R\)-PL descent theorem
\cite[Proposition~5.11]{GJR26}, not a claim that arbitrary singular germs
admit a finite-flat cover of this kind.

The published curve comparison also settles every strict analytic curve,
including singular and nonreduced curves
\cite[\S4.1.2 and Proposition~3.15]{GR25Classical}.  At that checkpoint, the residual P06.1 atom inside the scope of
\Cref{prob:frontier-pluriharmonic} began at non-quasi-smooth points in
analytic dimension at least two.  The supplied singular theorem below
subsequently removed that residual.  The exact frozen
statement and source-aware proof passed production Rethlas using
\texttt{gpt-5.6-sol} at \texttt{xhigh}, final \texttt{correct}, with zero
critical errors and zero gaps.  The controlling identifiers are:
\begin{center}
\small
\begin{tabular}{@{}ll@{}}
controlling run & \nolinkurl{20260829T043457Z_0b1b4d161e3b}\\
controlling item & \nolinkurl{pi_7cd35ebd843cf80a7cd696bb}
\end{tabular}
\end{center}

There is now a further singular higher-dimensional scope.  By
Book Theorem~M.1.2, Book Corollary~M.1.2, every classically pluriharmonic
function is \(\mathbf R\)-PL and GJR harmonic on a good strict product
\[
             C_1\times_K\cdots\times_K C_m\times_K Y,
\]
where the \(C_i\) are geometrically reduced good strict analytic curves and
\(Y\) is good, strict, and quasi-smooth.  The same conclusion holds
analytically locally on open analytic subspaces of such products.  This
includes the nodal and cuspidal cylinders and their finite products in the
characteristic-zero scope of Book Example~M.1.1.
After base change to a completed algebraic closure, the proof passes to the
product of the curve normalizations with \(Y\), applies
Book Theorem~M.1.1, and descends one finite normalization at
a time by the rational-section tube argument.  It is not a PL-descent
theorem for arbitrary finite maps.

The exact statement and source-aware proof passed the archived production
run \nolinkurl{20260829T191019Z_0b5e24145108} using
\texttt{gpt-5.6-sol} at \texttt{xhigh}, final \texttt{correct}, with zero
critical errors and zero gaps; the controlling item is
\nolinkurl{pi_6add90bcf0cfded12d6a074b}.  This result is retained as the
normalization-descent input for the stronger statement below.

The curve-factor hypotheses can now be made optimal in dimension one.  By
Book Theorem~M.1.3, the same conclusion holds
when every \(C_i\) is an arbitrary good strict analytic curve, with no
reducedness or geometric-reducedness assumption, and when the remaining
good strict factor \(Y\) is only required to have quasi-smooth reduction.
The case \(m=0\) includes nilpotent thickenings of quasi-smooth spaces.  The
proof combines invariance of \(\mathbf R\)-PL functions under nilpotent
closed immersions, an explicit nilpotent-kernel calculation for completed
tensor products, and the preceding successive curve-normalization theorem.
It also applies on the stated local open-product models.

The exact statement and proof have SHA-256 values
\begin{center}
\small
\begin{tabular}{@{}c@{}}
\nolinkurl{6954b3ebc1f9efb060b4858ce732d9720d62741ba2397c39573312dd6e3ba8f3}\\
\nolinkurl{e308ba5df336e38717ac27413c6f669a7fa37535bed613a79f095811c6dba4a7}
\end{tabular}
\end{center}
Two independent exact-byte audits passed.  Dependency-closed production
Rethlas verification with \texttt{gpt-5.6-sol} at \texttt{xhigh} returned
the controlling item
\begin{center}
\small
\nolinkurl{pi_7dceb6442b4d586c4475feca}.
\end{center}
Its disposition is \texttt{verified} and the final verdict is
\texttt{correct}, with zero critical errors and zero gaps.

There is also a nonproduct toric scope.  By
Book Theorem~M.1.4, a classically pluriharmonic germ
at the torus-fixed point of the affine toric variety associated to any
full-dimensional strongly convex rational cone is locally constant.  The
proof constructs cofinal weighted strict affinoid neighbourhoods with
singleton Shilov boundary and applies the classical maximum--minimum
principle.  Neither simpliciality nor a residue-characteristic hypothesis is
used.  The exact statement and proof have SHA-256 values
\begin{center}
\small
\begin{tabular}{@{}c@{}}
\nolinkurl{edbe69950a0a9c81ca2fd9077a58be96ec6bdb6c27ce9cf2524da5bf653a5ab6}\\
\nolinkurl{531e326b4aeb19067fcb88af2f6eb6477a169587d71a3249103777c9131cf8f4}
\end{tabular}
\end{center}
Two independent exact-byte audits passed.  Production Rethlas verification
with the same model and scrutiny returned the controlling item
\begin{center}
\small
\nolinkurl{pi_dfd932c583fee3acd3fefbe3}.
\end{center}
Its disposition is \texttt{verified} and the final verdict is
\texttt{correct}, with zero critical errors and zero gaps.

There is a broader non-toric vertex mechanism.  By
Book Theorem~M.1.5, if an affine
algebraic cone is defined by equations homogeneous for positive integral
weights, every classically pluriharmonic germ at its vertex is locally
constant.  For a point \(x\) in a sufficiently small weighted polydisc, the
map \(T_i\mapsto T_i(x)\tau^{w_i}\) from the closed unit disc over
\(\mathscr H(x)\) contracts \(x\) to the vertex.  Base-change and analytic
pullback preserve pluriharmonicity, and the disc maximum--minimum principle
then gives constancy.  This applies without normality, reducedness, or a
toric hypothesis; for example, when \(\operatorname{char}K\ne 3\), it treats
the non-toric normal isolated cubic cone \(x^3+y^3+z^3=0\) at its vertex.
It is only a weighted-cone-germ
criterion and supplies no such presentation for an arbitrary singular germ.

There is a relative non-toric extension along the zero section.  If a strict
affinoid family is cut out over a base by positive-weighted homogeneous fibre
equations, including equations with analytic base-dependent coefficients,
then Book Proposition~M.1.3 contracts each
point of a sufficiently small weighted tube to its point on the zero section
over its completed residue field.  Hence every classically pluriharmonic
function factors through the base on that tube.  A quasi-smooth base therefore
gives local \(\mathbf R\)-PL and GJR-harmonic regularity.  This needs neither
a product nor flat, reduced, normal, or irreducible fibres, but it does not
control arbitrary points away from the zero section or an arbitrary singular
base.

More generally,
Book Corollary~M.1.3 propagates the local
pluriharmonic \(\mathbf R\)-PL property from any verified base range through
each finite tower of such relative cones, along the composite zero section.
Intermediate bases may be singular, and coefficients at a later stage may
depend analytically on the earlier fibre variables.  The explicit two-stage
example of Book Example~M.1.2,
\[
 B_1:=K\langle x,y,z\rangle/(x^3+y^3+z^3),\qquad
 B_1\langle U,V,W\rangle/(UV-(x+y^2)W^2)
\]
has no common positive weighting for the displayed equations, even though
each stage has the required positive grading.  This demonstrates the extra
iterated scope without claiming that the resulting germ has no other
one-stage presentation.  The result still says nothing about points away
from the composite zero section or bases outside the initial verified
regularity range.

This range is stable under a genuine analytic standard-etale perturbation.
In the setting of
Book Theorem~M.1.6, the source equations
need not remain weighted homogeneous.  On shrinking the cone coordinates,
the coefficient errors of the standard-etale polynomial tend to zero in
the quotient Banach norms.  The contractive Newton isomorphism therefore
identifies the source tube with the same relative cone after base change to
its central standard-etale fibre.  If that fibre has the local
pluriharmonic \(\mathbf R\)-PL property---in particular, if the original
base is quasi-smooth, since the Jacobian criterion then makes the central
fibre quasi-smooth---the relative cone theorem gives the full local
\(\mathbf R\)-PL and GJR-harmonic conclusion at the source point.  The
explicit chart of
Book Example~M.1.3 contains the
nonhomogeneous equation \(T^2-T-x=0\).  This is still a central-fibre
criterion for the displayed strict standard-etale presentation; it neither
uniformizes an arbitrary singular germ nor treats an arbitrary point of the
source.

For an algebraic ordinary-etale morphism, this displayed analytic chart is
automatic at a source point over the zero section of a relative algebraic
positive-weighted cone with smooth base.  Indeed,
Book Theorem~M.1.7 applies the
algebraic etale structure theorem only to obtain an open immersion into a
standard-etale scheme, then separately retains a principal source open,
shrinks the target to a strict relative-cone affinoid, and imposes strict
bounds on the standard-etale coordinate and its inverse.  These bounds put
the chosen point in an actual ordinary analytic neighbourhood before the
preceding tube theorem is invoked.  This removes the displayed-chart
hypothesis for algebraic etale germs in that zero-section scope, but it says
nothing about points away from the zero section, genuinely log-etale maps,
or singular germs without a relative weighted-cone target.

The algebraic conclusion also propagates through arbitrary finite towers of
the displayed relative weighted cones.  The proof of
Book Theorem~M.1.8 does not incorrectly
reapply the one-step theorem over a singular intermediate base.  Instead it
strengthens the induction statement to every separated finite-type algebraic
etale source over a shorter tower.  The top central fibre is an ordinary
neighbourhood inside precisely such an algebraic etale source, so the
central-germ hypothesis of the analytic tube theorem follows by induction.
No common grading among stages, or flatness, reducedness, normality, or
irreducibility, is required.  The conclusion remains confined to the
composite zero section and does not create the required tower near an
arbitrary singular germ.

There is a further finite-quotient range.  By
Book Proposition~M.1.4, if a finite constant
group acts analytically on a separated strict affinoid space and a continuous
function on the analytic quotient pulls back to an \(\mathbf R\)-PL function,
then the function itself is \(\mathbf R\)-PL.  The proof combines
ordinary-open integral-PL decomposition, a common equivariant formal model,
tensor induction, and the possibly wild formal coarse quotient.  It assumes
neither flatness nor freeness, reducedness, normality, or invertibility of the
group order.  The exact statement and proof have SHA-256 values
\begin{center}
\small
\begin{tabular}{@{}c@{}}
\nolinkurl{e1daa4b3e3efd64fcc8f0ce05b9c1d5d6cb69b28686b939b48535fc56b0ccc07}\\
\nolinkurl{a6c24d1686c726dc72aebe55a1e83eb8d4fe1ed4025d4a8d7166da59123ce30e}
\end{tabular}
\end{center}
Two source-aware hostile audits passed.  Production Rethlas run
\nolinkurl{20260829T235341Z_e1daa4b3e3ef}, using
\texttt{gpt-5.6-sol} at \texttt{xhigh}, returned final
\texttt{correct}, with zero critical errors and zero gaps; the direct
endpoint assigned no item identifier.  Since classical pluriharmonicity is
stable under pullback, this settles P06.1 on finite constant-group quotients
of every separated strict affinoid source belonging to one of the verified
source scopes above.

The split simplicial toric range is now settled domain-wide, not merely at
the closed orbit.  Let \(\sigma\subset N_{\mathbf R}\) be any strongly
convex simplicial rational cone, possibly non-full-dimensional.  By
Book Theorem~M.1.9, every classically pluriharmonic
function on every good strict analytic domain in
\(X_{\sigma,N}^{\mathrm{an}}\) is \(\mathbf R\)-PL and GJR harmonic, over
an arbitrary complete nontrivially valued ground field.  The proof replaces
the lattice by the finite-index lattice generated by the primitive ray
vectors and a complement.  The resulting cover is smooth but need not be
flat.  After a finite scalar extension its prime-to-characteristic part is
a finite constant-group quotient, while its equal-characteristic primary
part is radicial and satisfies a uniform power condition.  The latter
regularity descends by Book Proposition~M.1.5: rational
domains and analytic units descend after taking one common Frobenius power.
Thus the argument includes wild diagonalizable quotient singularities and
does not invoke finite-flat descent.

The simpliciality and closed-orbit restrictions can now be removed.  By
Book Theorem~M.1.10, every classically pluriharmonic
function on every ordinary open analytic subspace of an arbitrary split
toric variety is \(\mathbf R\)-PL and GJR harmonic.  This includes every
positive-dimensional orbit of every non-simplicial cone.  The proof writes a
principal neighbourhood of an arbitrary orbit as a pointed transverse toric
factor times a split torus.  The parameterized theorem
Book Proposition~M.1.6 tests every completed-
residue-field fibre; the transverse semigroup Tate algebra has singleton
Shilov boundary, so the maximum--minimum principle forces constancy along
that fibre.  Fan gluing gives arbitrary split toric varieties, and
Book Corollary~M.1.7 descends the result for every
form that becomes split after the completion of an algebraic extension.

The orbitwise conclusion is stable under a good strict quasi-smooth
parameter.  More precisely,
Book Corollary~M.1.5 applies to every ordinary
open subspace of \(\mathcal X^{\mathrm{an}}\times_KY\), with
\(\mathcal X\) split toric and \(Y\) good, strict, and quasi-smooth, and
to the corresponding analytic local product models.  At an arbitrary toric
orbit, the principal-open decomposition supplies a pointed transverse toric
factor; the new parameter is the product of the orbit torus with \(Y\),
which remains quasi-smooth.  The fixed-fibre factorization and the
quasi-smooth theorem then apply.  This does not assert a result for a
nonsplit toric factor, a non-quasi-smooth parameter, or a germ without an
actual local product chart.
If the toric factor is a form that becomes split after the completion of an
algebraic extension, the same product conclusion follows from the
scalar-extension and ground-field descent in
Book Corollary~M.1.7.  Thus this further range still
excludes forms without such an algebraic splitting extension and does not
relax the quasi-smooth-parameter or local-product hypotheses.

The toric and curve-factor mechanisms may now occur simultaneously.
Book Corollary~M.1.9 permits an arbitrary
split toric variety, an arbitrary finite product of good strict analytic
curves, and a final good strict factor with quasi-smooth reduction.  At each
toric orbit, the principal-open decomposition exposes the pointed transverse
toric factor and absorbs the orbit torus into the final factor.  The
remaining parameter is then exactly in the scope of
Book Theorem~M.1.3.  Scalar extension and
\(\mathbf R\)-PL descent give the same conclusion for toric forms split by the
completion of an algebraic extension.  This does not cover a reduced
higher-dimensional singular factor outside the toric factor, a form without
such a splitting extension, or a nonproduct singular germ.

The same descent is analytic local.  If each point has an ordinary
neighbourhood which becomes an ordinary open subspace of a split toric
variety after a complete extension that is the completion of an algebraic
extension, then
Book Corollary~M.1.8 gives the same
conclusion, also after product with a good strict quasi-smooth parameter.
It retains the displayed local splitting chart and does not produce one for
an arbitrary analytic germ.

This analytic-local range admits the arbitrary curve factors as well.
Under the pointwise completed-algebraic splitting-chart hypothesis,
Book Corollary~M.1.10 applies the
split product theorem after scalar extension and descends the resulting
\(\mathbf R\)-PL function.  The displayed splitting chart remains an input:
the result does not construct it, allow a merely nonalgebraic splitting
extension, or identify a locally twisted nonproduct germ with a product.

The frozen statement/proof SHA-256 values are
\begin{center}
\small
\begin{tabular}{@{}c@{}}
\nolinkurl{05bda70eb1cf7ebbcc95d1831feadd946afa168f64a7298171381b9f1e865e60}\\
\nolinkurl{21c44c47a91a28469fb91a6c00569ed56d12699bae9c956c68c33ad493825b03}.
\end{tabular}
\end{center}
An independent exact-byte hostile audit passed.  The dependency-closed
production request using \texttt{gpt-5.6-sol} at \texttt{xhigh} returned
\texttt{correct}, with zero critical errors and zero gaps.  Its archived
request/response SHA-256 values are
\begin{center}
\small
\begin{tabular}{@{}c@{}}
\nolinkurl{def28bf9f16539fadab9b8b99dbd435d00eed8c3730c496d2b8832ee865a29c9}\\
\nolinkurl{718044118ee65e0d76b7c02b8b553552b8da974367513aa5ed4d57e8b6231f7d}.
\end{tabular}
\end{center}

The local product can also be constructed from an ordinary-etale source
chart rather than assumed.  In
Book Proposition~M.1.8, a separated finite-type etale map to
the product of a pointed split toric variety and a smooth variety becomes,
over a sufficiently deep weighted toric tube, an actual affinoid product
\(
 U\times_KD
\), where \(D\) is the retained quasi-smooth fixed-fibre factor.  The proof
uses one common Newton radius for four exact and perturbed root problems and
uses the two crossed uniqueness statements to identify the inverse maps.
It treats toric rank zero separately, so no empty maximum is hidden at the
dense orbit.  Singleton-Shilov factorization then gives
Book Theorem~M.1.11; the source-Zariski-local
consequence is Book Corollary~M.1.12.

An independent exact-byte hostile re-audit passed, and the dependency-closed
production request using \texttt{gpt-5.6-sol} at \texttt{xhigh} returned
\texttt{correct}, with zero critical errors and zero gaps.  The principal
and production receipt prefixes are
\begin{center}
  \small
  \begin{tabular}{@{}cc@{}}
    statement & \texttt{fbb294c435dffacf} \\
    proof & \texttt{10d188621a07242e} \\
    request & \texttt{dab3e30657a51183} \\
    response & \texttt{7e06480126ac47ad}
  \end{tabular}
\end{center}

The same ordinary-etale chart conclusion descends when the toric factor is an
algebraic-split form.  More precisely,
Book Corollary~M.1.16 pulls a
source-Zariski-local ordinary-etale chart to a complete extension which is
the completion of an algebraic extension, where the preceding split-chart
corollary applies, and then invokes \cite[Corollary~5.12]{GJR26} to descend
\(\mathbf R\)-PL regularity.  This is a scalar-extension consequence of the
verified split chart theorem; it does not assert a product construction
downstairs or supply an ordinary-etale chart for a general toroidal germ.

The manufactured ordinary-etale product also carries arbitrary finite curve
factors.  In
Book Corollary~M.1.18, the retained
quasi-smooth fixed fibre is grouped with the final factor; its reduction is
quasi-smooth, so the arbitrary-curve-factor theorem applies after
singleton-Shilov factorization.  The same completed-algebraic-extension
descent treats algebraic-split toric factors.  The statement still requires
the displayed ordinary-etale source chart and an external product, and makes
no claim for a genuinely log-etale or locally twisted nonproduct germ.

There is also a directly analytic standard-etale extension.  If the source
already has a strict affinoid chart over a weighted pointed split toric
domain times a good strict quasi-smooth affinoid parameter, with a monic
standard-etale Tate presentation, then
Book Proposition~M.1.9 constructs the same
actual product tube.  The singleton-Shilov factorization gives
Book Theorem~M.1.12, and the local-chart
form is Book Corollary~M.1.13.  This is a
direct reuse of the coefficient-contraction and common-Newton-radius proof;
it does not claim that an arbitrary analytic etale morphism admits the
stated standard-etale affinoid chart.

The all-orbits refinement Book Corollary~M.1.14
uses the principal toric decomposition at the image point and absorbs its
split-torus factor into the quasi-smooth parameter.  It therefore covers an
actual standard-etale strict affinoid chart over an arbitrary split toric
base whenever its target affinoid is an ordinary neighbourhood of that
image.  It still neither supplies such a chart for an arbitrary analytic
etale map nor treats a non-quasi-smooth parameter.

If that toric base becomes split over a complete analytic extension which is
the completion of an algebraic extension,
Book Corollary~M.1.17 applies the
all-orbits chart result after scalar extension and descends its
\(\mathbf R\)-PL conclusion.  It retains the ordinary-target chart
condition and does not cover a base without such an algebraic splitting
extension.

Arbitrary finite curve factors also pass through this strictly analytic
standard-etale tube.  In
Book Corollary~M.1.19,
the analytic coefficient contraction first manufactures the weighted toric
product.  Singleton-Shilov factorization then leaves the curve factors, the
quasi-smooth fixed fibre, and a final factor with quasi-smooth reduction, so
the arbitrary-curve theorem applies.  Completed-algebraic-extension descent
includes algebraic-split toric bases.  The conclusion still requires the
displayed monic standard-etale presentation, ordinary-target condition, and
external product.

Before the supplied singular descent theorem, these results narrowed
P06.1 to the following residual scopes:
arbitrary nonproduct non-quasi-smooth analytic germs in dimension at least
two outside the displayed product, quotient, and positive-weighted-cone
vertex ranges, products containing a reduced singular factor of analytic
dimension at least two outside those ranges, genuinely log-\(\acute{e}\)tale charts not covered by an
ordinary-etale or analytic standard-etale chart of the displayed kind, and
nonsplit local models outside the verified algebraic-split descent range.

Weakly smooth \(\dprimeddprime\)-harmonic functions are harmonic PL by
\cite[Proposition~10.17]{GJR26}.  The real-PL current problem is now settled:
Book Proposition~M.1.10 proves, on a good separated
boundaryless strict pure-dimensional space, that
\(\dprimeddprime[h]=0\) for an \(\mathbf R\)-PL function forces every residual
class \(L_h(x)\) to be numerically trivial, hence makes \(h\) harmonic and
classically pluriharmonic.  Its proof keeps the codimension-one calibrated
cells on a common analytic skeletal refinement and separates the
integral PL summands by \(\mathbf Q\)-linear independence; more precisely, the
coefficients of those calibrations in primitive lattice volumes are rational.
It does not pull back an arbitrary current to a curve.

The repaired exact-live statement and source-aware skeletal proof passed
production Rethlas using \texttt{gpt-5.6-sol} at \texttt{xhigh}, with final
verdict \texttt{correct}, zero critical errors, and zero gaps.  For readable
line breaking, the controlling and historical identifiers are displayed
separately:
\begin{center}
\small
\begin{tabular}{@{}ll@{}}
controlling run & \nolinkurl{20260825T061351Z_655e7f503e5d}\tabularnewline
controlling item & \nolinkurl{pi_eea98ea0a54108cae4e076be}\tabularnewline
frozen precursor & \nolinkurl{20260825T032958Z_3a27f19be958}\tabularnewline
rejected live repair & \nolinkurl{20260825T055921Z_655e7f503e5d}\tabularnewline
rejected separator proof & \nolinkurl{20260825T031243Z_1603e0a83037}
\end{tabular}
\end{center}
The frozen precursor is preserved but was superseded after the live proof was
repaired to use one common local chart, rational calibrated coefficients, and
a source-side cutoff.  The rejected live repair chose the affinoid before the
local common chart and relied on an undeclared internal comparison theorem.
The earlier target-side separator proof misused a source that separates a
chosen full-dimensional Abhyankar point, not the support of an arbitrary
complementary test form.  None of these three historical proofs is used.

Book Proposition~M.1.11 supplies a further
higher-dimensional continuous scope.  Let \(X\) be pure of dimension
\(n\geq 1\).  If \(h=\varphi\circ F\) factors through one pleasant full-rank
moment whose image is the nonempty connected open set
\(F(X)=F(\Sigma_F^{(n)})\), if \(F\) is proper in the compact-inverse-image
sense, and if the pushed-forward canonical calibration is the constant
positive multiplicity \(m|e_1\wedge\cdots\wedge e_n|\), then
\(\dprimeddprime[h]=0\) forces \(\varphi\) to be affine.  Hence \(h\) is ordinary
CLD-smooth, \(\mathbf R\)-PL, GJR harmonic, and classically pluriharmonic.  In
the boundaryless setting of \Cref{prob:frontier-pluriharmonic},
Book Corollary~M.1.20 derives the
constant calibration from harmony and gives the equivalence between current
vanishing and affinity of \(\varphi\).

The proof uses the exact proper-pushforward identity
\(F_*[h]=m[\varphi]\) and Euclidean distributional-Hessian recognition; it does
not pull an arbitrary CLD current back to curves.  The exact frozen statement
and source-aware proof passed production Rethlas run
\nolinkurl{20260829T031126Z_75e9cb498ce4}, item
\nolinkurl{pi_f96fe9f8904943602779f274}, using
\texttt{gpt-5.6-sol} at \texttt{xhigh}, final \texttt{correct}, with zero
critical errors and zero gaps.

There is a complementary invariant split-toric boundary result.  Let
\(V=\operatorname{trop}^{-1}(O)\) be a good strict open inverse image for
extended tropicalization in a split toric variety over a complete
nontrivially valued field.  If a finite continuous
maximal-compact-torus-invariant function \(h\) on \(V\) has
\(\dprimeddprime[h]=0\), then
Book Proposition~M.1.12
makes \(h\) locally CLD-smooth with zero Hessian, real PL, GJR harmonic, and
classically pluriharmonic, including at toric boundary strata.  The proof
uses the continuous invariant-function correspondence and current
pushforward calculation of \cite[(5.5) and Proposition~5.3.10(iv)]{BGJKK25}.
On a basic compactified tropical neighbourhood, finite continuity along the
normal rays forces every boundary-normal affine slope to vanish; the residual
slopes are logarithms of analytic units.  Thus
\cite[Proposition~4.2 and Theorem~4.3]{GR25Classical}, together with
\cite[Definition~7.2]{GJR26}, supplies the two asserted notions of
pluriharmonicity.

The exact statement and proof SHA-256 values are
\begin{center}
\small
\begin{tabular}{@{}c@{}}
\nolinkurl{8d69697806c33d59531b5d3247a2f9e0c1d07e313d329c15a2432646134a544a}\\
\nolinkurl{c17e99ebb54f45003d78f55670a978092294b44a7c4536b78d7207ec52cfda85}.
\end{tabular}
\end{center}
The dependency-closed source context and exact packet passed the independent
model-backed verifier as item
\nolinkurl{pi_30ec53e8836ba1d2b08fdf64}, final \texttt{correct}, with zero
critical errors and zero gaps.  This result does not cover arbitrary
continuous functions, arbitrary toroidal charts, or arbitrary analytic
germs, so it narrows but does not retire P06.2.

There is an algebraic-split scalar-stable extension of this invariant result.
If a good strict analytic form becomes a split toric open inverse image after
a complete analytic extension which is the completion of an algebraic
extension and the pullback of the finite continuous function is
maximal-compact-torus invariant with zero Hessian there, then
Book Corollary~M.1.22 descends
local \(\mathbf R\)-PL regularity, classical pluriharmonicity, and GJR
harmonicity.  This adds no automatic current-ascent claim: the invariant
zero-current condition is imposed upstairs as part of the hypothesis.
The same conclusion is available analytically locally by
Book Corollary~M.1.23,
provided the germ becomes a displayed split toric domain after a complete
algebraic extension and the upstairs hypotheses are imposed there.  This is
not a construction of such an algebraic-split model for an arbitrary analytic
germ.
The scalar-stable argument also applies to a locally trivial standard strict
open-polydisc bundle over the displayed algebraic-split form when the lifted
maximal compact torus action, invariance, and zero-current equation are all
imposed after that extension:
Book Corollary~M.1.68
first invokes the split-toric bundle result upstairs and then descends the
three regularity conclusions.  It asserts neither automatic current ascent
nor such a bundle presentation for a general germ.
The analytic-local companion
Book Corollary~M.1.70
needs these data only on a neighbourhood of each point of the total space.
It still imposes all lifted hypotheses there and does not produce a global or
local bundle chart from an arbitrary analytic germ.

There is now also a factorization-free local theorem on strict polydiscs.
By Book Theorem~M.1.14, over an algebraically
closed complete nontrivially valued field, every continuous function on a
strict open polydisc whose CLD \(\dprimeddprime\)-current vanishes is
constant.  The same holds on a strict closed polydisc when the current
vanishes on its analytic interior.  The proof restricts the current to every
translated punctured polydisc, uses top-form support and calibrated
integration on the split-torus Gauss section to obtain zero distributional
Hessian for every radial trace, kills all affine slopes by shrinking one
radius at a time, and then separates rigid points by translation-invariant
Gauss norms.  It neither pulls the current back to a curve nor assumes that
all analytic points lie on one moment skeleton.

The exact statement and proof have SHA-256 values
\begin{center}
\small
\begin{tabular}{@{}c@{}}
\nolinkurl{017f3da108ae737e4bbc1d490d6988faf76ce643b4cb52104aae53c4dcda75d6}\\
\nolinkurl{3763bf9a7e4b1f47602544ad2c3248d9b8abd7b9d34109e3c14259e422393bbd}
\end{tabular}
\end{center}
An independent exact-byte hostile audit passed.  The production provenance
is
\begin{center}
\small
\begin{tabular}{@{}ll@{}}
run & \nolinkurl{20260830T001728Z_017f3da108ae}\\
model and scrutiny & \texttt{gpt-5.6-sol}, \texttt{xhigh}.
\end{tabular}
\end{center}
The final verdict was \texttt{correct}, with zero critical errors and zero
gaps; the direct endpoint assigned no item identifier.

Over an arbitrary complete nontrivially valued field, the same local
split-torus calculation gives the exact residual statement
Book Proposition~M.1.14: a continuous zero-current function
has
\[
                         h(\eta_{a,\mathbf r})=h(a)
\]
at every \(K\)-rational centre and is therefore constant on all
\(K\)-points, all \(K\)-centred Gauss skeleta, and the closure of their
union.  Book Proposition~M.1.15 then gives full constancy if
the zero-current equation is explicitly known after one complete
algebraically closed extension.  On the open polydisc it also gives full
constancy without an extra scalar-extension hypothesis when both \(h\) and
\(-h\) are CLD-psh in the approximation sense, since psh functoriality
produces positive opposite Hessian currents upstairs.

The exact statement and proof have SHA-256 values
\begin{center}
\small
\begin{tabular}{@{}c@{}}
\nolinkurl{7510083fecda2d561c4ce2a1b35d08377268ef48adda2b242172428dbd39caec}\\
\nolinkurl{77a013d4e655d288d8803c9650df5446bed570de2644591edddfc116c323eab2}.
\end{tabular}
\end{center}
Two independent exact-byte hostile audits passed.  The production provenance
is
\begin{center}
\small
\begin{tabular}{@{}ll@{}}
run & \nolinkurl{20260830T015556Z_7510083fecda}\\
model and scrutiny & \texttt{gpt-5.6-sol}, \texttt{xhigh}.
\end{tabular}
\end{center}
The final verdict was \texttt{correct}, with zero critical errors, zero gaps,
and eight reference checks without an issue; the direct endpoint assigned no
item identifier.

The non-\(K\)-rational obstruction is now removed completely in analytic
dimension one on a strict disc.  By
Book Proposition~M.1.16, a continuous multiplier has a
continuous calibrated trace under every finite flat map of constant degree,
the associated current is the direct image current, and zero Hessian
descends.  Given any rigid point of a strict disc, its minimal polynomial
defines such a disc map.  Its zero fibre is the chosen rigid point with full
degree, while a graded Gauss-valuation calculation gives a second
single-point fibre of the same degree, including in ramified and inseparable
cases.  The rational-centred theorem applied to the trace and to the original
function therefore makes all rigid values equal; density proves
Book Proposition~M.1.17.

The exact statement and proof have SHA-256 values
\begin{center}
\small
\begin{tabular}{@{}c@{}}
\nolinkurl{67a8d8a21e23cb34cf9b8f6685efbf2c9eaa6ff484b3696cbc1867a8940e095c}\\
\nolinkurl{479643ee3aa534bbfc98adcf69837356d3cd8eff22d92accd1319aabb049553e}.
\end{tabular}
\end{center}
An independent exact-byte hostile audit passed.  The dependency-closed
production request archived as
\nolinkurl{20260830T033245Z_5c4711a474d8}, using
\texttt{gpt-5.6-sol} at \texttt{xhigh}, returned final
\texttt{correct}, zero critical errors, zero gaps, and thirteen checked
citation uses; the endpoint assigned no run or item identifier.

The arbitrary-field strict-open classification is the first product step.
By
Book Proposition~M.1.18, a continuous
zero-current function on a product with an arbitrary boundaryless pure
one-dimensional strict fibre has zero Hessian on that fibre over every base
point, over its actual completed residue field.  The proof spreads an
arbitrary compactly supported ordinary smooth fibre test by ambient moment
germs and a partition of unity, then uses fibre integration and the
projection formula.  Combining the arbitrary-field disc theorem with
Book Proposition~M.1.19 and iterated Gauss
replacement proves
Book Theorem~M.1.16: on a finite product
of strict open discs and annuli,
\[
                    h=c+\sum_{j=1}^{q}\lambda_j\log|S_j|.
\]
In particular, every continuous zero-current function on a strict open
polydisc is constant over an arbitrary complete nontrivially valued field.

The earlier frozen v3 statement/proof SHA-256 values were
\begin{center}
\small
\begin{tabular}{@{}c@{}}
\nolinkurl{985be85d97b583f72be5cee65988df6f31c1edb49a685121b834df179fa87fda}\\
\nolinkurl{42b592bf8ac26ed46c36928e5e90272225ae6883bd2bfe62d3ff2840f01d46e5}.
\end{tabular}
\end{center}
At that verification gate, an independent exact-byte audit and the
dependency-closed production request using \texttt{gpt-5.6-sol} at
\texttt{xhigh} returned \texttt{correct}, with zero critical errors and zero
gaps.  Its archived request/response SHA-256 values are
\begin{center}
\small
\begin{tabular}{@{}c@{}}
\nolinkurl{6261b61bc49669a74a6eee4a2d834c25c8741bef59b084b81eaa684629f2d604}\\
\nolinkurl{c8c951075a574c6a3badde2e2a22942718616cd21da0144cb8a983bc3cac822f}.
\end{tabular}
\end{center}
This v3 receipt is retained only as historical provenance.  A later
transitive source audit found that its slicing proof used a local-moment
result in the wrong direction.  The source-correct v4 proof and production
receipt below supersede it; no statement in the book depends on the v3
proof.

The one-dimensional endpoint variants are nevertheless now settled by a
different argument.  By
Book Proposition~M.1.20, if one or both endpoints
of a standard strict annulus are included and a continuous function has
zero CLD current on the absolute analytic interior, then
\[
                              h=a\log|T|+b
\]
on the whole domain.  The proof places the polynomial tube of each rigid
endpoint point inside the analytic interior, takes its finite-flat trace to
a strict open disc, and uses a singleton full-degree Gauss fibre plus
compact-complement convergence.  It never extends \(h\) beyond its stated
domain and makes no product or higher-dimensional endpoint claim.

The exact frozen statement/proof SHA-256 values are
\begin{center}
\small
\begin{tabular}{@{}c@{}}
\nolinkurl{a54723f040926d5a309494bb77d3c734c1c9d27bfb9fd18c56fc70ea4d03c9e1}\\
\nolinkurl{461afad5804a3269cb385550af6f4abcb5543d4f3151db56f986490c74e6947e}.
\end{tabular}
\end{center}
After an independent exact-byte release audit, the dependency-closed
production request using \texttt{gpt-5.6-sol} at \texttt{xhigh} returned
\texttt{correct}, with zero critical errors and zero gaps.  Its archived
request/response/header SHA-256 values are
\begin{center}
\small
\begin{tabular}{@{}c@{}}
\nolinkurl{34a4d9f0d4b10a1f5a4693adef7a233a8ee756813dfe1f21450dc0cfd2d67242}\\
\nolinkurl{6b5aa2445bbd8b1dbe594fc23571c979d64eb5dc839f038101ce97b93181cee4}\\
\nolinkurl{7d1ce47df4f2cceae087f95d72337b90ce2a3a2b6c51cabf015395597ff3e396}.
\end{tabular}
\end{center}
The direct endpoint assigned no run or item identifier.  No novelty or
priority claim is made.

The higher-dimensional standard strict endpoint variants are now settled as
well.  Let \(P\) be any finite product of standard strict open or closed
discs and standard strict open, closed, or half-closed annuli.  By
Book Theorem~M.1.15, if a continuous function
has zero CLD Hessian current on the absolute analytic interior of \(P\), then
\[
                         h=c+\sum_j\lambda_j\log|S_j|
                         \qquad\text{on }P.
\]
The proof only uses
\(
 \prod_i\operatorname{Int}(F_i)\subseteq\operatorname{Int}(P)
\), not equality.  It slices there, applies the one-dimensional
classification after finite residue extension at rigid base points, and
propagates the strict-core formula one endpoint coordinate at a time.  It
never restricts a current directly to a closed factor fibre and never
asserts invariance of analytic interiors under an arbitrary complete scalar
extension.

The principal v4 statement and proof have SHA-256 hashes
\begin{center}
\small
\begin{tabular}{@{}c@{}}
\nolinkurl{56ffaad6c42e0a14c56279e263d326120ce03f6ae132d8261deee56ff9ca7663}\\
\nolinkurl{2616ac69cd450010b8f0d108d030bd69d2ff2728cad9641ad7d9d0e1a981a0b7}.
\end{tabular}
\end{center}
An independent exact-byte hostile audit passed.  The dependency-closed
production request using \texttt{gpt-5.6-sol} at \texttt{xhigh} returned
\texttt{correct}, with zero critical errors and zero gaps.  Its archived
request/response/header SHA-256 values are
\begin{center}
\small
\begin{tabular}{@{}c@{}}
\nolinkurl{7b5bdbb4a15fb8d5202d57ddb4c49f84e1844908863c95580ceec7c1f1324171}\\
\nolinkurl{c0d9aaf5a8ef3400b6b20782a0a8d9073c24c200027c4451a9aec9e5432b446d}\\
\nolinkurl{1013a47eaabe8bdeec363937720e76b988fb76dbd90f2009ca10f6a513d0cb94}.
\end{tabular}
\end{center}

The arbitrary-radius continuation now has a sharp negative--positive answer.
Let \(R\notin\sqrt{|K^\times|}\) and consider the non-strict bidisc
\[
 P=\{|T|\leq R\}\times_K\{|S|\leq R^{-1}\}.
\]
By Book Proposition~M.1.21, the tube
\[
 W=\{|TS-1|<1\}\cap P
\]
is nonempty, relatively open in \(P\), and disjoint from its absolute
analytic interior.  Hence, for \(0<c<1\), the nonconstant continuous cutoff
\[
 h_c=\max\{0,1-c^{-1}|TS-1|\}
\]
has zero Hessian current on \(\operatorname{Int}(P)\).  Thus the
unrestricted arbitrary-radius extension is false already for a product of
two closed discs.  On the other hand,
Book Proposition~M.1.22 proves the full formula
\[
                    h=c_0+\sum_j\lambda_j\log|S_j|
\]
for every arbitrary-radius standard endpoint polycylinder whenever, after
one complete extension strictifying all radii, the pulled-back current
vanishes on the full upstairs absolute interior.  The proof applies the
strict-v4 theorem upstairs and descends only the resulting pointwise formula
through the surjective scalar-extension map; it does not assert automatic
ascent of the downstairs current.

The exact frozen statement/proof SHA-256 values are
\begin{center}
\small
\begin{tabular}{@{}c@{}}
\nolinkurl{27ff3ad7a6174b8b38d5890d8c891681edf416f02e2538017db22edf372c1ed8}\\
\nolinkurl{15ff9b3a28b63c1e3f37c104f27e30e77431f389ec694a1d3f3823232ef1f125}.
\end{tabular}
\end{center}
After an independent exact-byte hostile re-audit, the dependency-closed
production request using \texttt{gpt-5.6-sol} at \texttt{xhigh} returned
\texttt{correct}, with zero critical errors and zero gaps.  Its archived
request/response/header SHA-256 values are
\begin{center}
\small
\begin{tabular}{@{}c@{}}
\nolinkurl{bb90f025d3e6fb616918a364a9f1b014968a58164ff142557ec8188a45b40759}\\
\nolinkurl{839c5b3f0506f9789f2f14d5eaff8afc0c2bd20b17ba11bf7400b3a5e0a7582a}\\
\nolinkurl{3f31d65b71c13800f8486cf689d9fb508fa35c7b44e0b94084292abffd0dad66}.
\end{tabular}
\end{center}

These results do not close P06.2.  Outside the standard strict coordinate
product and proper full-rank factor scopes, current vanishing for an
arbitrary continuous function supplies neither translated torus charts
covering the germ, the required factor map, nor a legitimate pullback to all
curve probes.  Arbitrary analytic germs and nonproduct local domains remain
open.  For non-strict standard products, the unrestricted downstairs claim
is now disproved and the scalar-stable replacement is proved; what remains is
to characterize useful conditions weaker than scalar stability that still
control the invisible boundary.  The theorem also makes no inference from
an equation imposed only on the coordinatewise strict radial core.
Higher-dimensional standard strict product endpoints are no longer part of
the residual problem.  In the
published curve scope---an open subset of the analytification of a smooth
proper algebraic curve over an algebraically closed complete nontrivially
valued field---Wanner's comparison also settles it; see
Book Corollary~19.3.1.

There is also an exact recovery endpoint which does not require uniform
convergence once the approximants are already classically pluriharmonic.  If
a continuous candidate \(h\) is the pointwise limit of a sequence of
classically pluriharmonic functions, then
Book Proposition~15.14.2 makes \(h\) classically
pluriharmonic and upgrades the convergence locally to uniform convergence.
Thus a future proof of P06.2 may stop after producing such a pointwise
curve-test approximation.  The zero-current equation alone supplies no
such sequence, by Book Remark~15.14.1; nor are these
approximants ordinary CLD-smooth harmonic functions, so this recovery
criterion cannot replace the distinct approximation hypothesis used for
P21.3.

There is now a further non-coordinate product extension.  By
Book Proposition~M.1.23, a continuous zero-current function on
\(B\times_KD_R^-\), with \(B\) good, topologically separated, pure, and
strict, is the pullback of a unique continuous zero-current function on
\(B\).  The proof first uses the verified fibre-slicing theorem and the
arbitrary-field disc theorem to obtain fibrewise constancy, then descends
the current equation with a normalized compactly supported top-degree disc
form; it does not pull a continuous current back along the zero section.
Consequently, in the published curve scope every zero-current function on
\[
                         W\times_K\prod_{i=1}^rD_{R_i}^-
\]
is pulled back from \(W\) and is classically pluriharmonic by
Book Corollary~M.1.26; as a finite product of
strict analytic curves, it is also locally \(\mathbf R\)-PL and
Gubler--Jell--Rabinoff harmonic.  This narrows the residual without treating
annular fibrations, fibrations without local product charts, arbitrary
nonproduct local domains, or arbitrary analytic germs.

Finite globally coordinated annular products are now covered as well.  For
standard strict open annuli \(A_i=A^-_{r_i,R_i}\), a continuous zero-current
function on \(W\times_K\prod_{i=1}^sA_i\) has the unique form
\[
 p^*b+\sum_{i=1}^s(a_i\circ p)q_i^*\log|T_i|,
\]
where \(b\) has zero current on \(W\) and every \(a_i\) is locally constant,
by Book Theorem~M.1.18.  Exact
completed-residue-field slicing and two radial top-form averages first
produce a zero-current slope coefficient on the preceding base; a mixed
compact test gives it zero first current.  Induction uses the corresponding
first-current rigidity of the lower polyannulus product.  Thus the product
is classically pluriharmonic, locally \(\mathbf R\)-PL, and GJR harmonic.
The \(s=1\) instance is
Book Theorem~M.1.17.  Every annular coordinate
is global in this formula.  The regularity conclusion nevertheless
globalizes through analytic product charts:
Book Corollary~M.1.24 treats every locally
trivial finite standard strict open-polyannulus bundle over \(W\).  On each
displayed chart it applies the product theorem, so the function is
classically pluriharmonic, locally \(\mathbf R\)-PL, and GJR harmonic without
requiring the logarithmic coordinates or their coefficient functions to
agree globally.

For the globally coordinated product, the same formula persists through
every finite tower of locally trivial standard strict open-polydisc bundles.
Indeed, Book Corollary~M.1.50 first descends
a zero-current function to the global product, then applies the displayed
coefficient formula there.  Separately,
Book Corollary~M.1.51 gives
classical pluriharmonicity, local \(\mathbf R\)-PL regularity, and GJR
harmonicity for such a tower over any of the locally trivial polyannulus
bundles just described, but does not fabricate a global coefficient formula.
Thus only annular fibrations lacking the stated
analytic product charts, and arbitrary nonproduct germs, remain outside
these conclusions.

There is now a reverse-order tower range as well.  Start with a
boundaryless finite tower of locally trivial standard strict open-polydisc
bundles over \(W\), and only then form a globally coordinated finite product
with standard strict open annuli.  By
Book Corollary~M.1.49, every
continuous zero-current function on that product has the corresponding
unique logarithmic coefficient formula, now with its base term and its
locally constant slope coefficients on the top \(B_\ell\) of the disc
tower, and is classically pluriharmonic.  This uses first-current descent
through the disc tower and the first-current rigidity transferred from the
curve.  Its local curve-product presentation additionally makes it locally
\(\mathbf R\)-PL and GJR harmonic.  It still does not jointly trivialize
the disc tower over \(W\).

There is also a mixed chart range.  With globally coordinated standard
strict open-annulus and open-disc factors,
Book Corollary~M.1.27 first descends every
disc factor and gives the unique logarithmic coefficient formula over \(W\);
the result is classically pluriharmonic, locally \(\mathbf R\)-PL, and GJR
harmonic.  More generally, if every displayed local product has at least one
standard strict open-annulus factor and at least one standard strict
open-disc factor, then
Book Corollary~M.1.25 gives the same
three regularity conclusions on the total space.  It uses no global splitting
of either kind of coordinate and supplies no global logarithmic coefficient
formula in the locally trivial case.  Its finite locally trivial
open-polydisc-bundle towers are covered by
Book Corollary~M.1.52, with
the same three local regularity conclusions and still no global coefficient
formula.

Global product triviality is not needed in this curve range.  If a good,
separated strict analytic space is a locally trivial standard strict
open-polydisc bundle over such a \(W\), then
Book Corollary~M.1.28 gives the same classical
pluriharmonic, locally \(\mathbf R\)-PL, and GJR harmonic conclusions.  More precisely,
Book Proposition~M.1.29 first glues the local descents
to a global zero-current curve function, without a global section or product
decomposition.  This disc-bundle assertion by itself does not handle annular
fibres; the mixed chart result above does so only when the stated annular
and disc product coordinates are both present.  Neither statement covers a
disc fibration lacking the displayed local product charts.
The same argument is iterable: for every finite tower of such locally
trivial standard strict open-polydisc bundles over a published curve,
Book Corollary~M.1.48 descends a
continuous zero-current function all the way to the curve and hence proves
classical pluriharmonicity, local \(\mathbf R\)-PL regularity, and GJR
harmonicity on the total space.  Every level of the tower must carry its
displayed local product charts; this supplies no conclusion for an annular
fibre or for a disc fibration lacking those charts.

There is now also a proper-fibre bundle consequence.  On a locally trivial
projective-space bundle over a published curve,
Book Proposition~M.1.30 exhausts each relative standard
affine chart by strict open polydiscs, descends on every such chart, and uses
the common point \([1:\cdots:1]\) to identify the local descended functions.
The required locally holomorphically separated and paracompact base
hypotheses make the final current gluing legitimate.  Thus
Book Corollary~M.1.29 proves classical
pluriharmonicity, local \(\mathbf R\)-PL regularity, and GJR harmonicity for
every continuous zero-current function in this projective-bundle range.  It
is a proper-fibre globalization of the existing strict-disc descent, not a
solution for an arbitrary nonproduct germ or an arbitrary fibration.

The same finite-chart argument gives a strictly larger proper-fibre range.
If the fibre is the analytification of a positive-dimensional smooth proper
split toric variety, then
Book Proposition~M.1.31 uses the affine
spaces attached to the maximal cones of its complete regular fan.  Strict
open-polydisc exhaustions of those affine charts all contain the identity of
the dense split torus, so exact disc descent glues without an invariance
hypothesis.  Thus
Book Corollary~M.1.30 proves
classical pluriharmonicity, local \(\mathbf R\)-PL regularity, and GJR
harmonicity for every continuous zero-current function on such a locally
trivial bundle over a published curve.  This includes projective-space
bundles but also smooth proper split toric fibres which are not projective
spaces.  It does not treat singular or nonsplit toric fibres, nonproper
toric fibres, arbitrary fibrations, or arbitrary nonproduct germs.

There is another proper-fibre range with a geometrically connected smooth
proper algebraic curve as fibre.  In
Book Proposition~M.1.32, exact slicing is used
only over the dense set of \(K\)-rational rigid base points.  The fibre
there is over \(K\) itself, so the curve comparison and compact maximum
principle make the slice constant.  Density of rigid points in the entire
local product then gives the factorization through the base, and a normalized
compactly supported fibre top form descends the zero-current equation.  Thus
Book Corollary~M.1.31 proves classical
pluriharmonicity, local \(\mathbf R\)-PL regularity, and GJR harmonicity for
locally trivial bundles of this kind over a published curve.  The conclusion
does not extend here to an arbitrary proper curve fibre, a nonproper curve
fibre, or an arbitrary nonproduct germ.

The algebraic-closure restriction on the descent step now has an explicit
scalar-stable replacement.  Over an arbitrary complete nontrivially valued
field, Book Proposition~M.1.33 assumes
that, at every rigid base point, the fibre restriction still has zero
Hessian current after passage to a completed algebraic closure of the
completed residue field.  Wanner comparison and the compact maximum
principle then make that scalar-extended restriction constant.  A closed
point of the proper curve supplies a finite faithfully flat multisection;
equalizer descent and rigid-point density replace a rational section, with
no separability claim.  A normalized top form on the fixed fibre descends
the Hessian and first-current identities.  Thus every such function is the
pullback of a unique continuous zero-current function on the base.  This is
an explicit scalar-stability theorem, not an assertion that a downstairs
continuous current equality automatically ascends.

The exact statement/proof SHA-256 values are
\begin{center}
\small
\begin{tabular}{@{}c@{}}
\nolinkurl{cce5a5ff689f110e02564ee155b75027c3eec06d0e484af8e86dba0662fe9564}\\
\nolinkurl{b371852f0f87b4a1b382d6f67d792d8b34ba5299ee43e40323b01f7df8ede2d5}.
\end{tabular}
\end{center}
The repaired standalone packet passed the local model-backed verifier using
\texttt{gpt-6-astra} at \texttt{xhigh}, final \texttt{correct}, with zero
critical errors and zero gaps.  The rejected precursor and its incorrect
separability claim remain archived in the solver dossier, together with the
exact checked-item identifier.

The four exact bundle descents may now be interleaved.  In
Book Theorem~M.1.19, each stage of a finite
tower over a published curve may be a locally trivial standard strict
open-polydisc bundle, a projective-space bundle, a smooth proper split-toric
bundle, or a proper smooth connected curve bundle, in any order.  Descending
from the top stage to the curve
produces a unique continuous zero-current function there; curve comparison
and harmonic pullback then make the original function classically
pluriharmonic, locally \(\mathbf R\)-PL, and GJR harmonic.  Every
projective-space, split-toric, and proper-curve stage retains the locally
holomorphically separated and paracompact base assumptions of its source
proposition.  This
does not treat annular fibres, arbitrary fibrations, or nonproduct germs
without the displayed local trivializations.

One may nevertheless append globally coordinated annular factors at the
top.  If that top space is boundaryless, then
Book Corollary~M.1.32 gives a unique
coefficient formula for every continuous zero-current function on its
product with finitely many standard strict open annuli.  The zero-current
base coefficient and all locally constant annular slope coefficients
descend to the original curve; the resulting function is classically
pluriharmonic, locally \(\mathbf R\)-PL, and GJR harmonic.  This does not
supply locally twisted annular coordinates at the top or insert annular
stages at arbitrary intermediate levels of the tower.

Annular geometry may also occur at the bottom.  The general transfer
statement Book Proposition~M.1.34 isolates the
common exact descent, classical-pluriharmonic transfer,
\(\mathbf R\)-harmonic transfer, and first-current-rigidity transfer of the
four allowed bundle types over an arbitrary admissible base.  Applying it
to a globally coordinated curve-polyannulus product gives
Book Corollary~M.1.53: after any such mixed
tower, the top function still has the unique logarithmic coefficient formula
over the original curve.  Applying it instead to a locally trivial
curve-polyannulus bundle or a locally trivial mixed curve-polycylinder bundle
gives
Book Corollary~M.1.54, Book Corollary~M.1.55,
with classical pluriharmonicity, local \(\mathbf R\)-PL regularity, and GJR
harmonicity but no global coefficient formula.  Thus locally twisted
annular coordinates are now allowed at the bottom of the mixed exact tower.
The result still requires their displayed analytic local trivializations and
does not allow an annular stage to be interleaved at an arbitrary later
level, an annular fibration without such charts, or an arbitrary nonproduct
germ.

For globally coordinated annular products, the two endpoint orders can be
fully interleaved.  In
Book Theorem~M.1.20, each stage of a
finite boundaryless tower over the curve may be either one of the four exact
bundle types or a globally coordinated finite standard strict
open-polyannulus product.  An induction preserves classical
pluriharmonicity, \(\mathbf R\)-harmonicity, and first-current rigidity at
every level.  It gives a unique normal form on the tower top with one
logarithmic term for every annular coordinate introduced anywhere in the
tower; the zero-current base coefficient and all locally constant slope
coefficients live on the original curve.  Thus globally coordinated annular
stages may now occur at arbitrary levels.  For a global normal form, the
remaining obstruction is the transition monodromy of locally twisted annulus
bundles: the theorem does not choose global annular coordinates.  For one
locally trivial polyannulus bundle over the curve, however,
Book Proposition~M.1.26 now identifies the
coefficient system on overlaps.  The logarithmic coordinates transform by a
locally constant unimodular integer matrix and a vector of zero-current base
functions.  These data satisfy the affine cocycle law, define a rank-\(s\)
local system of lattices, and make the local slope vectors of every
zero-current function a global parallel section of its dual real local
system.  Thus the monodromy defect is encoded intrinsically even when it is
not trivialized.  The remaining translation defect is measured exactly by
the zero-current radiance class of
Book Definition~M.1.1.
Book Proposition~M.1.27 makes its vanishing
equivalent to a purely linear overlap law after zero-current affine gauge
and then gives one global coefficient formula for every zero-current
function.  The gauges are not asserted to be logarithms of analytic units,
so this criterion does not manufacture global annular coordinates.
Even when the full class does not vanish,
Book Proposition~M.1.28 gives the exact residual
answer slope by slope: the slopes of global zero-current functions are
precisely the parallel dual sections whose scalar contraction with the
radiance class vanishes.  Its exact sequence also shows that an injective
contraction, and in particular the absence of parallel dual sections,
forces unique descent of every zero-current function to the original curve.
When the radiance class vanishes, a chosen affine gauge splits this sequence.

For the three regularity conclusions, part of that monodromy obstruction can
be removed.  In
Book Corollary~M.1.33, one analytic open
cover of the original curve is allowed to coordinate every annular stage
after base change.  On each member of the cover, the whole restricted tower
is an instance of the interleaved theorem; the local conclusions therefore
glue even though the annular coordinates, radii, and product trivializations
need not agree on overlaps.  Thus curve-controlled locally twisted annular
stages may occur at arbitrary levels.  The result deliberately gives no
global coefficient formula.  It also does not treat annular trivializations
available only over arbitrary opens of intermediate tower levels, because
such opens need not retain the preceding tower or its paracompactness
hypotheses.

There is also a nonproduct endpoint-base range.  If the base is a positive-
dimensional standard strict endpoint polycylinder and the fibre is a locally
trivial standard strict open polydisc, then
Book Proposition~M.1.37 pulls every full
zero-current function back from the explicit affine combination of the base
annular logarithms.  Thus it is classically pluriharmonic, locally
\(\mathbf R\)-PL, and GJR harmonic.  This allows open, closed, and
half-closed annuli in the \emph{base}, but still does not treat an annular
fibre, a fibration without the displayed local product charts, or an equation
only on a strict radial core.
The same formula persists through every finite tower of the displayed locally
trivial strict open-polydisc bundles:
Book Corollary~M.1.56 first descends the
function to the endpoint-polycylinder base and then applies the endpoint
classification.  Each level must carry the stated local product charts; the
nontrivial tower result retains the formula together with local
\(\mathbf R\)-PL and GJR harmonicity, by propagating its local curve-product
charts.  This still does not cover annular fibres or fibrations without
them.

There is a further invariant nonproduct range.  Let the base be a good,
topologically separated, strict, pure positive-dimensional open inverse
image in the extended tropicalization of a split toric variety, and let a
locally trivial standard strict open-polydisc bundle carry a lift of its
maximal compact torus action.  If the continuous zero-current function is
invariant for that lifted action, then
Book Corollary~M.1.63 first descends it to
an invariant zero-current function on the toric base and then applies the
split-toric boundary theorem.  It is therefore locally real-PL,
Gubler--Jell--Rabinoff harmonic, and classically pluriharmonic.  The action,
invariance, and actual local open-polydisc trivializations are hypotheses:
this supplies no conclusion for noninvariant functions, nonsplit forms,
annular fibres, or arbitrary disc fibrations.
The same conclusion persists through every finite tower of such bundles
when the maximal compact torus acts on every level and all bundle maps are
equivariant: successive exact descents preserve invariance and reduce to
the same split-toric base theorem; see
Book Corollary~M.1.64.
The same bundle scope descends through a chosen complete extension that is
the completion of an algebraic extension when the base becomes the displayed
split toric domain and all three of the lifted action, invariance, and
zero-current conditions hold upstairs; see
Book Corollary~M.1.68.
The corresponding finite-tower statement is
Book Corollary~M.1.69;
it requires the full compatible tower, action, invariance, and current
equation after that chosen extension.
It does not manufacture the extension or bundle chart, and supplies no
automatic ascent of a downstairs current equality.
Its analytic-local form
Book Corollary~M.1.70
requires that displayed structure only near each point of the total space;
it makes no assertion when such neighbourhood data are absent.
The same localisation is valid for finite towers:
Book Corollary~M.1.71
requires a complete displayed tower near each point of the total space, with
the scalar-stable extension, the compatible action on every level, and
equivariance of every map.  It supplies the same three conclusions locally,
but does not create that tower or transfer a downstairs current equation
upstairs.

The former third atom in this environment is now resolved.
Book Theorem~19.3.4 proves that every
continuous current-psh function on every smooth proper algebraic curve over
an algebraically closed complete nontrivially valued field is, locally on a
strict affinoid neighbourhood, a uniformly convergent decreasing sequence of
ordinary CLD-smooth psh functions.  Equivalently, current positivity,
Thuillier subharmonicity, and the first ordinary-smooth decreasing-limit
stage coincide in this scope.  The proof passes through finite local tree
skeleta, rational strictification of graph Laplacians, Riemann--Roch movement
of residual divisors into fresh disc directions, and smooth convex caps.
The exact frozen statement and proof passed production Rethlas run
\nolinkurl{20260827T090645Z_767cbf81c723}, using
\texttt{gpt-5.6-sol} at \texttt{xhigh}, with final verdict
\texttt{correct}, zero critical errors, and zero gaps.  The direct endpoint
did not assign an item identifier.

The formerly written harmonic-approximant strengthening is false:
Book Proposition~19.3.2 constructs an integral
PL weakly smooth harmonic function on a smooth proper curve which is not even
a locally uniform limit of ordinary CLD-smooth harmonic functions.  Its
fixed-support argument uses zero curvature of every approximant and says
nothing against the now-proved psh regularization theorem, whose approximants
carry positive curvature in fresh directions.
The exact-live proposition and proof passed production Rethlas run
\nolinkurl{20260825T055907Z_b5c8f0719b60}, item
\nolinkurl{pi_d80993bb01a748f0422ae6a5}, with the same model and scrutiny,
final \texttt{correct}, zero critical errors, and zero gaps.  Thus weakly
smooth and CLD-smooth harmonic regularity on curves are genuinely different.
The supplied proof in Book Section~18.12 now covers all singular
points under the original P06.1 hypotheses, beyond the preceding product,
quotient and toroidal cases. It uses closed-subspace PL germs, finite and
proper descent, regular alteration covers and completed-algebraic scalar
descent. Its provenance is the author-supplied manuscript, separately
recorded from the historical production receipts above.
The real-PL zero-current subproblem, the proper
full-rank moment-factor subproblem, the algebraically closed strict-polydisc
subproblem, the arbitrary-field standard strict endpoint-polycylinder
subproblem,
the arbitrary-radius scalar-stable endpoint-polycylinder subproblem,
the algebraic-split scalar-stable invariant-toric subproblem,
its analytic-local algebraic-split analytic-form version,
the published smooth-proper-curve zero-current subproblem,
the arbitrary-field scalar-stable proper-smooth-curve-bundle descent
subproblem,
the smooth-proper-split-toric locally trivial bundle subproblem, and the
smooth-proper-curve ordinary psh
regularization subproblem are settled in their displayed scopes.  The
arbitrary continuous zero-current question remains open outside those
product scopes, on nonproduct local domains and arbitrary analytic germs;
for non-strict standard products, the corresponding downstairs statement is
false and only the scalar-stable repair is asserted.  A current equation only
on the coordinatewise strict radial core is also insufficient for the
endpoint theorem.
Accordingly P06.1 is retired and P06.2 remains active. Together with the
P07.1 counterexample below and the later supplied P08.1 counterexample, the
current inventory has \(42\) active atoms in \(12\) formal problem environments.

\clearpage\section{A geometric local theory on singular germs}
\label{family:6}
\noindent\textbf{Active items:} P12.2, P12.3, P12.4, P12.5, P12.6. \textbf{Status: Open in the residual scope.}

\subsection{Atomic targets}
\subsubsection*{P12.2 (T4): An independently geometric singular valuative psh cone}
\label{atom:P12.2}
\addcontentsline{toc}{subsection}{P12.2: An independently geometric singular valuative psh cone}
On the intrinsic normalized boundary of an excellent singular or nonnormal germ, construct a valuative psh class with independent geometric meaning and the correct smooth-germ specialization. Identify its relation to the existing ideal-generated cone. Defining another order-theoretic completion does not by itself supply that identification.

\subsubsection*{P12.3 (T4): An independently geometric singular envelope theory}
\label{atom:P12.3}
\addcontentsline{toc}{subsection}{P12.3: An independently geometric singular envelope theory}
Construct obstacle envelopes for the independently geometric singular psh cone, establish their positivity and the appropriate regularity, and compare them with the complete ideal-psh envelope already known on reduced-curve links.

\subsubsection*{P12.4 (T5): Model-independent bounded singular-local products}
\label{atom:P12.4}
\addcontentsline{toc}{subsection}{P12.4: Model-independent bounded singular-local products}
Construct bounded mixed Monge--Ampère products on singular and nonnormal normalized local boundaries, with model independence, the correct intersection comparison, locality and continuity. The existing compact boundary construction alone supplies no such product.

\subsubsection*{P12.5 (T5): Finite energy on singular local boundaries}
\label{atom:P12.5}
\addcontentsline{toc}{subsection}{P12.5: Finite energy on singular local boundaries}
Construct an energy functional, finite-energy class, and its Monge--Ampère calculus for excellent singular and nonnormal germs. Specify reference data, changes of primary gauge, and the relation between bounded approximation and finite-energy limits.

\subsubsection*{P12.6 (T4): Geometric finite-map functoriality and conductor descent}
\label{atom:P12.6}
\addcontentsline{toc}{subsection}{P12.6: Geometric finite-map functoriality and conductor descent}
Establish compatibility of independently geometric singular psh cones, envelopes and Monge--Ampère products with finite maps, normalization and conductor descent. Finite restriction, ideal-cone extrema, normal-domain norm traces, finite-flat determinant traces and their tower/gauge rules are known. They do not identify the geometric operators; an unrestricted continuous degree trace is obstructed.

\subsection{What constitutes a solution}
The requested geometric cone, envelopes, bounded products, finite energy and map compatibility are separate targets. Prove their smooth-germ specialization, primary-gauge changes and normalization/conductor rules. Do not identify an ideal-generated order completion with a geometric pluripotential theory without an additional comparison theorem.

\subsection{Full statement and detailed state of the art}
\label{sec:ch72-singular-local}

Finite-map progress is now available in Book Section~70.2:
continuous positive trace over normal targets without flatness, fiber
extrema for arbitrary finite maps onto their exact images, and the exact
conductor criterion for continuous descent. The ideal-generated cones
are explicit; their identification with a geometric singular psh class
and compatibility with envelopes and Monge--Amp\`ere products are not
claimed. This settles the valuative transfer layer of P12.6 and narrows
its remaining geometric target.

\begin{problem}\label{prob:frontier-singular-local}
Extend the normalized local theory of
Book Chapter~70 from smooth complex germs to excellent
singular and nonnormal germs. On the normalized boundary of
Book Section~70.1.1, construct, when possible,
psh functions, envelopes, bounded products, finite energy, and finite-map
functoriality.
\end{problem}

The normalizing-boundary part is resolved in arbitrary dimension.
Book Theorem~70.1.2 gives the intrinsic centered ray
space and its nonempty compact sections for all \(\mathfrak m\)-primary
ideals, with exact change \textup{(Book 70.7a)}.
Book Theorem~70.1.3 presents each section as a
quotient of the separated normalization branches. Fibers are singletons
where the conductor has finite value; on its infinite-value stratum,
equality of restrictions to the conductor quotient is the precise
identification rule. Nonprimary ideals define the partial sections
\textup{(Book 70.7e)}, which can be noncompact.
Nilpotents are invisible to these semivaluations, and a zero-dimensional
local reduction has empty centered cone. This settles P12.1, while
P12.2--P12.6 retain their separate psh, envelope, product, energy and
finite-map targets. It does not assert an intersection theory or
resolution in positive characteristic.

In dimension one, this construction reduces to a finite branch set. By
Book Theorem~70.1.4, the link normalized by an
\(\mathfrak m\)-primary ideal \(\mathfrak a\) is the finite discrete set of
normalization branches \(\mathfrak q\) above the closed point, with point
\[
  \operatorname{ord}_{\mathfrak q}/
  \operatorname{ord}_{\mathfrak q}(\mathfrak a).
\]
Changing the primary ideal rescales this same branch set by the exact
cocycle \textup{(Book Chapter~70, \(\dagger \))}.  In particular, branches
cannot be indexed only by irreducible components: the irreducible node in
Book Example~70.1.4 has one minimal prime and two
normalization branches.

There is also an exact link-level finite-map formula.
Book Proposition~70.1.2 applies to an arbitrary finite
local homomorphism of reduced Nagata curve germs, without injectivity.  When
a source branch has ramification index \(e\) over a target branch, both its
order and the order of the extended ideal \(\mathfrak aS\) acquire the
factor \(e\), so normalization cancels it.  The resulting restriction maps
are compatible with composition and change of primary ideal.  They need not
be surjective; this remains a link-level result.

There is now also an exact conductor-controlled ideal-level range.
Book Proposition~70.1.3 shows that every
sufficiently deep integral vector of branch orders is realized by one
downstairs nonzerodivisor generating an \(\mathfrak m\)-primary ideal.
Consequently Book Theorem~70.1.5 identifies the
principal-ideal profiles with the strictly negative rational orthant, their
uniform closure with all nonpositive functions on the finite link, and their
raw obstacle envelope with \(\min\{\phi,0\}\).  The monotone boundary is
sharp: all nonpositive profiles admit increasing approximation from below,
whereas a profile with a zero branch coordinate cannot be a decreasing limit
of primary-ideal profiles.  These conclusions concern only the explicitly
defined ideal-function family.

There is now also a complete order-theoretic calculus for the book-local
ideal-closure candidate.  By
Book Theorem~70.1.6, its continuous ideal-psh
class is every finite nonpositive function on the branch link, and its
decreasing completion is every extended nonpositive function.  Equivalently,
it is the intrinsic cone of nonpositive degree-one homogeneous functions on
the unnormalized branch-valuation rays.  Changing the primary ideal merely
changes the cross-section and obeys an exact cocycle; finite local maps pull
the class back, preserve principal models, and commute with the envelope
\(\min\{\phi,0\}\).  Their pullback image is exactly the fiber-constant
subclass; pullback is injective precisely for a surjective branch map and
surjective precisely for an injective branch map, while surjective branch
maps act isometrically on the continuous class.  On the completed lattices,
fiberwise supremum and infimum give explicit left and right adjoints to
pullback, with \(-\infty\) and zero as the two empty-fiber conventions; the
resulting source-side composites are the least fiber-constant majorant and
greatest fiber-constant minorant, so their fixed points are exactly the
functions descending through the branch map.
The completed class retains
the identically
\(-\infty\) bottom element because a nonsurjective branch map can pull a
nonbottom function to it; this occurs explicitly in
Book Example~70.1.8 for the projection of a node onto
one branch.

The reverse direction is now exact for every finite map whose branch map is
surjective.  For a base branch \(\mathfrak q\), set
\(\delta_{\mathfrak q}=\dim_{K_{\mathfrak q}}
(S\otimes_RK_{\mathfrak q})\).  By
Book Theorem~70.1.7,
\[
 \delta_{\mathfrak q}
   =\sum_{\mathfrak q'\mid\mathfrak q}
      e_{\mathfrak q'/\mathfrak q}f_{\mathfrak q'/\mathfrak q}.
\]
The corresponding unnormalized weighted traces compose in every
branch-surjective finite tower.  Dividing on each base branch by
\(\delta_{\mathfrak q}\) gives a normalized average between the fiberwise
infimum and supremum and a left inverse to pullback.  Conversely, if a base
branch is missed, pullback is not injective and hence cannot have any left
inverse.  Thus branch surjectivity is the exact retraction criterion.  For a
finite real-valued obstacle \(\psi\), the same normalized average satisfies
the sharp Jensen relation
\[
 \operatorname{Av}^{\mathrm{gen}}(P^{\mathrm{id}}\psi)
   \leq P^{\mathrm{id}}(\operatorname{Av}^{\mathrm{gen}}\psi);
\]
equality holds exactly when no branch fiber contains both a strictly negative
and a strictly positive obstacle value.  On finite-valued functions,
\(\rho^*\operatorname{Av}^{\mathrm{gen}}\) is moreover a positive norm-one
idempotent projection onto the fiber-constant pullback image.  Its finite-set
dual lifts every signed base measure by distributing the mass at
\(\mathfrak q\) among the branches \(\mathfrak q'\mid\mathfrak q\) with
weights \(e_{\mathfrak q'/\mathfrak q}f_{\mathfrak q'/\mathfrak q}/
\delta_{\mathfrak q}\); pushing this lift forward recovers the original
measure, the lift is an isometry for total variation, and its pairing with a
source function equals the pairing of the base measure with the normalized
average.  Composing lift and pushforward gives the adjoint positive norm-one
projection onto measures having those canonical conditional weights on each
fiber.  If the base carries any full-support probability measure and the
source carries its canonical lift, the function projection is precisely the
conditional expectation onto the fiber-constant functions: it obeys the
pullback-module identity, contracts every \(L^p\) norm, is the orthogonal
projection in \(L^2\), and satisfies the exact within-fiber variance
decomposition.  Equality in the \(L^p\) contraction for
\(1<p<\infty\) occurs exactly on the descent image.  This is finite-link
Markov duality, not a construction of a singular Monge--Amp\`ere measure.  For a
source nonzerodivisor, its generic norm orders
\(\sum f_{\mathfrak q'/\mathfrak q}w_{\mathfrak q'}(g)\) become, after one
sufficiently divisible common scaling, a deep integral branch vector; the
conductor theorem realizes it by a single downstairs regular principal
ideal.  This gives exact principal-model descent for the normalized average.
Normalized averages, and dually the signed-measure lifts, compose if and only
if the next-stage generic degree is constant over the intermediate branches;
the unnormalized trace needs no such condition.  The
one-branch--two-branch--three-branch tower in
Book Example~70.1.5 shows this is sharp: its direct
normalized average is \(-2\), while its iterated normalized average is
\(-7/4\); both unnormalized traces equal \(-6\).  On the dual side, a unit
base mass lifts directly to \((1/3,1/3,1/3)\), but successive lifting gives
\((1/2,1/4,1/4)\).

Finite flatness makes this generic descent canonical.  If the local map has
constant rank \(d\),
Book Theorem~70.1.8 proves the branch degree
identity
\(\sum_{\mathfrak q'\mid\mathfrak q}e_{\mathfrak q'/\mathfrak q}
f_{\mathfrak q'/\mathfrak q}=d\) and the determinant-norm order formula
\[
 w_{\mathfrak q}(\mathrm N_{S/R}(g))
   =\sum_{\mathfrak q'\mid\mathfrak q}
      f_{\mathfrak q'/\mathfrak q}w_{\mathfrak q'}(g).
\]
Consequently the normalized weighted branch average is an order-preserving
retraction of pullback, lies between the fiberwise infimum and supremum,
commutes with a change of normalizing ideal, and composes in finite flat
towers.  On principal models it is literally
\(u^{\mathfrak aS}_{g,N}\mapsto
u^{\mathfrak a}_{\mathrm N_{S/R}(g),dN}\).  Thus the source conductor
realization descends through one determinant norm.  This is a genuine
finite-map pushforward without the conductor scaling.  Flatness is used for
this canonical global determinant norm, not for the preceding generic trace
on an arbitrary branch-surjective finite map.

The label ``ideal-psh'' is a scope marker, not a comparison theorem.  No
argument here identifies this closure with an independently constructed
singular nef-b-divisor, CLD, or complex-analytic psh cone, and it supplies no
Monge--Amp\`ere product or energy class.  Thus the curve case of the explicit
ideal-closure and envelope calculus is complete, but it does not by itself
discharge the full psh or envelope atoms.

The most direct comparison can nevertheless be computed exactly whenever
the finite normalization is a union of \(q\) closed discs, is an isomorphism
away from one point, and its conductor contains sufficiently high separated
branch powers.  By
Book Theorem~70.1.9, a homogeneous
profile \(c\in[-\infty,0]^q\) has a literal logarithmic radial extension
which is classically psh if and only if it has one which is CLD-v2 psh, and
the exact subcone is
\[
                         \{0\}\cup[-\infty,0)^q.
\]
The positive direction forms the normalized tuple whose retained components
are the monomials \(t_i^{k_i}\), whose \(-\infty\)-components are zero, and
descends it through the conductor to one analytic function; scaled
logarithms and a decreasing limit then give all remaining slopes.  For the
converse, a
zero branch forces a nonnegative nodal value by upper semicontinuity; any
simultaneously negative branch then contradicts its closed-disc Shilov
maximum principle.  Hence the mixed-zero ideal profile \((-1,0)\) is
excluded, while every profile with all coordinates strictly negative or
\(-\infty\) does extend.  The conductor realizes the required branchwise
monomials by single downstairs analytic functions, so this supplies the
exact branch-compatibility rule throughout the stated disc-normalized
curve range.  The ordinary multinode theorem
Book Theorem~70.1.10 records the explicit
coordinate-ring instance.  At the opposite extreme, the unibranch cusp in
Book Example~70.1.9 has conductor
\(t^2K\langle t\rangle\), and its cone
\(\{0\}\cup[-\infty,0)=[-\infty,0]\) contains every nonpositive radial
profile; the mixed-zero obstruction is therefore genuinely multibranch.
The result does not decide whether a different
boundary-trace comparison exists, or treat a singular germ outside this
normalization-and-conductor range.

There is a genuine, but different, global boundary case.  Over an
algebraically closed field of characteristic zero, the global
trivially-valued framework used in \cite[\S2]{BJKS2} treats an integral
projective variety without initially assuming normality: it supplies a
global psh class, its finite-energy space and energy, and a
Monge--Amp\`ere operator with dense finite-energy image.  Exact
surjectivity there is equivalent to the envelope property, which that
source records in the smooth case.  This global construction neither
selects a normalized link for an excellent local germ nor supplies the
curve-local conductor realization above, local bounded products, or the
required local finite-map formula.
It therefore does not discharge any part of
\Cref{prob:frontier-singular-local}.

The remaining programme tests higher-dimensional normal crossings, quotient
singularities, and purely inseparable universal homeomorphisms.  Any claim of
independence of psh or intersection constructions from boundary choices
must respect the established radial comparison and conductor quotient.
Relative b-divisors can be defined on suitable model systems, but beyond
the normalized boundary and, in dimension one, the branch link, its ideal-psh candidate and envelope,
its arbitrary finite-map pullback, its branch-surjective generic trace and
conductor descent, its finite conditional-expectation calculus, its
finite-flat canonical norm pushforward, and the
principal-ideal range above, the independently
geometric psh and intersection-theoretic package is
presently \emph{open}.

\index{local volume!resolved questions}
Two older local-volume questions should not be carried forward as open.
For the Boucksom--de Fernex--Favre local volume, Yuchen Zhang proves a
positive dimension-dependent lower bound for every nonzero volume of an
isolated Gorenstein singularity and constructs a non-
\(\mathbf Q\)-Gorenstein isolated singularity of volume zero that admits no
boundary making the pair log canonical
\cite{Zhang14LocalVolume}.  Thus the first conjectural
direction is confirmed and the second implication is disproved.  These
results delimit the singular theory above; they are resolved history, not
additional problems.

\clearpage\section{Finite-energy equations on polarized affine cones}
\label{family:7}
\noindent\textbf{Active items:} P13.1, P13.2. \textbf{Status: Open in the residual scope.}

\subsection{Atomic targets}
\subsubsection*{P13.1 (T4): Comparison for nonproportional Reeb polarizations}
\label{atom:P13.1}
\addcontentsline{toc}{subsection}{P13.1: Comparison for nonproportional Reeb polarizations}
For a normal polarized affine cone \((X,T;\xi)\) over an algebraically closed trivially valued characteristic-zero field, compare \(\PSH(\xi)\) and \(\PSH(\xi^{\prime})\) for nonproportional Reeb fields. Obtain the estimates strong enough for uniform differentiability of energy. Identities along one proportional Reeb ray do not answer this question.

\subsubsection*{P13.2 (T4): The image of the fixed-Reeb Monge--Ampère operator}
\label{atom:P13.2}
\addcontentsline{toc}{subsection}{P13.2: The image of the fixed-Reeb Monge--Ampère operator}
Characterize precisely the image of \(\MA_\xi:\mathcal E^1(\xi)/\mathbf R\to\mathcal M^1(\xi)\) and the measures admitting a solution. Retain the invariant carrier \(X_0^T\) and the affine-cone probability normalization. The operator is already bi-Lipschitz onto its image, so uniqueness there is established; the residual is the image and existence characterization.

\subsection{What constitutes a solution}
For P13.1 the estimates must compare nonproportional Reeb fields and have the uniformity actually needed for energy differentiation. For P13.2 state necessary and sufficient conditions on measures on the invariant carrier, respecting the existing probability convention; bi-Lipschitzness onto the image does not imply surjectivity.

\subsection{Full statement and detailed state of the art}
\label{sec:ch72-affine-cones}

\begin{problem}\label{prob:frontier-affine-cone-ma}
Let \((X,T;\xi)\) be a normal polarized affine cone over an algebraically
closed trivially valued field of characteristic zero.  Develop a comparison
calculus for \(\PSH(\xi)\) and \(\PSH(\xi')\) when \(\xi\) and \(\xi'\) are
nonproportional Reeb fields, strong enough to prove the required uniform
differentiability of the energy.  For a fixed \(\xi\), characterize the image
of
\[
 \MA_\xi:\mathcal E^1(\xi)/\mathbf R
 \longrightarrow \mathcal M^1(\xi),
\]
and determine exactly which \(\mu\in\mathcal M^1(\xi)\) admit a solution of
\(\MA_\xi(\varphi)=\mu\), unique modulo constants.
\end{problem}

The fixed-polarization theory is already substantial:
Book Theorem~70.9.1 constructs the mixed operator,
proves decreasing-net continuity, and makes the single-potential operator
bi-Lipschitz onto its image.  Here \(\mathcal M^1(\xi)\) uses the invariant
carrier \(X_0^T\) of Book Definition~70.9.1: all affine-cone psh
functions factor through torification, and all canonical model and
finite-energy measures are supported on that closed carrier.  A measure
on the full link outside \(X_0^T\) cannot be an output; energy tests see
only its torification.  Thus the remaining image problem concerns
finite-energy measures on \(X_0^T\), with the affine-cone probability
normalization \textup{(Book 70.62)} explicitly retained.
The proportional-rescaling identities of
\cite[Proposition~3.3]{Wu24AffineCones} do not compare nonproportional Reeb
rays.  Rational approximation likewise does not yield constants uniform in
two varying polarizations.  The available arXiv source is still
version~1 at this cutoff.  Its ``Outlook and relation to other works''
identifies uniform differentiability as the key missing input and locates
it in understanding \(\PSH(\xi)\) versus \(\PSH(\xi')\) when the Reeb
fields are not multiples \cite{Wu24AffineCones}.  It therefore supplies
neither a cross-polarization estimate nor surjectivity onto
\(\mathcal M^1(\xi)\).

The natural strategy is to prove Reeb-uniform Hodge-index and
\(I\)--\(J\) estimates on a compact slice of the Reeb cone, use them to
differentiate the energy along rational approximants, and then run the
finite-energy variational method.  At present the first step is unavailable:
the known addition formula stays on one Reeb ray, and an irrational output
measure may already be carried by a genuinely quasimonomial point.  No
argument in this edition closes either the cross-polarization estimate or the
image characterization.  Both residuals are therefore \emph{open} and are
recorded as one coupled, unresolved Rethlas packet in the editorial ledger;
no verification verdict is claimed for either residual.

\clearpage\section{Pseudoeffective DXZ--MP realization}
\label{family:8}
\noindent\textbf{Active items:} P14.1, P14.4, P14.5, P14.6, P14.8. \textbf{Status: Open in the residual scope.}

\subsection{Atomic targets}
\subsubsection*{P14.1 (T4): Intrinsic raw-transform image beyond the known regimes}
\label{atom:P14.1}
\addcontentsline{toc}{subsection}{P14.1: Intrinsic raw-transform image beyond the known regimes}
For a pseudoeffective class \(\alpha\) on a compact K\"ahler manifold, intrinsically characterize the image of the perturbation-independent raw ideal-valuative DXZ transform on the MP spectrum. The full K\"ahler image and the projective ample rational-line-bundle regime are established. The big non-nef image remains open; the full real N\'eron--Severi DXZ--BJ comparison does not assert the same literal ideal formula outside the rational-line-bundle case.

\subsubsection*{P14.4 (T4): A P-sensitive big non-nef positivity theory}
\label{atom:P14.4}
\addcontentsline{toc}{subsection}{P14.4: A P-sensitive big non-nef positivity theory}
Construct a positivity theory for big non-nef classes retaining \(P\)-singularity data, movable-positive-part data, or equally informative additional structure. A relatively nef total-space model cannot represent a non-nef generic class. The raw ideal-valuative transform is only \(\mathcal I\)-level data and cannot by itself distinguish \(P\)-types or recover raw mass.

\subsubsection*{P14.5 (T5): Monge--Ampère measures in the big non-nef realization}
\label{atom:P14.5}
\addcontentsline{toc}{subsection}{P14.5: Monge--Ampère measures in the big non-nef realization}
Construct the Monge--Ampère measure theory corresponding to a valid big non-nef, \(P\)-sensitive realization, with the correct mass and comparison with existing positive-volume DXZ and MP operators in their common scopes.

\subsubsection*{P14.6 (T5): Energy in the big non-nef realization}
\label{atom:P14.6}
\addcontentsline{toc}{subsection}{P14.6: Energy in the big non-nef realization}
Construct the corresponding intrinsic energy theory and its continuity, differentiation and measure compatibility. State which additional singularity data and reference choices it needs; positive-volume \(\mathcal I\)-level energy alone does not supply a \(P\)-sensitive theory.

\subsubsection*{P14.8 (T5): Zero-mass and formal-data reconstruction}
\label{atom:P14.8}
\addcontentsline{toc}{subsection}{P14.8: Zero-mass and formal-data reconstruction}
Go beyond the established realizable-\(\mathcal I\)-data dichotomy: formulate and solve arbitrary formal-data realization, a general all-semivaluation reconstruction formula, or the remaining zero-mass strict-interior reconstruction range with precise admissibility conditions. A canonical normalized \(\mathcal I\)-model can be recovered from realizable data, but raw representatives and raw mass cannot descend to the \(\mathcal I\)-quotient. The terminal DXZ left-limit result does not create a second zero-mass interval.

\subsection{What constitutes a solution}
Identify the data retained by the proposed realization and prove its order, homogeneity, approximation and operator comparisons in the named positivity class. Distinguish \(\mathcal I\)-data from \(P\)-data. Known upper semicontinuity, the full K\"ahler image and strict-interior filtration maximality must not be relisted as unresolved. P14.8 must respect the verified raw-data obstruction.

\subsection{Full statement and detailed state of the art}
\label{sec:ch72-mp-pseff}

\begin{problem}\label{prob:frontier-mp-pseff}
Let \(X\) be a compact Kähler manifold and let
\(\alpha\in H^{1,1}_{\mathrm{BC}}(X,\mathbf R)\) be pseudoeffective.
Construct the strongest intrinsic realization on the Mesquita--Piccione
tropical spectrum of the DXZ inverse-system space
\(\PSH^{\NA}(X,\alpha)\), and identify its order, homogeneity,
approximation, energy, and Monge--Ampère structures whenever these exist.
\end{problem}

The verified core is now Book Chapter~71.  Every compatible DXZ
family defines a perturbation-independent raw ideal-valuative function on
\(X^\beth\).  Its divisorial restriction is the sourced DXZ
analytification, which is injective and order-reflecting.  Moreover,
Book Theorem~71.5.2, Book Corollary~71.5.1
proves upper semicontinuity of this raw transform at every semivaluation,
including nondivisorial points, by decreasing continuous finite-ideal
approximants.  That candidate has passed its independent mathematical and
static audits.  An initial exact V14 packet, together with seven cited
line-range extracts from five hash-pinned primary-source files, received a
source-closed production pass.  A later combined exact-live audit nevertheless
detected one omitted reduction: the cited DEL00 subadditivity theorem is stated
for psh weights, whereas the levels here are only quasi-psh.  The live proof now
writes each level locally as a psh function plus a smooth term and uses
invariance of multiplier ideals under the resulting locally bounded
perturbation before applying subadditivity.  The repaired exact-current V15
packet, with the perturbation-independent transform and the five one-curve
conclusions supplied as exact premise cards, received final verdict
\texttt{correct}, with zero critical errors and zero gaps, for item
\nolinkurl{pi_4e14dc8a6a26aa190cd191b5}.

The 11 September V14 preproduction release bound the exact then-current
Chapter~71 statement, proof, and example, the current Codex executable, a
single static model catalog, and all fixed source routes.  Its frozen
source-stage tree passed a 58-row manifest, 197 regression tests, all release
gates, and an independent \texttt{gpt-6-astra}/xhigh hostile audit with zero
critical errors and zero material gaps.  A fresh final-path smoke then crossed
the cache horizon with one 310.005920-second MCP call, exactly one structured
post-call status, and the required process chronology; an independent replay
again found zero critical errors and zero gaps.  The exact target, the failed
V13 execution-envelope receipt, the successful V14 smoke, and both audit
summaries are archived in
\nolinkurl{solver_dossiers_20260911/p142_v14_preproduction_gate_20260911}.
The exact production request and response, the source-extract materializer,
and their SHA-256 ledger are archived in
\nolinkurl{solver_dossiers_20260919/}
\nolinkurl{p142_v14_source_closed_production_20260919}.
That archive is retained as the superseded pre-repair pass.  The exact-current
V15 target, premise cards, failed first review, accepted retry, and hash ledger
are archived in
\nolinkurl{solver_dossiers_20260919/}
\nolinkurl{p142_v15_qpsh_repair_production_20260919}.
The repaired V15 production pass and the present live integration retire
P14.2.  The full
K\"ahler comparison below is therefore unconditional within its stated
K\"ahler hypotheses.
Translation holds on all semivaluations, while
the other pointwise operation formulas are retained on the rational
divisorial locus with the exact endpoint hypotheses of
\cite{Xia23Operations}.  The K\"ahler finite-energy correspondence with MP
is proved.  The domination-preserving rational-divisorial approximation in
Book Theorem~71.7.2, combined with the P14.2
upper-semicontinuity theorem and the existing divisorial comparison, now
identifies the literal raw transform with the corresponding MP potential at
every semivaluation in the finite-energy K\"ahler domain; see
Book Corollary~71.7.1. Its source inputs,
extended-value arithmetic, and draft build are reproducibly checked.  The
frozen V3 source-bound packet embeds immutable excerpts of MP24 and DXZ23,
with their full-file hashes.  Following its initial review, the same frozen
statement and proof received a production review by \texttt{gpt-5.6-sol} at
\texttt{xhigh}, with final verdict \texttt{correct}, zero critical errors,
and zero gaps (\texttt{pi\_75b42272e15105bb6be0525c}).  Its exact archived
HTTP receipt and strict protocol audit are recorded in
\nolinkurl{solver_dossiers_20260909/p143_v3_source_bound_packet_20260909/}
\nolinkurl{PRODUCTION_RECEIPT_SOL_20260909.md}.  This verifies exactly the
conditional dominated-approximation and all-point comparison packet:
P14.2 upper semicontinuity and the Chapter~71 finite-energy divisorial
normalization remain explicitly named premises.  This finite-energy result
did not by itself establish general MP positivity or retire P14.3.
The historical local V4 exact-live audit in
\nolinkurl{solver_dossiers_20260909/p143_v3_source_bound_packet_20260909/}
\nolinkurl{LIVE_SOURCE_AUDIT_V4_20260909} records then-current Chapter~71
lines, the accepted V3 packet, and this boundary by local hashes and a range
checker.  It is a local integrity gate, not a replayable release pin or a
production verifier receipt; the production receipt above likewise does not
alter the P14.3 status.
The full MP section theorem now removes the opposite image obstruction.
By the all-point upper-semicontinuity theorem,
Book Theorem~71.7.3 constructs, for every
\(\phi\in\PSH_{\mathrm{MP}}(\alpha)\), a unique DXZ object whose raw
ideal-valuative transform is \(\phi\) on all semivaluations.  Choose any
decreasing rational PL MP model net \(\phi_i\downarrow\phi\).  Its
finite-energy DXZ preimages decrease automatically by divisorial order
reflection; their endpoints decrease to the finite value \(\phi(0)\), so
their DXZ infimum exists.  The divisorial infimum formula identifies its raw
transform with \(\phi\) on rational divisorial points, and the dominated
comparison theorem promotes this equality to all of \(X^{\bethan}\).  The
resulting section preserves and reflects order and is independent of the
model net.

The first \texttt{gpt-6-astra}/xhigh review rejected an earlier wording
which allowed a model net with limit identically \(-\infty\); the constant
net \(\phi_n=-n\) made the endpoint bound fail.  After the nontrivial-limit
condition was made explicit, the repaired packet received final
\texttt{correct} verdicts from both \texttt{gpt-6-astra}/xhigh and the
\texttt{gpt-5.6-sol}/xhigh production review.  A later adversarial manuscript
audit detected that the written proof treated positive rational divisorial
points but had not separately displayed the equality at the trivial
semivaluation.  The final proof now uses the endpoint formula for decreasing
DXZ infima to obtain
\(\boldsymbol\Gamma_{\max}=\phi(0)\) and hence
\(\mathfrak T_{\boldsymbol\Gamma}(0)=\phi(0)\) before applying the positive
divisorial formula.  A fresh read-only
\texttt{gpt-5.6-sol}/xhigh verification of these final files returned
\texttt{correct}, again with zero critical errors and zero gaps, and bound
the verdict to their SHA-256 values.  The rejected and accepted requests,
source checks, and exact hashes are archived in
\nolinkurl{solver_dossiers_20260913/p143_model_monotone_closure_20260913}.
This proves that the raw-transform image contains the entire MP cone,
including possibly infinite-energy potentials.

The reverse inclusion is now supplied by the K\"ahler truncation theorem
Book Theorem~71.7.4.  For an arbitrary DXZ object
\(\boldsymbol\Gamma\), the truncations
\(\boldsymbol\Gamma^{(j)}=\boldsymbol\Gamma\vee(\mathbf0-j)\) have a literal
zero tail below \(-j\), hence full positive volume and finite energy, and they
decrease to \(\boldsymbol\Gamma\).  Their finite-energy MP images decrease to
a nontrivial limit.  The compactness--directedness argument in
Book Proposition~60.4.2 proves that this limit is MP psh;
the P14.2 upper-semicontinuity theorem and the dominated divisorial comparison
identify it with the literal raw transform at every
semivaluation.  Combined with the preceding section, this gives a full order
isomorphism between the DXZ and MP cones in every K\"ahler class, without
projectivity, rationality, or a finite-energy restriction.

The exact statement, proof, and cited MP24, DXZ23, and Xia source ranges are
archived in
\nolinkurl{solver_dossiers_20260919/}
\nolinkurl{p143_full_kahler_truncation_integration_20260919}.
The source-bound V4 production review returned final verdict
\texttt{correct}, zero critical errors, and zero gaps.
Its checked item is \nolinkurl{pi_88c7024842d3fd7d91b55d28}; it treated
P14.2 as an explicit premise.  After the P14.2 repair, all three current target proofs were
rematerialized byte-for-byte.  The unconditional packet supplies seventeen
exact internal premise statements, the accepted P14.2 V15 receipt, and one
consistent extract family for each of MP24, DXZ23, and Xia.  Its repaired
production review by \texttt{gpt-5.6-sol} at \texttt{xhigh} returned final
verdict \texttt{correct}, zero critical errors, and zero gaps for item
\nolinkurl{pi_90c4637f53444d88823e88f3}.  The exact-live packets, the diagnostic
failed review, the accepted response, and the SHA-256 ledger are archived in
\nolinkurl{solver_dossiers_20260919/}
\nolinkurl{p143_unconditional_post_p142_production_20260919}.
Thus P14.3 has no independent residual in the K\"ahler setting and is retired
from the active inventory; the result is unconditional under the stated
K\"ahler hypotheses.

On the finite-energy domain,
Book Corollary~71.7.2 identifies the restriction with the
finite-energy MP correspondence.  The new theorem gives the intrinsic
full-cone K\"ahler image characterization asked for in P14.1; that atom
remains active only beyond the K\"ahler and established projective regimes.
If \(X\) is smooth projective and
the class is pseudoeffective
in the real N\'eron--Severi space, the full DXZ--Boucksom--Jonsson
correspondence is known; for a \(\mathbf Q\)-line-bundle class, the raw
ideal-valuative function is exactly the Boucksom--Jonsson psh metric on
every semivaluation by \cite[formula~(3.8) and Theorem~3.12]{Xia23Operations}.
For a connected smooth projective variety and every ample real
N\'eron--Severi class, Book Theorem~71.7.5 proves the
unrestricted MP--Boucksom--Jonsson identification; for an ample
\(\mathbf Q\)-class its composition with \cite[Theorem~3.14]{DXZ23}
also identifies the full DXZ and MP cones without an energy hypothesis.
This source-explicit book-level deduction retires the former P00.1 atom;
it is not presented as a novelty or priority claim.

The problem remains open in the following precise senses: characterize the
image intrinsically outside the K\"ahler and established smooth-projective
regimes, and construct a \(P\)-sensitive positivity, measure, and energy
theory for big non-nef classes.  One cannot obtain the latter by replacing
\(P\)-equivalence with \(\mathcal I\)-equivalence.

The strict-interior raw-filtration question of Darvas--Xia--Zhang is now
settled in Book Theorem~71.9.1.  Under the hypotheses of
\cite[Definition~4.1]{DXZ23}, it proves
\[
 u_\tau^{\mathcal F}
 =P[u_\tau^{\mathcal F}]
 =P[u_\tau^{\mathcal F}]_{\mathcal I}
 \qquad\text{for every }\tau<\tau_L(\mathcal F),
\]
without finite generation.  The second equality was already
\cite[Lemma~4.5]{DXZ23}; the new proof supplies equality with the raw
function by a fixed-degree birational estimate.  This retires only P14.7.
For a later approximation theorem for graded-series equilibrium metrics,
see \cite[Theorem~4.3]{Finski26Kodaira}.

The terminal left-limit consequence in
Book Corollary~71.9.1 does not reconstruct an arbitrary
zero-mass potential.  The separate atom P14.8 therefore remains active.  A
frozen exact-dichotomy candidate now specifies its input as a realizable
\(\mathcal I\)-class in a compact K\"ahler class, its target as the normalized
\(\mathcal I\)-model representative, and its reconstruction map as the
corresponding \(\mathcal I\)-envelope.  Its independently audited proof gives
the positive reconstruction result for that canonical representative and a
polar Cantor-family obstruction to reconstructing a raw representative or
its raw nonpluripolar mass.  An exact primary-source v4 packet has received a
zero-finding independent verification of this proof.  A later three-source
v5 successor makes the generalized-Cantor capacity interface explicit through
the frozen Kellay--Mashreghi formula and has also received a final
zero-finding review.  The source-bound V6 successor appends the exact
Boucksom sequential/valuative-Lelong, BEGZ locality/polar-set, and
Kellay--Mashreghi capacity excerpts to its frozen proof.  Its independent
model-backed review returned final \texttt{correct}, with zero critical
errors and zero gaps, for item
\nolinkurl{pi_af3d4e7db16b025b3ed90770}.  Its statement and proof SHA-256
values are
\begin{center}
\small
\begin{tabular}{@{}c@{}}
\nolinkurl{4e4256d738ec65dea040e82b3b40bffb233197e9f466c07166760196cc13c04e}\\
\nolinkurl{a51d3e150958265f8b2aa0067d8cf7d64a6c68e14d65e6520fed52bb94489648}.
\end{tabular}
\end{center}
The exact reconstruction and the full Cantor family are now recorded in
Book Corollary~51.8.1, Book Example~51.9.1.  The separate
exact-live V7 audit, pinned to commit
\nolinkurl{ff333ee24648faa9ae23cf8d87a4bea1ec62e961}, checks that these
exact source excerpts give the same dichotomy without enlarging its scope.
After its initial \texttt{gpt-6-astra}/xhigh source-and-boundary review, the
same frozen V7 statement and proof received production
\texttt{gpt-5.6-sol}/xhigh review, final \texttt{correct}, with zero critical
errors and zero gaps, for \nolinkurl{pi_c2863cf20cfa0c4bba834a16}.  The
request, response, strict receipt audit, immutable source-range checker, and
scope boundary are preserved in
\nolinkurl{solver_dossiers_20260909/p148_v6_source_closed_packet_20260909/}
\nolinkurl{LIVE_INTEGRATION_AUDIT_V7_20260909}.  This discharges the
production-review gate for that frozen conditional comparison.  P14.8 still
does not retire or change the inventory count: the packet neither realizes
arbitrary formal data nor proves an all-semivaluation formula or a second
zero-mass interval strictly before \(\tau_L(\mathcal F)\).

The V2 release removed the stale-source ambiguity at its recorded
checkpoint.  It is pinned to commit
\nolinkurl{01ba38fbd824f75ba4150f1913150950f9d4c19d}, whose Chapter~51
source has SHA-256
\nolinkurl{7f2b4890d6332d042113d89bfa921f0f58bbfcbd63c5a18a9eff21469a0fc40b}
and uses the correct arbitrary-divisor error \(A(E)/m\).  The frozen source
stage passed 197 tests and a zero-finding \texttt{gpt-6-astra}/xhigh audit;
the final path then passed its long smoke.  Two earlier attempts were
strictly rejected and were not reused.  The accepted
\texttt{gpt-5.6-sol}/xhigh run
\nolinkurl{20260911T121602Z_4e4256d738ec_e04ebb78} returned final
\texttt{correct}, with zero critical errors and zero gaps; an independent
hostile replay closed all 87 tool-call bindings and again found no critical
error or material gap.  The release, rejected-run separation, and exact
receipts are archived in
\nolinkurl{solver_dossiers_20260911/p148_current_source_release_v2_20260911}.
Thus the realizable-\(\mathcal I\)-data dichotomy is production-verified
against that pinned source: the normalized canonical \(\mathcal I\)-model is
determined without a positive-mass assumption, whereas the raw representative
and its raw mass are not.  This still does not realize arbitrary formal data,
prove a general all-semivaluation formula, or retire the broader P14.8 atom.

A later complex result closes one narrower envelope-only gap: in a big
class, for every complex potential \(u\), the Darvas--Di Nezza--Lu ceiling
equals \(P_\theta[u]\), and the latter is idempotent even when its
Monge--Amp\`ere mass vanishes \cite[Theorem~1.2]{PSWZ26}.  It constructs
neither a non-Archimedean realization or raw transform nor a DXZ
filtration interval, and it does not recover a raw representative or raw
mass from an \(\mathcal I\)-class.  It therefore does not retire P14.8.

Two unchanged extensions are excluded by
Book Proposition~71.1.1, Book Proposition~71.10.2.
A relatively \emph{nef} total-space class restricts to a nef class on every
noncentral fibre, so vertical corrections and central modifications that
remain relatively nef cannot represent a big non-nef generic class.  If
\(\operatorname{vol}(\alpha)=0\), every top nonpluripolar mass vanishes and
the existing probability and normalized energy formulas are undefined.

There is a separate datum obstruction to an \(\mathcal I\)-only repair of
the big non-nef case.  The ideal-valuative transform records
\(\mathcal I\)-equivalence, whereas \(P\)-equivalence is strictly finer
(Book Section~71.2, Book Corollary~51.8.2).  Already
in an ample class on \(\mathbf P^1\), Book Example~51.9.1 gives
\(\mathcal I\)-equivalent positive-mass potentials which are not
\(P\)-equivalent and whose raw nonpluripolar masses are respectively
\(\tfrac12\) and \(1\); consequently raw mass does not descend to the
\(\mathcal I\)-quotient (Book Corollary~51.9.1).  Thus a
future big non-nef theory whose state consists only of the existing
ideal-valuative data cannot itself be \(P\)-sensitive, nor can it recover
raw mass without additional data.  This is a design constraint, not a
nonexistence theorem: a construction retaining a \(P\)-class, a movable
positive part, or comparable information remains open, as does the
positivity problem itself.  In particular, the surviving positive-volume
DXZ energy of Book Proposition~71.10.1 is an
\(\mathcal I\)-model/test-curve statement and does not supply that missing
\(P\)-level datum.

\clearpage\section{Zero-volume energy, measures and convex bodies}
\label{family:9}
\noindent\textbf{Active items:} P15.1, P15.2, P15.3. \textbf{Status: Open in the residual scope.}

\subsection{Atomic targets}
\subsubsection*{P15.1 (T5): Path-independent unnormalized zero-volume energy}
\label{atom:P15.1}
\addcontentsline{toc}{subsection}{P15.1: Path-independent unnormalized zero-volume energy}
For a pseudoeffective class \(\alpha\) with \(\operatorname{vol}(\alpha)=0\), choose the minimum necessary reference scale and construct a path-independent unnormalized energy theory. If K\"ahler perturbations and a leading volume coefficient are used, prove independence of the stated perturbation choices. Division by zero volume is not an admissible normalization.

\subsubsection*{P15.2 (T5): Path-independent zero-volume Monge--Ampère theory}
\label{atom:P15.2}
\addcontentsline{toc}{subsection}{P15.2: Path-independent zero-volume Monge--Ampère theory}
Construct a path-independent zero-volume Monge--Ampère theory with its carrier, dimension, mass and reference choices explicitly stated. A family of perturbative probability measures is not automatically a canonical limiting operator, and raw mass cannot be recovered from \(\mathcal I\)-equivalence alone.

\subsubsection*{P15.3 (T4): P-sensitive lower-dimensional transcendental Okounkov bodies}
\label{atom:P15.3}
\addcontentsline{toc}{subsection}{P15.3: P-sensitive lower-dimensional transcendental Okounkov bodies}
Construct a transcendental Okounkov body retaining a prescribed \(P\)-singularity type and encoding lower-dimensional numerical data at zero volume. A class-level compact body of zero ambient volume is already known; it does not identify numerical dimension or retain the requested \(P\)-type.

\subsection{What constitutes a solution}
Name the reference scale, dimensional carrier and retained singularity data, then prove independence of all claimed perturbation choices. Do not normalize zero mass to a probability. A class-level body with zero ambient volume does not supply a \(P\)-sensitive numerical-dimension theorem or an energy/operator construction.

\subsection{Full statement and detailed state of the art}
\label{sec:ch72-zero-volume}

\begin{problem}\label{prob:frontier-zero-volume}
After choosing the minimum necessary reference scale, is there a
path-independent unnormalized energy and Monge--Ampère theory for a
pseudoeffective zero-volume class?  Is there a \(P\)-sensitive
transcendental Okounkov body encoding its lower-dimensional numerical data?
\end{problem}

The unchanged probability signature is \emph{disproved}: total mass zero
cannot be normalized to mass one.  Kähler perturbations give a family of
probabilities, not automatically a canonical limit.  Moreover
Book Theorem~69.6.1 proves that a body depending only on
\(\mathcal I\)-equivalence cannot recover actual nonpluripolar mass.

The attempted positive construction rescales by the first nonzero coefficient
in \(\operatorname{vol}(\alpha+\varepsilon\omega)\), then tests independence
of the perturbation path and comparison with numerical dimension.  Nef
zero-volume classes and toric examples provide the first nontrivial cases.
Any surviving convex body must use \(P\)-type or richer data.  No general
path-independent construction is presently verified.

A distinct published class-level endpoint must be kept separate.
Darvas--Reboulet--Witt Nystr\"om--Xia--Zhang attach a natural compact
convex set to a pseudoeffective transcendental cohomology class and a
chosen flag, and prove that its full \(n\)-dimensional volume equals
\(\operatorname{vol}(\alpha)/n!\) \cite[pp.~135--136]{DRWNXZ}.  Hence a
zero-volume class has such a canonical class-level body of zero
\(n\)-volume.  This does not identify its lower-dimensional volume or
dimension with the numerical dimension, attach a body to a prescribed
\(P\)-singularity type, or construct an unnormalized energy or
Monge--Amp\`ere operator.  It therefore narrows the starting point for
P15.3 but does not close any of the three zero-volume atoms.

Xia's zero-volume benchmark uses the limit partial body obtained by
intersecting the bodies in ample perturbations and conjectures that its
dimension equals the numerical dimension
\cite[Section~5.6 and Conjecture~5.29]{Xia21}.
The same paper's Chebyshev question has a positive-volume hypothesis
\cite[Section~6.4, p.~1325]{Xia21}: fix a big line bundle on a smooth
projective variety, its smooth curvature form \(\theta\), an admissible
flag and a continuous metric weight \(v\).  Assume
\(\int_X\theta_{P_\theta[\varphi]_{\mathcal I}}^n>0\), replace
\(\varphi\) by that ideal-model projection, and choose
\(\eta\in\PSH(X,\theta)\) such that \(\theta_\eta\) is a K\"ahler
current and \(\eta\preceq\varphi\).  Is the transform
\(\lim_{\varepsilon\downarrow0}
c_{[(1-\varepsilon)\varphi+\varepsilon\eta]}[v]\)
on \(\operatorname{Int}\Delta_\nu(\theta,\varphi)\) independent of
\(\eta\)?  This is an adjacent positive-volume question, rather than a
zero-volume atom P15.1--P15.3.  Both source questions remain published
benchmarks; the stronger \(P\)-sensitive zero-volume body requested above
is formulated in this book and should not be attributed to that source.

The third conjecture of that source is the non-toric mixed-volume formula
\cite[Conjecture~8.8]{Xia21}.  On 24 August 2026, the author supplied the TeX
proof underlying the unpublished manuscript \cite{XiaMVOB}.  After repairs
of its positive-Cao-volume reduction, convex sandwich, pointwise potential
sandwich, rational analytic-resolution step, degenerate finite-level case,
multigraded saturation step, and effective-shift body inclusion, the exact
repaired proof received a production \texttt{gpt-5.6-sol}/\texttt{xhigh}
verdict of \emph{correct}, with no critical error and no gap.  Thus
Book Theorem~55.5.1 proves Conjecture~8.8 for singular psh metrics on
big line bundles over a smooth projective complex variety when every Cao
(equivalently, ideal-model) volume is positive;
Book Theorem~55.5.2 gives the movable-product specialization for
big line bundles with minimal singularities.  This is an author-supplied,
repaired, and independently verified proof in an unpublished work in
progress, not a claim that the result has appeared in the published
literature.  It asserts a supremum over surjective rank-\(n\) valuations,
not attainment by one valuation.  Zero volume, transcendental classes,
singular varieties, and attainment remain outside the theorem; in
particular, the genuinely different problem
\Cref{prob:frontier-zero-volume} remains open.
This source-status item is not an additional problem environment in the
literal count at the beginning of the chapter.

\clearpage\section{The general non-Archimedean Monge--Ampère equation}
\label{family:10}
\noindent\textbf{Active items:} P16.1, P16.2, P16.3, P16.4, P16.5, P16.6, P16.7. \textbf{Status: Open in the residual scope.}

\subsection{Atomic targets}
\subsubsection*{P16.1 (T5): Finite-energy equations on singular or unibranch varieties}
\label{atom:P16.1}
\addcontentsline{toc}{subsection}{P16.1: Finite-energy equations on singular or unibranch varieties}
Extend the finite-energy non-Archimedean Monge--Ampère operator and its existence-and-uniqueness theorem to singular or unibranch varieties, replacing each use of smoothness in compactness, envelopes, energy differentiation and comparison. State all retained polarization and field hypotheses.

\subsubsection*{P16.2 (T5): Proper nonprojective algebraizable equation theories}
\label{atom:P16.2}
\addcontentsline{toc}{subsection}{P16.2: Proper nonprojective algebraizable equation theories}
Construct an appropriate replacement positivity datum and the associated psh, finite-energy, comparison and equation theory on proper nonprojective algebraizable spaces. A proper algebraic scheme with an ample line bundle is already projective, so retaining that datum while dropping projectivity is not a new case.

\subsubsection*{P16.3 (T5): Finite-energy equations in positive residue characteristic}
\label{atom:P16.3}
\addcontentsline{toc}{subsection}{P16.3: Finite-energy equations in positive residue characteristic}
Construct the required finite-energy operator and prove existence and uniqueness when the original residue-characteristic restriction is removed. Replace the resolution, model-regularity and envelope inputs wherever they are used rather than merely changing the standing field hypothesis.

\subsubsection*{P16.4 (T4): Finite-energy equations over arbitrary complete valued fields}
\label{atom:P16.4}
\addcontentsline{toc}{subsection}{P16.4: Finite-energy equations over arbitrary complete valued fields}
Extend the operator and existence theorem to nondiscrete valuations and arbitrary complete non-Archimedean fields outside the established invariant-toric, skeleton-supported endpoint. A definition of the finite-energy space alone does not establish the operator or the equation.

\subsubsection*{P16.5 (T5): Finite-energy equations for nonample positivity classes}
\label{atom:P16.5}
\addcontentsline{toc}{subsection}{P16.5: Finite-energy equations for nonample positivity classes}
Extend the equation theory to big or pseudoeffective nonample classes, specifying the psh cone, reference singularity type, energy and total mass. At zero volume, the existing probability normalization must be replaced rather than retained unchanged.

\subsubsection*{P16.6 (T4): Arbitrary-support finite-energy measures}
\label{atom:P16.6}
\addcontentsline{toc}{subsection}{P16.6: Arbitrary-support finite-energy measures}
Prove the operator and existence-and-uniqueness theorem for arbitrary finite-energy measure support beyond the known toric skeletal regime. Measures must satisfy the correct mass and energy conditions. Do not reduce arbitrary support to a single skeleton or replace the genuine Monge--Ampère measure by another inverse measure construction.

\subsubsection*{P16.7 (T5): Finite-energy equations on nonalgebraizable analytic spaces}
\label{atom:P16.7}
\addcontentsline{toc}{subsection}{P16.7: Finite-energy equations on nonalgebraizable analytic spaces}
Construct the finite-energy operator and its existence-and-uniqueness theory on genuinely nonalgebraizable analytic spaces, with an explicit positivity datum and the necessary compactness, envelope and comparison replacements.

\subsection{What constitutes a solution}
For each axis, state the retained hypotheses and construct the finite-energy operator together with the corresponding existence-and-uniqueness theorem. Specify the measure mass and finite-energy condition, uniqueness modulo constants and the replacement inputs for compactness, envelopes, energy differentiation and comparison. An endpoint solving a special toric measure class does not remove a general axis.

\subsection{Full statement and detailed state of the art}
\label{sec:ch72-general-ma}

\begin{problem}\label{prob:frontier-general-ma}
Extend the finite-energy existence and uniqueness theory for the NA
Monge--Ampère equation beyond the currently proved combinations of smoothness,
projectivity, residue characteristic, valuation, positivity class, support,
and algebraizability hypotheses.
\end{problem}

This is not one yes-or-no assertion.  Positive characteristic, singular or
unibranch varieties, nondiscrete valuations, big classes, arbitrary support,
and nonalgebraizable analytic spaces require different inputs.  The known
variational proofs use compactness, envelope regularity, differentiability of
energy, and a comparison principle; deleting a geometric hypothesis is valid
only after replacing every place where it enters.  The relevant positive
results and their distinct domains are surveyed in
Book Chapter~28, Book Chapter~29.

The projectivity axis has one logical constraint.  A proper algebraic
\(K\)-scheme endowed with an ample line bundle is already projective over
\(K\) \cite[Tag~0B5N]{StacksProject}.  Thus P16.2 cannot mean that one
retains the BFJ line-bundle polarization \(L\) while simply deleting the word
``projective.''  A nonprojective proper algebraizable extension must first
replace that positivity datum---and then construct the corresponding psh,
finite-energy, and comparison theory---rather than merely rerun the original
argument.  No canonical replacement with the required equation theory is
asserted here.

One exact source-stated axis is especially clean.  Pille--Schneider observes
that the finite-energy space can be defined over an arbitrary
non-Archimedean field, but that the extension of the Monge--Amp\`ere operator
is not known there in general; the discretely valued characteristic-zero and
trivially valued cases are the established endpoints
\cite[Remark~4.4]{PS23Hybrid}.  This should be solved before bundling all
remaining axes into a single theorem.

\begin{proposition}[Arbitrary-field invariant toric skeletal endpoint]
\label{prop:p164-arbitrary-field-toric-skeletal-endpoint}
Let \(K\) be a complete nontrivially valued non-Archimedean field,
\(X_\Sigma\) a proper toric \(K\)-variety, and \(L\) an ample toric line
bundle.  If \(\mu\) is a positive Radon measure on \(X_\Sigma^{\an}\), is
supported on the canonical skeleton \(N_\Sigma\), gives zero mass to
\(N_\Sigma\setminus N_{\mathbf R}\), and has total mass
\(\deg_L(X_\Sigma)\), then there is a toric continuous singular psh metric
on \(L\) whose Monge--Amp\`ere measure is \(\mu\).  It is unique up to
multiplication of the metric by a positive constant.
\end{proposition}



Thus nondiscreteness and the arbitrary complete ground field cause no
obstruction in this invariant toric, skeleton-supported range.
Here ``continuous singular'' has the source's extended-valued meaning:
local weights are continuous into \(\mathbf R\cup\{-\infty\}\), as in
\cite[\S5.3.1]{BGJKK25}.  It does not imply finite real-valued weights or
finite energy.  This is not an arbitrary-field finite-energy theorem in the
BFJ sense: the result permits
a singular potential, requires torus invariance and skeletal support, and
does not construct the general extended Monge--Amp\`ere operator.  In
particular, P16.4 and P16.6 retain the respectively nonskeletal and
noninvariant finite-energy residual scopes.

There is also a distinct, nonproper algebraic endpoint.  Let \(n>d\geq1\),
let \(\overline X\) be a smooth projective Fano complex variety of dimension
\(n\), and let \(D\) be a reduced simple-normal-crossing anticanonical
divisor whose dual complex is the standard \((d-1)\)-simplex.  For the
affine Calabi--Yau complement \(X=\overline X\setminus D\), Wang's
Theorem~A constructs a continuous non-Archimedean psh function on
\(X^{\an}\), over the trivially valued field \(\mathbf C\), whose local
Monge--Amp\`ere measure is the canonical Lebesgue measure on the essential
skeleton \cite{Wang25NACalabiYau}.  The source is a version-one preprint and
uses a local CLD/tropical--convex construction.  In particular, it neither
constructs the BFJ finite-energy operator nor treats arbitrary measures or
arbitrary affine varieties.  Since \(X\) is affine and hence not proper,
this endpoint does not answer P16.2; it only records a special nonprojective
algebraic solution after the properness requirement is also dropped.

There is a separate finite-atomic inverse statement which must not be
called a finite-energy Monge--Amp\`ere theorem.  Under its standing
assumptions that \(K\) is maximally complete and
\(\lvert K^\times\rvert=\mathbf R_{>0}\), Fang fixes an irreducible
projective \(K\)-variety with an ample line bundle and a finite set of graded
birational points on its affine cone.  For every probability vector on that
finite set, the source constructs a Shilov-finite, possibly singular metric
having the prescribed \emph{equidistribution} measure
\cite[Proposition~4.8]{Fang25Shilov}.  The source itself says that comparing
this Banach-algebra measure with the genuine Monge--Amp\`ere measure requires
the separate continuity property of an envelope metric.  Thus it supplies
neither a continuous semipositive solution nor the BFJ finite-energy operator
for the prescribed atomic measure, and does not narrow P16.6.

\index{Monge--Amp\`ere equation!skeletal factorization}
There is also a resolved negative regularity chronology.  The first
skeletal-factorization question of Gubler--Jell--K\"unnemann--Martin already
fails on their toric surface \cite[Appendix~A]{GJKM19}.
Boucksom--Jonsson then construct a finite
atomic Monge--Amp\`ere measure supported on a dual complex whose solution is
invariant under no retraction attached to any SNC model; they also construct
a divisorial Green function that is not \(\mathbf Q\)-PL
\cite[Theorems~A and~B]{BJGreen24}.  Hence support on a dual complex does not
imply factorization through some refined skeleton.  This is a theorem-level
obstruction to a proof strategy, not an open regularity expectation.

The Rethlas programme treats each axis separately and first searches the
post-2020 literature for an existing theorem.  Alterations and universal
homeomorphisms can address some algebraic descent questions, but they do not
by themselves supply multiplier ideals or differentiability in positive
characteristic.  The unrestricted formulation is \emph{open}; several
strictly delimited axes are already known.

\clearpage\section{Residual hybrid entropy and YTD questions}
\label{family:11}
\noindent\textbf{Active items:} P19.1, P19.2, P19.3, P19.4. \textbf{Status: Open in the residual scope.}

\subsection{Atomic targets}
\subsubsection*{P19.1 (T5): Entropy slopes of arbitrary finite-energy weak rays}
\label{atom:P19.1}
\addcontentsline{toc}{subsection}{P19.1: Entropy slopes of arbitrary finite-energy weak rays}
Prove the remaining entropy-slope statement for arbitrary finite-energy weak rays on a compact K\"ahler manifold, outside the stated smooth polarized projective maximal-ray results. State the non-Archimedean datum, the entropy normalization and the precise ray hypotheses.

\subsubsection*{P19.2 (T4): Approximation preserving energy, entropy and singularities}
\label{atom:P19.2}
\addcontentsline{toc}{subsection}{P19.2: Approximation preserving energy, entropy and singularities}
Construct model approximations in the residual hybrid programme that preserve all required energy, entropy and singularity data. Convergence of energy alone is insufficient, and the published/source-stated projective Fujita regimes are established subcases.

\subsubsection*{P19.3 (T5): Hybrid compactness and measures for degenerating families}
\label{atom:P19.3}
\addcontentsline{toc}{subsection}{P19.3: Hybrid compactness and measures for degenerating families}
Prove the compactness and measure-continuity statements needed along genuinely degenerating families of compact K\"ahler manifolds. A fixed quasi-projective variety hybridized with explicit scale factors is a different known endpoint.

\subsubsection*{P19.4 (T5): Equivariant, transcendental and singular hybrid extensions}
\label{atom:P19.4}
\addcontentsline{toc}{subsection}{P19.4: Equivariant, transcendental and singular hybrid extensions}
Establish the required automorphism-equivariant transcendental and singular extensions, with the exact stability and metric signatures specified. Retain the residual distinctions from the cited smooth polarized projective and static K\"ahler-class correspondences.

\subsection{What constitutes a solution}
Freeze the exact ray, approximation, family or equivariant signature and prove the residual statement with its entropy and energy normalizations. Keep recent preprint claims in their stated smooth polarized scope. The four targets are independent residual hybrid goals; they are not prerequisites for the already cited static uniform-model-stability-to-cscK implication.

\subsection{Full statement and detailed state of the art}
\label{sec:ch72-hybrid-ytd}

\begin{problem}\label{prob:frontier-hybrid-ytd}
Develop the remaining entropy-slope, approximation, compactness, and
automorphism statements for arbitrary finite-energy weak rays and
degenerating families in the hybrid K-stability/YTD programme on a compact
Kähler manifold.
\end{problem}

The static projective entropy and stability gap must not be listed as wholly
open at this cutoff.  Boucksom--Jonsson's version-two notes of 19 May 2026
give a structural complex-to-non-Archimedean overview and only sketch a
version of the cscK YTD argument \cite{BJ26ComplexNA}.  They are used here
as orientation, not as a separate approximation or hybrid-continuity theorem
that narrows any atom in \Cref{prob:frontier-hybrid-ytd}.

Trusiani's version-two preprint of 17 August 2026,
\cite[Theorem~B and Corollary~A]{Trusiani26}, retains the special Fujita
approximations for big line bundles and the claimed solution of the
Boucksom--Jonsson entropy-regularization conjecture; its Theorem~A retains
the automorphism-relative uniform YTD theorem for smooth polarized
projective varieties.  Corollary~A also records the resulting equality of
the Archimedean entropy slope along a maximal finite-energy ray with the
non-Archimedean entropy, in that projective framework.  The new Appendix by
S\'ebastien Boucksom supplies a non-Archimedean proof over
\(\Bbbk((t))\), and Theorem~C adds the weighted-relative correspondence,
including Calabi extremal metrics when the weight is log-concave.  These are
recent preprint claims, rather than results independently reproved in this
book.  Boucksom--Jonsson's 18 September 2025
preprint \cite{BJ25WeightedYTD} gives a
non-Archimedean formulation for cscK and weighted extremal metrics on
polarized projective manifolds.  Darvas--Zhang's 3 November 2025 version-two
preprint \cite[Theorems~1.6 and~5.5]{DarvasZhang25YTD} gives a uniform
YTD correspondence for a unique cscK metric on an arbitrary compact
Kähler manifold using \(K^\beta\)-stability and, in the polarized
projective case, a reductivity-plus-log-discrepancy-model formulation.
Mesquita--Piccione--Witt Nyström prove that uniform model stability implies
existence in a Kähler class and establish the transcendental finite-energy
Calabi--Yau homeomorphism \cite[Theorems~A and~B]{MPWN25}.  Finally,
Trusiani's Corollary~B identifies the standard automorphism-relative
projective stability formulations in the smooth polarized setting
\cite{Trusiani26}.  Thus the projective entropy slope, the weighted-relative
correspondence in the stated smooth polarized scope, and comparison of the
listed projective signatures are no longer open problems.

There is a complementary hybrid endpoint outside the compact proper setting.
For each of the three affine/Calabi-model situations stated in Wang's
Theorem~B, explicit scale factors make the Archimedean potential and its
non-Archimedean counterpart into one continuous function on the hybrid space;
the correspondingly rescaled Archimedean Monge--Amp\`ere measures converge
weakly to the non-Archimedean measure \cite[Theorem~B]{Wang25NACalabiYau}.
This is a comparison for a fixed quasi-projective variety over the hybrid
parameter, not a compactness or measure-continuity theorem for a degenerating
family of compact K\"ahler manifolds.  It therefore does not discharge
P19.3.

Consequently, the atoms in \Cref{prob:frontier-hybrid-ytd} are not asserted
to be logical prerequisites for the cited uniform-model-stability-to-cscK
implication.  They record the distinct residual hybrid programme: statements
for arbitrary finite-energy weak rays, log-discrepancy-model regularization,
and genuinely degenerating families that are not supplied merely by the
static Kähler-class theorem.

Book Chapter~66, Book Chapter~68 separates hybrid convergence from
entropy and coercivity.  The BHJ erratum is enforced in
Book Chapter~65.  The remaining targets are narrower: direct
algebraic proofs of the known projective equivalences; \(K^\beta\)-stability
in the presence of nontrivial automorphisms; equivariant transcendental and
singular extensions; regularization of rays attached to log-discrepancy
models; and the genuinely hybrid compactness and measure-continuity
statements needed to pass through a degenerating family.  Boucksom--Jonsson
also isolate bounded support size for divisorial test measures, algebraic
equivariance proofs, and extension to polarized klt pairs
\cite[Questions and further remarks]{BJ25WeightedYTD}.  Those residual
problems remain \emph{open}; the smooth polarized projective entropy and
weighted-relative theorems are dated Trusiani version-two claims, not
missing inputs.

\clearpage\section{Continuous-pair Bedford--Taylor independence}
\label{family:12}
\noindent\textbf{Active items:} P21.3. \textbf{Status: Open in the residual scope.}

\subsection{Atomic targets}
\subsubsection*{P21.3 (T2): Bedford--Taylor independence for continuous equal-curvature pairs}
\label{atom:P21.3}
\addcontentsline{toc}{subsection}{P21.3: Bedford--Taylor independence for continuous equal-curvature pairs}
For a complete non-Archimedean field and a good separated strict boundaryless pure \(n\)-dimensional analytic space, \(n\ge2\), decide whether every positive \(\dprime\)- and \(\ddprime\)-closed CLD \((p,p)\)-current, \(0\le p\le n-1\), has Bedford--Taylor independence on every ordinary open for continuous psh pairs with equal Hessian currents. The full quantified equations appear below. A negative answer must give an analytic counterexample in this scope, not just a defective abstract tropical pushforward.

\subsection{What constitutes a solution}
Prove the quantified every-open, every-current continuous-pair assertion, or supply an analytic counterexample with all ambient hypotheses. A fixed polycylinder, a cover of one specified open and a genuine ordinary-open basis give different conclusions. Weak approximability is a sufficient condition, not a density theorem for arbitrary currents.

\subsection{Full statement and detailed state of the art}
\label{sec:ch72-bedford-taylor}

The two original realization questions have been settled and retired.  The
full unbounded cross cannot literally be the image of a harmonic
tropicalization whose source is compact, by
Book Proposition~69.2.3.  Conversely,
Book Proposition~69.2.4 constructs a reduced
boundaryless analytic curve with a proper noncompact tropicalization onto
the cross, exhausted by compact harmonic-tropicalization charts, and an
effective cycle whose current has weights \((2,2,1,1)\).  That construction
realizes the psh approximants and the defect after proper pushforward, but
the two curvature currents are unequal before pushforward.  It therefore
does not answer the analytic Bedford--Taylor question upstairs.

The normal-curve theorem
Book Proposition~16.4.8 settles the
ordinary current problem on every connected good separated normal
boundaryless pure analytic curve over a complete non-Archimedean field,
possibly trivially valued: every positive closed current of bidegree
\((0,0)\) is \(c\delta_Y\).  It requires neither strictness nor a smoothness
hypothesis and gives the full CLD Bedford--Taylor predicate.

The componentwise theorem
Book Proposition~16.4.9 settles the displayed
continuous-pair problem for every good separated strict boundaryless pure
one-dimensional analytic space over either valuation regime, without
reducedness or a smoothness hypothesis on component normalizations.  If
\((X_i)_i\) are the irreducible components of \(X_{\mathrm{red}}\), every
positive closed current of bidegree \((0,0)\) has the unique locally finite
expression
\[
                         T=\sum_i c_i\delta_{X_i},
                         \qquad c_i\geq0.
\]
Nilreduction preserves the ordinary CLD calculus and current topology while
the fundamental integration current retains generic component
multiplicities.  The two local estimates
Book Lemma~16.4.3, Book Lemma~16.4.4 prevent cancellation at conductor
and component-intersection points in the nontrivial and trivial valuation
regimes, respectively.

A separate comparison now applies in every dimension \(n\geq1\) and over
every complete non-Archimedean field:
Book Corollary~16.14.2 proves on every ordinary open
that equal canonical curvatures of an arbitrary continuous
\(\operatorname{CPsh}\) pair are equal against every reduced analytic
component current.  Generic nonreduced multiplicities are retained, and
divisorial component intersections are allowed.  Its input is the
ordinary-local full-\(\operatorname{CPsh}\) analytic-nullity theorem, not a
classification of arbitrary positive closed currents.  Its linear companion
Book Corollary~16.14.3 therefore covers every
nonnegative component-diagonal current
\(\sum_i c_i\delta_{W_i}\) on the same ordinary open.  Its smooth-factor
extension Book Corollary~16.14.4 also covers every
fixed smooth positive closed decoration
\(\bigl(\sum_i c_i\delta_{W_i}\bigr)\wedge\gamma\) in bidegree
\((q,q)\), \(0\leq q\leq n-1\).  This remains a component-diagonal
current with a prescribed smooth factor: it neither treats a general
positive closed current nor supplies a variable or singular
positive-bidegree factor.  More generally,
Book Corollary~16.14.5 covers the weak limits, with
\(q\) fixed, of finite nonnegative sums of such smooth-component currents.
This is a stated weak-approximability condition, not a density theorem for
positive closed currents.  Its local counterpart
Book Corollary~16.14.6 permits the component
families and the approximating nets to vary over an ordinary open cover.  It
is still a stated local weak-approximability condition; it does not claim
that an arbitrary positive closed current admits such a cover.



There is now also a higher-dimensional local positive result.  Suppose that
a good topologically separated pure boundaryless analytic space has a
genuine basis of ordinary open sets on which every continuous function with
zero CLD Hessian is ordinary CLD-smooth with zero smooth Hessian.  Then
Book Proposition~16.4.14 shows that every positive
\(\mathrm d'\)- and \(\mathrm d''\)-closed CLD current has
continuous-pair Bedford--Taylor independence.  By
Book Corollary~16.4.2, the hypothesis holds when
the basis members are standard strict open polycylinders.  On one such
polycylinder, Book Proposition~16.4.15 gives the
same identity for pairs defined on the whole polycylinder, without claiming
the every-open predicate.

There is a further nonproduct basis range.  If every ordinary point
neighbourhood belongs to a basis of boundaryless strict finite towers of
locally trivial standard open-polydisc bundles over a standard strict
endpoint polycylinder, then
Book Corollary~M.1.59 gives continuous-pair
Bedford--Taylor independence for every positive closed CLD current.  The
endpoint-tower classification makes the equal-curvature difference an
ordinary CLD-smooth zero-Hessian function, so the proof uses the local smooth
pluriharmonic criterion rather than any pullback of the current.  It requires
an actual ordinary-open basis of the displayed boundaryless towers and does
not cover a single chart in an otherwise arbitrary germ, annular fibres, or
fibrations without the stated local product charts.

The endpoint base now also supports proper and mixed fibres.  In
Book Theorem~M.1.25, each stage may be a
locally trivial standard strict open-polydisc bundle, a projective-space
bundle, a smooth proper split-toric bundle, or---over an algebraically
closed field---a proper smooth connected curve bundle, in any finite order.
Every continuous zero-current function
on the entire top-level space is the pullback of a unique logarithmic-affine
function on the endpoint polycylinder, hence is ordinary CLD-smooth with
zero smooth Hessian.  The whole-space, cover, and basis consequences
Book Corollary~M.1.60, Book Corollary~M.1.61, Book Corollary~M.1.62
therefore give Bedford--Taylor independence for arbitrary positive closed
ordinary CLD currents in the corresponding scopes.  The cover proof first
globalizes the smooth zero-Hessian difference and does not invoke current
sheaf gluing outside its stated hypotheses.  Projective, split-toric, and
proper-curve stages retain the exact local holomorphic-separation and
paracompactness requirements of their descent theorems.

This narrows the higher-dimensional residual; it does not turn a
polycylinder cover into a basis and does not settle arbitrary analytic
germs.

There is now a second, genuinely fibred basis range controlled by the
residual geometry of a curve rather than by an explicit polycylinder
formula.  For a smooth proper curve over an algebraically closed complete
nontrivially valued field,
Book Proposition~19.3.1 proves that a continuous
zero-current function is ordinary CLD-smooth with zero smooth Hessian
exactly when its type--\(2\) residual slope classes vanish in
\(\operatorname{Pic}^0(C_x)\otimes_{\mathbf Z}\mathbf R\).  Thus vanishing
of these groups at every type--\(2\) point removes the ordinary-smooth
obstruction completely.  This includes the cases in which the residue field
is algebraic over a finite field and in which the base is a Mumford curve.

The exact descent through locally trivial standard strict open-polydisc,
projective-space, smooth proper split-toric, and proper smooth connected
curve bundles may then be iterated in any finite order.  The resulting
theorem
Book Theorem~M.1.21 makes every continuous
zero-current function on the whole top-level tower an ordinary CLD-smooth
pullback with zero smooth Hessian.  Hence
Book Corollary~M.1.34 gives whole-space
continuous-pair Bedford--Taylor independence against every positive closed
CLD current; the cover and basis versions
Book Corollary~M.1.35, Book Corollary~M.1.36
give the corresponding local and every-open conclusions.  Projective,
split-toric, and proper-curve stages retain the locally holomorphically
separated and paracompact base hypotheses of their descent theorems, and no nonsmooth
current is pulled back.  A single mixed-tower chart inside an unrelated
open set is still insufficient, and no assertion is made beyond the
displayed residual-Picard bases.

The residual-Picard range now also admits locally twisted annular geometry
at the bottom.  In
Book Theorem~M.1.22, the bottom
space is an analytically locally trivial bundle over the same curve domain
whose fibres are finite products of standard strict open annuli and,
optionally, standard strict open discs; any finite mixed tower of the four
exact bundle types may then be placed above it.  Every continuous
zero-current function on the top is ordinary CLD-smooth with zero smooth
Hessian.  Indeed, on each bottom chart its annular part is a sum of a
residual-Picard-smooth curve function and locally constant multiples of the
smooth affine functions \(\log|T_i|\), while the disc factors and all later
stages are handled by exact descent and pullback.  Consequently,
Book Corollary~M.1.37, Book Corollary~M.1.38, Book Corollary~M.1.39
give the whole-space, specified-cover, and every-open basis versions of
Bedford--Taylor independence against arbitrary positive closed ordinary CLD
currents.  Neither integral annular monodromy nor the zero-current radiance
class is required to vanish, because ordinary smoothness is checked
chartwise.  Accordingly, the theorem produces no global logarithmic
coefficient formula and still requires the displayed local product charts
and the residual-Picard hypothesis.

The annular stages may now also occur at arbitrary levels among the four
exact bundle types.  In
Book Theorem~M.1.23, one analytic
open cover of the residual-Picard curve base makes every base-changed tower
a finite interleaving of globally coordinated standard strict
open-polyannulus products with standard strict open-polydisc,
projective-space, smooth proper split-toric, and proper smooth connected
curve bundles.  The annular radii, coordinates, and product trivializations
may vary with the member of the cover and need not agree on overlaps.  The
interleaved coefficient formula is therefore used only locally: its curve
coefficient is ordinary CLD-smooth by the residual-Picard criterion, while
every logarithmic annular term and all of its later pullbacks are ordinary
CLD-smooth with zero smooth Hessian.  Hence every continuous zero-current
function on the top is ordinary CLD-smooth with zero smooth Hessian.  The
whole-space, specified-cover, and every-open basis consequences are
Book Corollary~M.1.40, Book Corollary~M.1.41, Book Corollary~M.1.42.
This removes the restriction that all annular geometry occur at the bottom,
but gives no global logarithmic coefficient formula and retains both the
curve-local tower charts and the residual-Picard hypothesis.

The common-cover requirement can be replaced by a genuinely stage-local
condition.  Book Proposition~M.1.35
shows that
an analytically locally trivial finite strict open-polyannulus bundle
preserves ordinary CLD-smooth zero-Hessian regularity and first-current
rigidity whenever each member of its chosen trivializing cover already has
the corresponding zero-current regularity and rigidity.  Applying this at
annular levels and the exact descent theorem at the other levels gives
Book Theorem~M.1.24: annular charts may
now live over arbitrary intermediate bases and need not descend from one
cover of the residual-Picard curve.  Its whole-space, specified-cover, and
every-open basis Bedford--Taylor consequences are
Book Corollary~M.1.43, Book Corollary~M.1.44, Book Corollary~M.1.45.
No compatibility of annular coordinates or global logarithmic coefficient
formula is used.  The stage-local hypotheses cannot be recovered from the
corresponding whole-space conclusions in general.  Indeed,
Book Proposition~M.1.36 constructs a
smooth proper surface \(X=(\mathbf P^1_K)^{\mathrm{an}}\times_KD^{\mathrm{an}}\)
on which every global continuous zero-current function is constant, while
an ordinary product open \(\Omega\subset X\) carries a continuous
zero-current function which is not ordinary CLD-smooth.  The concrete
good-reduction elliptic instance is
Book Corollary~M.1.46.  Moreover,
\(X\to(\mathbf P^1_K)^{\mathrm{an}}\) is itself the proper smooth curve-bundle
stage allowed by the mixed exact-tower theorem.  Thus the failure occurs on
an exact intermediate base, not merely on an unrelated space, and the
stage-local chart assumptions are genuinely necessary for this route.
This strictly sharpens the boundary of the verified residual-geometric
class but neither proves nor disproves P21.3.

There is also a verified source-scoped split-toric subcase, including toric
boundary strata.  On an open set of the form
\(V=\operatorname{trop}^{-1}(O)\) for extended tropicalization in a split
toric variety, the
proposition
\(Book Proposition~16.13.1\) gives
continuous-pair Bedford--Taylor independence for every positive closed CLD
current when the difference of the finite continuous psh potentials is
invariant under the maximal compact torus.  The proof uses the
invariant-function correspondence and function-current pushforward
calculation of \cite[(5.5) and Proposition~5.3.10(iv)]{BGJKK25}; finite
continuity along the compactified normal rays forces the zero-Hessian
difference to be smooth at the boundary.  The exact stored two-individually
invariant-potential statement and proof passed the controlling local
independent verifier run with final verdict \texttt{correct}, zero critical
errors, and zero gaps (item \texttt{pi\_96813b21562d920cf411c5a0}).  The
present difference-invariant strengthening is a direct book-level reuse of
that proof: its tropical factorization, pushforward, and boundary-ray steps
use only the difference.  More precisely, its multiplier identity requires
only a finite continuous pair with equal current Hessian and invariant
difference; the psh hypothesis is retained in the displayed consequence
solely because P21.3 asks about psh pairs.  This neither treats pairs with noninvariant
difference, nonsplit tori, arbitrary toroidal charts, nor general analytic
germs.  Scalar extension alone does not enlarge this P21.3 subcase:
the available base-change results identify only specified form/measure degrees
in the CLD calculus, or the top-degree Monge--Amp\`ere measures of continuous
semipositive line-bundle metrics on projective varieties
\cite[Propositions~15.3.1--15.3.2]{CLD25};
\cite[Proposition~8.4.2]{BGJKK25}.
They do not construct, or give faithful descent for, an arbitrary positive
closed CLD current and its mixed current
\(\mathrm d'\mathrm d''(uT)\).  Thus the scalar-stable potential statements
in P06.1/P06.2 do not by themselves yield an algebraic-split-form P21.3
corollary for an arbitrary positive closed current.  There is nevertheless
one paired scalar-stable class.  If the current is locally
weak-form-approximable and the difference satisfies the analytic-local
algebraic-split equivariant tower hypotheses after the displayed scalar extension, then
Book Corollary~M.1.72 gives the
Bedford--Taylor equality.  The scalar-stable condition is used only to make
the downstairs difference weakly smooth and weakly harmonic; its constant
approximation net and the weak-form closure bridge then control the ordinary
order-zero product.  This neither supplies current ascent or descent nor
applies to an arbitrary positive closed current or an arbitrary continuous
zero-current difference.
There is, however, a genuine nonproduct extension at the original split
level.  On every finite tower of equivariant locally trivial standard strict
open-polydisc bundles over such a split-toric base, the exact bundle descent
makes an invariant equal-curvature difference the pullback of an invariant
zero-current base function.  The boundary-ray argument makes that base
function ordinary CLD-smooth with zero Hessian, and CLD-form pullback
preserves this property through the tower.  Therefore
Book Corollary~M.1.65 gives
continuous-pair Bedford--Taylor independence for every positive closed CLD
current and every finite continuous psh pair whose difference is invariant
on a boundaryless top-level tower of dimension at least two.  It requires the
actual split base, an action on every tower level, and equivariance of every
bundle map; it neither treats pairs with noninvariant difference nor descends
arbitrary currents through an algebraic-splitting extension.
This conclusion is local in the spatial variable.  For one specified open
set, an ordinary cover by such boundaryless split-toric equivariant towers,
with the difference invariant for the action on each member, already gives
the equality by
Book Corollary~M.1.66; no basis of the
ambient space is needed.  A basis of actual towers gives the corresponding
conclusion simultaneously for every open set by
Book Corollary~M.1.67.  Neither
criterion turns unrelated local charts into a global toric action or a global
invariance condition.  The same cover-versus-basis distinction holds for the
endpoint-tower criteria
Book Corollary~M.1.58, Book Corollary~M.1.59.

The additional supplied proof
Book Theorem~16.16.1 establishes zero-Hessian
rigidity and continuous-pair multiplication on full products of open disks
and annuli over an algebraically closed nontrivially valued field. This
is an established positive range of P21.3; arbitrary ordinary open domains
and the remaining field and singularity scopes below stay open.

\begin{problem}
\label{prob:frontier-cross-fan-upstairs}
Let \(K\) be a complete non-Archimedean field, let \(n\geq2\), and let \(X\)
be a good separated strict boundaryless pure \(n\)-dimensional
\(K\)-analytic space.  Determine whether every positive
\(\mathrm d'\)- and \(\mathrm d''\)-closed CLD current \(T\) of bidegree
\((p,p)\), \(0\leq p\leq n-1\), has the following property: for every open
\(W\subseteq X\) and all continuous psh functions \(U,V\) on \(W\),
\begin{equation}
                  \mathrm d'\mathrm d''[U]
                  =\mathrm d'\mathrm d''[V]
 \label{eq:p21p3-post-nonreduced-residual-equal-curvature}
\end{equation}
implies
\begin{equation}
 \mathrm d'\mathrm d''\bigl(U\,T|_W\bigr)
 =\mathrm d'\mathrm d''\bigl(V\,T|_W\bigr).
 \label{eq:p21p3-post-nonreduced-residual-bt}
\end{equation}
If this fails, construct a counterexample.  A counterexample could provide an
analytic weighted-gluing defect before tropical pushforward.
\end{problem}

No one-dimensional case remains inside the displayed ambient hypotheses.
Indeed, Book Proposition~16.4.9 treats either
valuation, arbitrary reducibility and nilpotent structure, and normalizations
with no smoothness assumption.  Non-good, non-separated, non-strict, and
boundary cases lie outside the displayed problem rather than forming hidden
residual atoms.  Within dimension at least two, spaces with the strict-open
polycylinder basis of
Book Corollary~16.4.2 are also removed from the
residual scope, as are fixed-domain pairs on one standard strict open
polycylinder.  The invariant split-toric family, including its toric boundary
strata, of
Book Proposition~16.13.1 is likewise
removed.  The atom remains open for general spaces in the displayed ambient
class.


The problem uses the book's continuous-potential convention.  Literal
\cite[Definition~8.6.2]{CLD25} quantifies over the full CLD psh class, which
may include discontinuous functions.  The normal-curve theorem establishes
that stronger statement, while the reducible and nonreduced componentwise
theorem establishes the continuous-pair predicate used here.  The remaining
discontinuous-potential question in the latter setting is neither asserted to
be equivalent to the displayed problem nor separately counted here.  Thus
P21.3 remains one genuinely open residual T2 atom, now only in analytic
dimension at least two and outside the verified basis, fixed-domain, and
invariant split-toric regimes above.  The book-wide inventory remains
12 open problem environments and 42 active atomic questions, with
tier totals \(0/3/7/15/17\) for T1 through T5.

The four-punctured smooth-curve obstruction remains useful as a warning about
the method.  Its coarse tropical weights \((2,2,1,1)\) violate refined
balancing after a third analytic unit is appended, so a coarse tropical
pushforward cannot replace analytic classification.  The well-defined
residual atom is precisely the continuous-pair CLD question in the ranges
not covered by
Book Proposition~16.4.9, Book Corollary~16.4.2, Book Proposition~16.4.15, Book Proposition~16.13.1.  Equality only after a
noninjective tropical pushforward does not decide it.

There is an exact sufficient route for the remaining analytic question.  Put
\(h=U-V\).  If, locally on \(W\), the continuous zero-Hessian function
\(h\) is a locally uniform limit of smooth functions \(h_j\) satisfying
\(\mathrm d'\mathrm d''h_j=0\), then
Book Corollary~16.13.1 gives
\(\mathrm d'\mathrm d''(hT)=0\), hence the required equality for every
positive closed \(T\).  Thus the unresolved continuous-pair problem is
concentrated at higher-dimensional gluing loci where the displayed smooth
pluriharmonic approximation is not known.  Merely knowing
\(\mathrm d'\mathrm d''[h]=0\) does not supply the approximants
(Book Corollary~16.13.1); the abstract cross-fan shows why
cellwise affine data alone cannot replace them.  Nor do the bare real-PL or
GJR-harmonic conclusions occurring in P06.1/P06.2: that G-local,
weakly smooth regularity does not provide ordinary CLD-smooth zero-Hessian
approximants.  Indeed,
Book Proposition~19.3.2 gives an integral PL
weakly smooth harmonic function on a curve with no locally uniform ordinary
CLD-smooth harmonic approximation.  This obstructs this approximation route,
not the separately proved one-dimensional P21.3 classification.  This bridge
is related to, but is not identified with, P06.2.
At the level of strong currents there is nevertheless a formal compatibility:
Book Proposition~18.9.2 gives
\(\mathrm d'\mathrm d''(h\wedge S)=0\) for a weakly smooth zero-Hessian
\(h\) and a closed strong current \(S\).  This is deliberately not used to
settle P21.3.  The latter uses the order-zero ordinary product \(hT\); even
when an ordinary current \(T\) has a strong lift, its equality with the
restriction of \(h\wedge S\) requires a separate compatibility argument, as
spelled out in Book Section~18.9.
There is now one genuine ordinary-current bridge inside this distinction.
If \(T\) is locally weak-form-induced in the sense of
Book Definition~18.9.2, and the difference \(h=U-V\) is a
weakly smooth zero-Hessian function, then
Book Proposition~18.9.3 proves the required equality.  Its
proof identifies the order-zero product \(hT\) with the restriction of the
strong product by uniform ordinary-smooth approximation on compact supports
and the \(C^0\) continuity of weakly smooth integration.  It covers neither an
arbitrary strong lift nor an arbitrary positive closed current, and it still
assumes weak smoothness of \(h\); it is therefore a current-and-pair subcase,
not a resolution of the remaining gluing problem.
The bridge now persists through two natural but explicitly conditional
closures.  If the ordinary current is locally a weak limit of positive closed
locally weak-form-induced currents, and \(h\) is locally a uniform limit of
weakly smooth zero-Hessian functions, then
Book Proposition~18.9.4 gives the same ordinary-current
equality.  The proof first passes the form-induced identity along the current
net and then along the function net, using the order-zero
multiplication-continuity lemma at each step.  It is not a density theorem
for arbitrary positive closed CLD currents or arbitrary continuous
zero-Hessian differences, so it narrows a paired closure class without
removing the unresolved higher-dimensional gluing step.
The exact P06.2
proper-full-rank moment classification supplies this bridge for a genuine
pairwise subcase: if \(h\) has such a factorization locally, then
Book Corollary~M.1.21 gives the required equality
for every positive closed current.  P06.2 supplies no such local
factorization for an arbitrary \(h\), however, so this new exclusion is a
condition on a pair of potentials rather than a resolution of the
higher-dimensional gluing problem.

\clearpage
\section*{Settled questions excluded from the active count}
\addcontentsline{toc}{section}{Settled questions excluded from the active count}
The active list does not revive questions already retired in the repository.
Important excluded signatures include the ample-class full-cone comparison;
real-PL positive-current/residue-nefness comparison; ordinary CLD-smooth
curve regularization; singular classical-pluriharmonic real-PL regularity
(P06.1); the uniform hierarchy bound (P07.2) and the supplied unrestricted
first-stage counterexample (P07.1); global ordinary smooth psh metric
approximation (P08.1, negative); the unrestricted proper analytic
parlinear/global-dpsh extension (P08.2, negative); delta-metricity and signed
PL decomposition/positive-cone obstructions; generalized-Lelong separation;
the global BJ Radon--Riesz implication (negative); excellent-germ normalized
boundaries (P12.1); raw-transform upper semicontinuity (P14.2), full
K\"ahler DXZ--MP comparison (P14.3), and strict-interior raw-filtration
maximality (P14.7); Banach-family raw Monge--Ampère continuity (P17.1);
both listed Favre endpoint-continuity signatures (negative); real-toric
mixed products; and the original cross-fan realization questions.
Local recursive projection and algebraic-chart comparisons P05.3/P05.4
are also retired, with the corrected operator-domain assumptions retained.

A negative unrestricted statement can coexist with the open task of
classifying its valid range. Thus the unconditional nodal-envelope failures
do not retire P03.1--P03.3; nor do zero-volume probability and
\(\mathcal I\)-only mass obstructions retire a theory with new appropriate
data. The recent smooth polarized projective entropy and YTD claims are
kept in their exact preprint scopes and are not relisted as residual atoms.

\section*{Source and verification note}
\addcontentsline{toc}{section}{Source and verification note}
The controlling sources are Chapter 72 and its active inputs, the atomic
inventory, and the dated editorial source/verification ledgers in the
repository \url{https://github.com/MingchenXia/nonarchimedean-pluripotential-theory}.
The exact repository URL is specified in the export manifest; this catalogue
is a local source export, not a claim of remote publication.
The extraction retains all twelve formal problem statements and the full
state-of-knowledge discussions for those families, while omitting proof
bodies. External book references are resolved from the successful current
book build; all bibliography entries used here are embedded verbatim.
The 42 atomic IDs, placement and tiers are checked against the exact active
inventory. This catalogue supplies no new independent mathematical verdict
and no new literature search beyond the dated source coverage.

\clearpage
\phantomsection
\addcontentsline{toc}{section}{References}
\begin{thebibliography}{EGA IV$_4$}
\bibitem[BGGJK25]{BGJKK25}
Jos{\'e}~Ignacio Burgos~Gil, Walter Gubler, Philipp Jell, and Klaus
  K{\"u}nnemann.
\newblock Pluripotential theory for tropical toric varieties and
  non-{A}rchimedean {M}onge--{A}mp{\`e}re equations.
\newblock {\em Kyoto J. Math.}, 65(1):55--152, 2025.

\bibitem[Bin76]{Bingener76Stein}
J\"{u}rgen Bingener.
\newblock Schemata {\"u}ber {S}teinschen {A}lgebren.
\newblock {\em Schr. Math. Inst. Univ. M\"{u}nster (2)}, page~52, 1976.

\bibitem[BJ23]{BJKS2}
S{\'e}bastien Boucksom and Mattias Jonsson.
\newblock A non-{A}rchimedean approach to {K}-stability, {II}: Divisorial
  stability and openness.
\newblock {\em J. Reine Angew. Math.}, 805:1--53, 2023.

\bibitem[BJ24b]{BJGreen24}
S{\'e}bastien Boucksom and Mattias Jonsson.
\newblock Non-{A}rchimedean {G}reen's functions and {Z}ariski decompositions.
\newblock {\em Comptes Rendus Math\'ematique}, 362(S1):5--42, 2024.

\bibitem[BJ25c]{BJ25WeightedYTD}
S{\'e}bastien Boucksom and Mattias Jonsson.
\newblock On the {Y}au--{T}ian--{D}onaldson conjecture for weighted csc{K}
  metrics, 2025.
\newblock Version 1, 18 September 2025.

\bibitem[BJ26]{BJ26ComplexNA}
S\'{e}bastien Boucksom and Mattias Jonsson.
\newblock From complex to non-{A}rchimedean geometry: an approach to the {YTD}
  conjecture, 2026.
\newblock Version 2, 19 May 2026.

\bibitem[BMPS26]{BMPS26}
Ana~Mar{\'i}a Botero, Enrica Mazzon, and L{\'e}onard Pille-Schneider.
\newblock Monge--{Amp\`e}re measures on balanced polyhedral spaces, 2026.
\newblock Manuscript dated 10 March 2026.

\bibitem[CLD25]{CLD25}
Antoine Chambert-Loir and Antoine Ducros.
\newblock Formes diff{\'e}rentielles r{\'e}elles et courants sur les espaces de
  {B}erkovich, 2025.
\newblock Version 2, 25 July 2025.

\bibitem[DRWN{\etalchar{+}}26]{DRWNXZ}
Tam\'as Darvas, R\'emi Reboulet, David Witt~Nystr\"om, Mingchen Xia, and Kewei
  Zhang.
\newblock Transcendental {O}kounkov bodies.
\newblock {\em J. Differential Geom.}, 132(1):135--178, 2026.

\bibitem[DXZ25]{DXZ23}
Tam\'{a}s Darvas, Mingchen Xia, and Kewei Zhang.
\newblock A transcendental approach to non-{A}rchimedean metrics of
  pseudoeffective classes.
\newblock {\em Comment. Math. Helv.}, 100(2):269--322, 2025.

\bibitem[DZ25]{DarvasZhang25YTD}
Tam\'as Darvas and Kewei Zhang.
\newblock A {YTD} correspondence for constant scalar curvature metrics, 2025.
\newblock Version 2, 3 November 2025.

\bibitem[Fan25]{Fang25Shilov}
Yanbo Fang.
\newblock Metrized ample line bundles in non-{A}rchimedean geometry {II}, 2025.
\newblock Version 1, 26 November 2025.

\bibitem[Fin26]{Finski26Kodaira}
Siarhei Finski.
\newblock {K}odaira--{I}itaka dimension and multiplicity: an analytic
  perspective, 2026.
\newblock Version 1, 23 March 2026.

\bibitem[GJKM19]{GJKM19}
Walter Gubler, Philipp Jell, Klaus K{\"u}nnemann, and Florent Martin.
\newblock Continuity of plurisubharmonic envelopes in non-{A}rchimedean
  geometry and test ideals.
\newblock {\em Ann. Inst. Fourier (Grenoble)}, 69(5):2331--2376, 2019.
\newblock With an appendix by Jos{\'e} Ignacio Burgos Gil and Mart{\'i}n
  Sombra.

\bibitem[GJR26]{GJR26}
Walter Gubler, Philipp Jell, and Joseph Rabinoff.
\newblock Forms on {B}erkovich spaces based on harmonic tropicalizations.
\newblock {\em Bull. Soc. Math. France}, 154(1):135--238, 2026.

\bibitem[GR25a]{GR25Classical}
Walter Gubler and Joseph Rabinoff.
\newblock Classically psh and pluriharmonic functions on {B}erkovich spaces,
  2025.
\newblock Version 2, 17 September 2025.

\bibitem[MPWN25]{MPWN25}
Pietro Mesquita-Piccione and David Witt~Nyström.
\newblock A transcendental non-{A}rchimedean {C}alabi--{Y}au theorem with
  applications to the csc{K} problem, 2025.
\newblock Version 1, 11 September 2025.

\bibitem[PS23]{PS23Hybrid}
L{\'e}onard Pille-Schneider.
\newblock Global pluripotential theory on hybrid spaces.
\newblock {\em J. \'{E}c. polytech. Math.}, 10:601--658, 2023.

\bibitem[PSWZ26]{PSWZ26}
Kai Pang, Haoyuan Sun, Zhiwei Wang, and Xiangyu Zhou.
\newblock Capacity stability of complex {M}onge--{A}mp\`ere equations with
  moving prescribed singularities, 2026.
\newblock Version 1, 14 July 2026.

\bibitem[Stacksb]{StacksProject}
The {Stacks Project Authors}.
\newblock Stacks project.
\newblock \url{https://stacks.math.columbia.edu}, 2026.
\newblock Accessed 29 August 2026.

\bibitem[Sup26b]{SuppliedP05220260925}
Withdrawn draft: Orthogonality on one-dimensional balanced fans, 2026.
\newblock Author-supplied withdrawn manuscript dated and received 25 September
  2026; replacement withdrawal version preserved in the book repository.

\bibitem[Tru26]{Trusiani26}
Antonio Trusiani.
\newblock A solution to the {Y}au--{T}ian--{D}onaldson conjecture through
  special {F}ujita approximations, 2026.
\newblock Version 2, 17 August 2026; includes an appendix by S\'{e}bastien
  Boucksom.

\bibitem[Vil21]{Vil20}
Christian Vilsmeier.
\newblock A comparison of the real and non-archimedean {M}onge-{A}mp\`ere
  operator.
\newblock {\em Math. Z.}, 297(1-2):633--668, 2021.

\bibitem[Wan19]{Wanner19}
Veronika Wanner.
\newblock Comparison of two notions of subharmonicity on non-archimedean
  curves.
\newblock {\em Math. Z.}, 293:443--474, 2019.

\bibitem[Wan25]{Wang25NACalabiYau}
Ying Wang.
\newblock Non-{A}rchimedean {C}alabi--{Y}au potentials on certain affine
  varieties, 2025.
\newblock Version 1, 24 October 2025.

\bibitem[Wu24]{Wu24AffineCones}
Yueqiao Wu.
\newblock Some non-{A}rchimedean pluripotential theory on polarized affine
  cones, 2024.
\newblock Version 1, 19 June 2024.

\bibitem[Xia25a]{Xia23Operations}
Mingchen Xia.
\newblock Operations on transcendental non-{A}rchimedean metrics.
\newblock In {\em Convex and complex: perspectives on positivity in geometry},
  volume 810 of {\em Contemp. Math.}, pages 271--293. Amer. Math. Soc.,
  Providence, RI, 2025.

\bibitem[Xia25b]{Xia21}
Mingchen Xia.
\newblock Partial {O}kounkov bodies and {D}uistermaat-{H}eckman measures of
  non-{A}rchimedean metrics.
\newblock {\em Geom. Topol.}, 29(3):1283--1344, 2025.

\bibitem[Xia26a]{XiaMVOB}
Mingchen Xia.
\newblock Mixed volumes of {O}kounkov bodies.
\newblock Work in progress, July 2026; available at
  \url{https://mingchenxia.github.io/Papers/MVOB.pdf}. The author-supplied
  source was repaired and independently audited on 24 August 2026. The full
  repaired proof is reproduced in this book; its production Rethlas proof
  digest is \texttt{9227826d...58cad7b}, with final verdict correct and no
  reported gap, 2026.

\bibitem[Zha14]{Zhang14LocalVolume}
Yuchen Zhang.
\newblock On the volume of isolated singularities.
\newblock {\em Compositio Mathematica}, 150(8):1413--1424, 2014.
\end{thebibliography}
\end{document}
